% !TEX program = lualatex
% ---------------------------------------------------------------------------------------------
%  Readback note 22: Two consecutive theta lifts: admissible combinations, arrow pairs and zero terms
%
%  Build:  python3 tools/rb.py note <slug>
%
%  One unit of a Lean artifact, rewritten as ordinary mathematics. The prose is hand-written. Lean names,
%  line ranges and counts come from data.tex; quoted Lean statements are copied from the artifact by the build.
%
%  Rules of the text: standard terminology; every symbol is explained where it first appears and listed again
%  in the notation at the end; no letter has two meanings; a proof is complete or says what it takes for granted;
%  Lean names appear only in the LEAN lines and in the concordance.
%
%  The signature rule (the build enforces it): a function, relation or predicate that the artifact defines is
%  written only through its notation macro, which the build makes from the spec (`notation`) with one slot for
%  each explicit argument of the Lean definition. No argument is fixed and dropped, also not "for brevity" and
%  not for one section. Every argument and every hypothesis of a stated declaration stands inside its
%  statement, or in a Setting that the statement cites by number; only implicit and instance arguments (the
%  ambient objects) may stand in the conventions.
% ---------------------------------------------------------------------------------------------
\documentclass[11pt,a4paper]{article}
\ifdefined\pdfvariable\pdfvariable suppressoptionalinfo \numexpr 1+2+4+8\relax\fi
\usepackage[a4paper,hmargin=35mm,top=22mm,bottom=27mm,footskip=12mm,nohead]{geometry}

\usepackage{fontspec}
\directlua{
  luaotfload.add_fallback("gfallback", {"STIXTwoMath-Regular.otf:mode=node;", "DejaVuSans.ttf:mode=node;"})
  luaotfload.add_fallback("gmonofallback", {"DejaVuSansMono.ttf:mode=node;", "STIXTwoMath-Regular.otf:mode=node;"})
}
\setmainfont{NewCM10}[Extension=.otf,UprightFont=*-Book,ItalicFont=*-BookItalic,
  BoldFont=*-Bold,BoldItalicFont=*-BoldItalic,SmallCapsFont=*-Book,SmallCapsFeatures={Letters=SmallCaps},
  RawFeature={fallback=gfallback}]
\setsansfont{NewCMSans10}[Extension=.otf,UprightFont=*-Book,ItalicFont=*-BookOblique,
  BoldFont=*-Bold,BoldItalicFont=*-BoldOblique,RawFeature={fallback=gfallback}]
\setmonofont{JetBrainsMono}[Extension=.otf,UprightFont=*-Regular,ItalicFont=*-Italic,
  BoldFont=*-SemiBold,BoldItalicFont=*-SemiBoldItalic,Scale=MatchLowercase,
  RawFeature={-calt;-liga;fallback=gmonofallback}]
\usepackage{amsmath}
\usepackage{amsthm}
\usepackage[math-style=TeX]{unicode-math}
\setmathfont{NewCMMath-Book.otf}
\emergencystretch=2.5em\relax

\usepackage{xcolor}
\usepackage{tikz}
\usetikzlibrary{calc,arrows.meta,positioning,backgrounds}
\usepackage[protrusion=true,expansion=true]{microtype}
\usepackage{enumitem}
\usepackage{needspace}
\usepackage{booktabs}
\usepackage{array}
\usepackage{longtable}
\usepackage{fancyhdr}
\usepackage{lastpage}
\usepackage{multicol}
\usepackage[hidelinks,pdfusetitle]{hyperref}

\definecolor{ink}{HTML}{0F172A}
\definecolor{inksoft}{HTML}{475569}
\definecolor{inkfaint}{HTML}{94A3B8}
\definecolor{rulegrey}{HTML}{CBD5E1}
\definecolor{arith}{HTML}{C2410C}
\definecolor{est}{HTML}{7C3AED}
\definecolor{ana}{HTML}{0891B2}
\definecolor{ok}{HTML}{047857}
\definecolor{leancol}{HTML}{3730A3}

\title{Two consecutive theta lifts: admissible combinations, arrow pairs and zero terms}
\author{grh-compact}
\date{}

% ------------------------------------------------------------------------------------- macros
\newcommand{\leanfont}{\ttfamily\fontsize{8.6}{10.5}\selectfont\color{leancol}}
\newcommand{\lean}[1]{\csname LN@#1\endcsname}
\newcommand{\fillin}[1]{{\color{arith}\itshape[#1]}}  % a place of the template that is not written yet; the build reports every one that is left
\newcommand{\leanquote}[1]{\csname LQ@#1\endcsname}
% generated by rb.py note; do not edit
\newcommand{\ConcDecls}{98}
\newcommand{\Pin}{72fead10}
\newcommand{\PinText}{that of commit {\leanfont 72fead10}}
\newcommand{\ArtifactName}{grh-compact}
\expandafter\newcommand\csname LN@IsAdmissibleComb\endcsname{\href{https://github.com/lsernau/grh-lean/blob/72fead102f663d5c7a407964a13dc0cf9d977207/grh-compact/src/New/Seesaw/AdmissibleComb.lean\#L12-L13}{{\leanfont IsAdmissibleComb}}}
\expandafter\newcommand\csname LN@ofsum\endcsname{\href{https://github.com/lsernau/grh-lean/blob/72fead102f663d5c7a407964a13dc0cf9d977207/grh-compact/src/New/Seesaw/AdmissibleComb.lean\#L20-L22}{{\leanfont IsAdmissibleComb.\allowbreak{}of\_sum}}}
\expandafter\newcommand\csname LN@single\endcsname{\href{https://github.com/lsernau/grh-lean/blob/72fead102f663d5c7a407964a13dc0cf9d977207/grh-compact/src/New/Seesaw/AdmissibleComb.lean\#L24-L25}{{\leanfont IsAdmissibleComb.\allowbreak{}single}}}
\expandafter\newcommand\csname LN@bind\endcsname{\href{https://github.com/lsernau/grh-lean/blob/72fead102f663d5c7a407964a13dc0cf9d977207/grh-compact/src/New/Seesaw/AdmissibleComb.lean\#L32-L40}{{\leanfont IsAdmissibleComb.\allowbreak{}bind}}}
\expandafter\newcommand\csname LN@map\endcsname{\href{https://github.com/lsernau/grh-lean/blob/72fead102f663d5c7a407964a13dc0cf9d977207/grh-compact/src/New/Seesaw/AdmissibleComb.lean\#L42-L42}{{\leanfont IsAdmissibleComb.\allowbreak{}map}}}
\expandafter\newcommand\csname LN@smul\endcsname{\href{https://github.com/lsernau/grh-lean/blob/72fead102f663d5c7a407964a13dc0cf9d977207/grh-compact/src/New/Seesaw/AdmissibleComb.lean\#L44-L46}{{\leanfont IsAdmissibleComb.\allowbreak{}smul}}}
\expandafter\newcommand\csname LN@smulofeqzeroorisOfFinOrder\endcsname{\href{https://github.com/lsernau/grh-lean/blob/72fead102f663d5c7a407964a13dc0cf9d977207/grh-compact/src/New/Seesaw/AdmissibleComb.lean\#L48-L50}{{\leanfont IsAdmissibleComb.\allowbreak{}smul\_of\_eq\_zero\_or\_isOfFinOrder}}}
\expandafter\newcommand\csname LN@ofeqisOfFinOrdersmul\endcsname{\href{https://github.com/lsernau/grh-lean/blob/72fead102f663d5c7a407964a13dc0cf9d977207/grh-compact/src/New/Seesaw/AdmissibleComb.lean\#L52-L54}{{\leanfont IsAdmissibleComb.\allowbreak{}of\_eq\_isOfFinOrder\_smul}}}
\expandafter\newcommand\csname LN@mapbilinear\endcsname{\href{https://github.com/lsernau/grh-lean/blob/72fead102f663d5c7a407964a13dc0cf9d977207/grh-compact/src/New/Seesaw/AdmissibleComb.lean\#L56-L59}{{\leanfont IsAdmissibleComb.\allowbreak{}map\_bilinear}}}
\expandafter\newcommand\csname LN@smulmem\endcsname{\href{https://github.com/lsernau/grh-lean/blob/72fead102f663d5c7a407964a13dc0cf9d977207/grh-compact/src/New/Seesaw/AdmissibleComb.lean\#L64-L67}{{\leanfont IsAdmissibleComb.\allowbreak{}smul\_mem}}}
\expandafter\newcommand\csname LN@pure\endcsname{\href{https://github.com/lsernau/grh-lean/blob/72fead102f663d5c7a407964a13dc0cf9d977207/grh-compact/src/New/Seesaw/AdmissibleComb.lean\#L71-L77}{{\leanfont IsAdmissibleComb.\allowbreak{}pure}}}
\expandafter\newcommand\csname LN@normUnitsmul\endcsname{\href{https://github.com/lsernau/grh-lean/blob/72fead102f663d5c7a407964a13dc0cf9d977207/grh-compact/src/New/Seesaw/AdmissibleComb.lean\#L79-L83}{{\leanfont IsAdmissibleComb.\allowbreak{}normUnit\_smul}}}
\expandafter\newcommand\csname LN@selfOrConj\endcsname{\href{https://github.com/lsernau/grh-lean/blob/72fead102f663d5c7a407964a13dc0cf9d977207/grh-compact/src/New/Seesaw/AdmissibleEvaluation.lean\#L20-L20}{{\leanfont selfOrConj}}}
\expandafter\newcommand\csname LN@mulVecmemweightedTranslations\endcsname{\href{https://github.com/lsernau/grh-lean/blob/72fead102f663d5c7a407964a13dc0cf9d977207/grh-compact/src/New/Seesaw/AdmissibleEvaluation.lean\#L22-L26}{{\leanfont mulVec\_mem\_weightedTranslations}}}
\expandafter\newcommand\csname LN@starmemweightedTranslations\endcsname{\href{https://github.com/lsernau/grh-lean/blob/72fead102f663d5c7a407964a13dc0cf9d977207/grh-compact/src/New/Seesaw/AdmissibleEvaluation.lean\#L28-L30}{{\leanfont star\_mem\_weightedTranslations}}}
\expandafter\newcommand\csname LN@invxFinmulVecmemofspansLattice\endcsname{\href{https://github.com/lsernau/grh-lean/blob/72fead102f663d5c7a407964a13dc0cf9d977207/grh-compact/src/New/Seesaw/AdmissibleEvaluation.lean\#L32-L42}{{\leanfont inv\_xFin\_mulVec\_mem\_of\_spansLattice}}}
\expandafter\newcommand\csname LN@IsNeighbourVector\endcsname{\href{https://github.com/lsernau/grh-lean/blob/72fead102f663d5c7a407964a13dc0cf9d977207/grh-compact/src/New/Seesaw/AdmissibleEvaluation.lean\#L60-L62}{{\leanfont IsNeighbourVector}}}
\expandafter\newcommand\csname LN@IsAdmissibleEval\endcsname{\href{https://github.com/lsernau/grh-lean/blob/72fead102f663d5c7a407964a13dc0cf9d977207/grh-compact/src/New/Seesaw/AdmissibleEvaluation.lean\#L67-L72}{{\leanfont IsAdmissibleEval}}}
\expandafter\newcommand\csname LN@existsisAdmissibleEvalstd\endcsname{\href{https://github.com/lsernau/grh-lean/blob/72fead102f663d5c7a407964a13dc0cf9d977207/grh-compact/src/New/Seesaw/AdmissibleEvaluation.lean\#L81-L160}{{\leanfont exists\_isAdmissibleEval\_std}}}
\expandafter\newcommand\csname LN@isAdmissibleEvalfinEval\endcsname{\href{https://github.com/lsernau/grh-lean/blob/72fead102f663d5c7a407964a13dc0cf9d977207/grh-compact/src/New/Seesaw/AdmissibleEvaluation.lean\#L167-L215}{{\leanfont isAdmissibleEval\_finEval}}}
\expandafter\newcommand\csname LN@integralthetaeqlift\endcsname{\href{https://github.com/lsernau/grh-lean/blob/72fead102f663d5c7a407964a13dc0cf9d977207/grh-compact/src/New/Seesaw/DoubleLiftValue.lean\#L27-L53}{{\leanfont integral\_theta\_eq\_lift₂}}}
\expandafter\newcommand\csname LN@integralconstTermAtTmul\endcsname{\href{https://github.com/lsernau/grh-lean/blob/72fead102f663d5c7a407964a13dc0cf9d977207/grh-compact/src/New/Seesaw/DoubleLiftValue.lean\#L58-L118}{{\leanfont ArrowPair.\allowbreak{}integral\_constTermAtT\_mul}}}
\expandafter\newcommand\csname LN@zeroTermFin\endcsname{\href{https://github.com/lsernau/grh-lean/blob/72fead102f663d5c7a407964a13dc0cf9d977207/grh-compact/src/New/Seesaw/EdgePair.lean\#L24-L27}{{\leanfont zeroTermFin}}}
\expandafter\newcommand\csname LN@zeroTermFinspec\endcsname{\href{https://github.com/lsernau/grh-lean/blob/72fead102f663d5c7a407964a13dc0cf9d977207/grh-compact/src/New/Seesaw/EdgePair.lean\#L28-L37}{{\leanfont zeroTermFin\_spec}}}
\expandafter\newcommand\csname LN@zeroTermFinfinHeisOpfinFrameVector\endcsname{\href{https://github.com/lsernau/grh-lean/blob/72fead102f663d5c7a407964a13dc0cf9d977207/grh-compact/src/New/Seesaw/EdgePair.lean\#L39-L41}{{\leanfont zeroTermFin\_finHeisOp\_finFrameVector}}}
\expandafter\newcommand\csname LN@zeroTermFinsupport\endcsname{\href{https://github.com/lsernau/grh-lean/blob/72fead102f663d5c7a407964a13dc0cf9d977207/grh-compact/src/New/Seesaw/EdgePair.lean\#L46-L54}{{\leanfont zeroTermFin\_support}}}
\expandafter\newcommand\csname LN@zeroTermFinfinHeisOpvacElem\endcsname{\href{https://github.com/lsernau/grh-lean/blob/72fead102f663d5c7a407964a13dc0cf9d977207/grh-compact/src/New/Seesaw/EdgePair.lean\#L58-L75}{{\leanfont zeroTermFin\_finHeisOp\_vacElem}}}
\expandafter\newcommand\csname LN@GradeZeroPart\endcsname{\href{https://github.com/lsernau/grh-lean/blob/72fead102f663d5c7a407964a13dc0cf9d977207/grh-compact/src/New/Seesaw/EdgePair.lean\#L90-L96}{{\leanfont GradeZeroPart}}}
\expandafter\newcommand\csname LN@ratCoordMatrixsum\endcsname{\href{https://github.com/lsernau/grh-lean/blob/72fead102f663d5c7a407964a13dc0cf9d977207/grh-compact/src/New/Seesaw/EdgePair.lean\#L99-L100}{{\leanfont ratCoordMatrix\_sum}}}
\expandafter\newcommand\csname LN@splitting\endcsname{\href{https://github.com/lsernau/grh-lean/blob/72fead102f663d5c7a407964a13dc0cf9d977207/grh-compact/src/New/Seesaw/EdgePair.lean\#L107-L114}{{\leanfont GradeZeroPart.\allowbreak{}splitting}}}
\expandafter\newcommand\csname LN@hom\endcsname{\href{https://github.com/lsernau/grh-lean/blob/72fead102f663d5c7a407964a13dc0cf9d977207/grh-compact/src/New/Seesaw/EdgePair.lean\#L118-L122}{{\leanfont GradeZeroPart.\allowbreak{}hom}}}
\expandafter\newcommand\csname LN@stackedT\endcsname{\href{https://github.com/lsernau/grh-lean/blob/72fead102f663d5c7a407964a13dc0cf9d977207/grh-compact/src/New/Seesaw/EdgePair.lean\#L126-L126}{{\leanfont stackedT}}}
\expandafter\newcommand\csname LN@stackedTFin\endcsname{\href{https://github.com/lsernau/grh-lean/blob/72fead102f663d5c7a407964a13dc0cf9d977207/grh-compact/src/New/Seesaw/EdgePair.lean\#L130-L130}{{\leanfont stackedTFin}}}
\expandafter\newcommand\csname LN@inclXmulstringMapmulinclXinv\endcsname{\href{https://github.com/lsernau/grh-lean/blob/72fead102f663d5c7a407964a13dc0cf9d977207/grh-compact/src/New/Seesaw/EdgePair.lean\#L133-L135}{{\leanfont inclX\_mul\_stringMap\_mul\_inclX\_inv}}}
\expandafter\newcommand\csname LN@ArrowPair\endcsname{\href{https://github.com/lsernau/grh-lean/blob/72fead102f663d5c7a407964a13dc0cf9d977207/grh-compact/src/New/Seesaw/EdgePair.lean\#L141-L153}{{\leanfont ArrowPair}}}
\expandafter\newcommand\csname LN@homZeroPart\endcsname{\href{https://github.com/lsernau/grh-lean/blob/72fead102f663d5c7a407964a13dc0cf9d977207/grh-compact/src/New/Seesaw/EdgePair.lean\#L159-L160}{{\leanfont ArrowPair.\allowbreak{}homZeroPart}}}
\expandafter\newcommand\csname LN@splittinga\endcsname{\href{https://github.com/lsernau/grh-lean/blob/72fead102f663d5c7a407964a13dc0cf9d977207/grh-compact/src/New/Seesaw/EdgePair.lean\#L164-L164}{{\leanfont ArrowPair.\allowbreak{}splitting}}}
\expandafter\newcommand\csname LN@finEval\endcsname{\href{https://github.com/lsernau/grh-lean/blob/72fead102f663d5c7a407964a13dc0cf9d977207/grh-compact/src/New/Seesaw/EdgePair.lean\#L168-L169}{{\leanfont ArrowPair.\allowbreak{}finEval}}}
\expandafter\newcommand\csname LN@arrowTensorfinEval\endcsname{\href{https://github.com/lsernau/grh-lean/blob/72fead102f663d5c7a407964a13dc0cf9d977207/grh-compact/src/New/Seesaw/EdgePair.lean\#L171-L174}{{\leanfont ArrowPair.\allowbreak{}arrowTensor\_finEval}}}
\expandafter\newcommand\csname LN@zeroTermsplittingtripleWeil\endcsname{\href{https://github.com/lsernau/grh-lean/blob/72fead102f663d5c7a407964a13dc0cf9d977207/grh-compact/src/New/Seesaw/EdgePair.lean\#L179-L207}{{\leanfont ArrowPair.\allowbreak{}zeroTerm\_splitting\_tripleWeil}}}
\expandafter\newcommand\csname LN@finEvalinclXFin\endcsname{\href{https://github.com/lsernau/grh-lean/blob/72fead102f663d5c7a407964a13dc0cf9d977207/grh-compact/src/New/Seesaw/EdgePair.lean\#L211-L225}{{\leanfont ArrowPair.\allowbreak{}finEval\_inclXFin}}}
\expandafter\newcommand\csname LN@middlePair\endcsname{\href{https://github.com/lsernau/grh-lean/blob/72fead102f663d5c7a407964a13dc0cf9d977207/grh-compact/src/New/Seesaw/EdgePair.lean\#L237-L245}{{\leanfont middlePair}}}
\expandafter\newcommand\csname LN@lastPair\endcsname{\href{https://github.com/lsernau/grh-lean/blob/72fead102f663d5c7a407964a13dc0cf9d977207/grh-compact/src/New/Seesaw/EdgePair.lean\#L249-L256}{{\leanfont lastPair}}}
\expandafter\newcommand\csname LN@isUnimodulargradedLatticeBasisgram\endcsname{\href{https://github.com/lsernau/grh-lean/blob/72fead102f663d5c7a407964a13dc0cf9d977207/grh-compact/src/New/Seesaw/FramesForY.lean\#L15-L24}{{\leanfont isUnimodular\_gradedLatticeBasis\_gram}}}
\expandafter\newcommand\csname LN@existsframegradedGram\endcsname{\href{https://github.com/lsernau/grh-lean/blob/72fead102f663d5c7a407964a13dc0cf9d977207/grh-compact/src/New/Seesaw/FramesForY.lean\#L29-L36}{{\leanfont exists\_frame\_gradedGram}}}
\expandafter\newcommand\csname LN@divmemintegralAway\endcsname{\href{https://github.com/lsernau/grh-lean/blob/72fead102f663d5c7a407964a13dc0cf9d977207/grh-compact/src/New/Seesaw/IntegralAway.lean\#L21-L23}{{\leanfont div\_mem\_integralAway}}}
\expandafter\newcommand\csname LN@localLevelleintegralAboveofcoprime\endcsname{\href{https://github.com/lsernau/grh-lean/blob/72fead102f663d5c7a407964a13dc0cf9d977207/grh-compact/src/New/Seesaw/IntegralAway.lean\#L26-L27}{{\leanfont localLevel\_le\_integralAbove\_of\_coprime}}}
\expandafter\newcommand\csname LN@normfinMultipliereqone\endcsname{\href{https://github.com/lsernau/grh-lean/blob/72fead102f663d5c7a407964a13dc0cf9d977207/grh-compact/src/New/Seesaw/IntegralAway.lean\#L29-L32}{{\leanfont Parameters.\allowbreak{}norm\_finMultiplier\_eq\_one}}}
\expandafter\newcommand\csname LN@normpadicdet\endcsname{\href{https://github.com/lsernau/grh-lean/blob/72fead102f663d5c7a407964a13dc0cf9d977207/grh-compact/src/New/Seesaw/IntegralAway.lean\#L34-L36}{{\leanfont AutChar.\allowbreak{}HasWeight.\allowbreak{}norm\_padic\_det}}}
\expandafter\newcommand\csname LN@gradedLatticeBasis\endcsname{\href{https://github.com/lsernau/grh-lean/blob/72fead102f663d5c7a407964a13dc0cf9d977207/grh-compact/src/New/Seesaw/LatticeAndTorsionDefs.lean\#L16-L19}{{\leanfont gradedLatticeBasis}}}
\expandafter\newcommand\csname LN@SpansLattice\endcsname{\href{https://github.com/lsernau/grh-lean/blob/72fead102f663d5c7a407964a13dc0cf9d977207/grh-compact/src/New/Seesaw/LatticeAndTorsionDefs.lean\#L25-L26}{{\leanfont Grading.\allowbreak{}SpansLattice}}}
\expandafter\newcommand\csname LN@integralAway\endcsname{\href{https://github.com/lsernau/grh-lean/blob/72fead102f663d5c7a407964a13dc0cf9d977207/grh-compact/src/New/Seesaw/LatticeAndTorsionDefs.lean\#L30-L31}{{\leanfont integralAway}}}
\expandafter\newcommand\csname LN@IsAdmissibleCoeffs\endcsname{\href{https://github.com/lsernau/grh-lean/blob/72fead102f663d5c7a407964a13dc0cf9d977207/grh-compact/src/New/Seesaw/LatticeAndTorsionDefs.lean\#L35-L38}{{\leanfont IsAdmissibleCoeffs}}}
\expandafter\newcommand\csname LN@weightedTranslations\endcsname{\href{https://github.com/lsernau/grh-lean/blob/72fead102f663d5c7a407964a13dc0cf9d977207/grh-compact/src/New/Seesaw/LatticeAndTorsionDefs.lean\#L46-L48}{{\leanfont weightedTranslations}}}
\expandafter\newcommand\csname LN@translations\endcsname{\href{https://github.com/lsernau/grh-lean/blob/72fead102f663d5c7a407964a13dc0cf9d977207/grh-compact/src/New/Seesaw/LatticeAndTorsionDefs.lean\#L52-L52}{{\leanfont translations}}}
\expandafter\newcommand\csname LN@IsSplitSquarefree\endcsname{\href{https://github.com/lsernau/grh-lean/blob/72fead102f663d5c7a407964a13dc0cf9d977207/grh-compact/src/New/Seesaw/LatticeAndTorsionDefs.lean\#L58-L59}{{\leanfont IsSplitSquarefree}}}
\expandafter\newcommand\csname LN@isAdmissibleCombweilvacElem\endcsname{\href{https://github.com/lsernau/grh-lean/blob/72fead102f663d5c7a407964a13dc0cf9d977207/grh-compact/src/New/Seesaw/MovedVacuumComb.lean\#L12-L22}{{\leanfont isAdmissibleComb\_weil\_vacElem}}}
\expandafter\newcommand\csname LN@inFinFin\endcsname{\href{https://github.com/lsernau/grh-lean/blob/72fead102f663d5c7a407964a13dc0cf9d977207/grh-compact/src/Lib/Weil/Construction/Projections.lean\#L32-L33}{{\leanfont Space.\allowbreak{}finHeisOp\_finHeisOp}}}
\expandafter\newcommand\csname LN@inPsiFin\endcsname{\href{https://github.com/lsernau/grh-lean/blob/72fead102f663d5c7a407964a13dc0cf9d977207/grh-compact/src/Lib/Adelic/Adeles/AdditiveCharBasic.lean\#L15-L16}{{\leanfont isOfFinOrder\_ψf}}}
\expandafter\newcommand\csname LN@inFinVac\endcsname{\href{https://github.com/lsernau/grh-lean/blob/72fead102f663d5c7a407964a13dc0cf9d977207/grh-compact/src/Lib/Weil/Construction/Projections.lean\#L131-L134}{{\leanfont Space.\allowbreak{}finHeisOp\_vacElem}}}
\expandafter\newcommand\csname LN@inFinMem\endcsname{\href{https://github.com/lsernau/grh-lean/blob/72fead102f663d5c7a407964a13dc0cf9d977207/grh-compact/src/Lib/Weil/Construction/FinitePart.lean\#L22-L23}{{\leanfont Space.\allowbreak{}finHeisOp\_mem\_finitePart}}}
\expandafter\newcommand\csname LN@inOmFin\endcsname{\href{https://github.com/lsernau/grh-lean/blob/72fead102f663d5c7a407964a13dc0cf9d977207/grh-compact/src/Lib/Weil/Construction/Projections.lean\#L61-L62}{{\leanfont Space.\allowbreak{}omega\_finHeisOp}}}
\expandafter\newcommand\csname LN@inOmInt\endcsname{\href{https://github.com/lsernau/grh-lean/blob/72fead102f663d5c7a407964a13dc0cf9d977207/grh-compact/src/Lib/Weil/Construction/Projections.lean\#L91-L96}{{\leanfont Space.\allowbreak{}omega\_integrated}}}
\expandafter\newcommand\csname LN@inRhoInt\endcsname{\href{https://github.com/lsernau/grh-lean/blob/72fead102f663d5c7a407964a13dc0cf9d977207/grh-compact/src/Lib/Weil/Construction/Projections.lean\#L53-L55}{{\leanfont Space.\allowbreak{}exists\_rho\_integrated\_eq}}}
\expandafter\newcommand\csname LN@inVacSum\endcsname{\href{https://github.com/lsernau/grh-lean/blob/72fead102f663d5c7a407964a13dc0cf9d977207/grh-compact/src/Lib/Weil/Construction/Projections.lean\#L215-L220}{{\leanfont Space.\allowbreak{}omega\_vacElem\_sum}}}
\expandafter\newcommand\csname LN@inNormSq\endcsname{\href{https://github.com/lsernau/grh-lean/blob/72fead102f663d5c7a407964a13dc0cf9d977207/grh-compact/src/Lib/Weil/Construction/FiniteMultiplier.lean\#L181-L185}{{\leanfont norm\_finMultiplier\_sq}}}
\expandafter\newcommand\csname LN@inSmulVac\endcsname{\href{https://github.com/lsernau/grh-lean/blob/72fead102f663d5c7a407964a13dc0cf9d977207/grh-compact/src/Lib/Weil/Construction/FinitePart.lean\#L57-L61}{{\leanfont Space.\allowbreak{}exists\_eq\_smul\_omega\_vacElem}}}
\expandafter\newcommand\csname LN@inDual\endcsname{\href{https://github.com/lsernau/grh-lean/blob/72fead102f663d5c7a407964a13dc0cf9d977207/grh-compact/src/Lib/Weil/Construction/FiniteVacuum.lean\#L62-L64}{{\leanfont stdLattice.\allowbreak{}mem\_stdLattice\_iff\_psi₀}}}
\expandafter\newcommand\csname LN@inProj\endcsname{\href{https://github.com/lsernau/grh-lean/blob/72fead102f663d5c7a407964a13dc0cf9d977207/grh-compact/src/Lib/Weil/Construction/Projections.lean\#L149-L152}{{\leanfont Space.\allowbreak{}projOn\_finHeisOp}}}
\expandafter\newcommand\csname LN@inThArrow\endcsname{\href{https://github.com/lsernau/grh-lean/blob/72fead102f663d5c7a407964a13dc0cf9d977207/grh-compact/src/Lib/Weil/DualPair/TwoEdgesOnePair.lean\#L112-L115}{{\leanfont theta\_arrowProduct}}}
\expandafter\newcommand\csname LN@inHomArrow\endcsname{\href{https://github.com/lsernau/grh-lean/blob/72fead102f663d5c7a407964a13dc0cf9d977207/grh-compact/src/Lib/Weil/DualPair/TwoEdgesOnePair.lean\#L117-L121}{{\leanfont homHeisOp\_arrowProduct}}}
\expandafter\newcommand\csname LN@inFinArrow\endcsname{\href{https://github.com/lsernau/grh-lean/blob/72fead102f663d5c7a407964a13dc0cf9d977207/grh-compact/src/Lib/Weil/DualPair/TwoEdgesOnePair.lean\#L123-L128}{{\leanfont finHeisOp\_arrowProduct}}}
\expandafter\newcommand\csname LN@inArrowPair\endcsname{\href{https://github.com/lsernau/grh-lean/blob/72fead102f663d5c7a407964a13dc0cf9d977207/grh-compact/src/Lib/Weil/DualPair/TwoEdgesOnePair.lean\#L133-L135}{{\leanfont arrowProduct\_twistedPairWeil}}}
\expandafter\newcommand\csname LN@inTensorLift\endcsname{\href{https://github.com/lsernau/grh-lean/blob/72fead102f663d5c7a407964a13dc0cf9d977207/grh-compact/src/Lib/Weil/DualPair/TwoEdgesOnePair.lean\#L151-L153}{{\leanfont arrowTensor\_map\_liftWeilFin}}}
\expandafter\newcommand\csname LN@inLeftInv\endcsname{\href{https://github.com/lsernau/grh-lean/blob/72fead102f663d5c7a407964a13dc0cf9d977207/grh-compact/src/Lib/Weil/DualPair/TwoEdgesOnePair.lean\#L137-L143}{{\leanfont leftInverse\_arrowTensor}}}
\expandafter\newcommand\csname LN@inRange\endcsname{\href{https://github.com/lsernau/grh-lean/blob/72fead102f663d5c7a407964a13dc0cf9d977207/grh-compact/src/Lib/Weil/DualPair/TwoEdgesOnePair.lean\#L145-L149}{{\leanfont range\_arrowTensor}}}
\expandafter\newcommand\csname LN@inThRat\endcsname{\href{https://github.com/lsernau/grh-lean/blob/72fead102f663d5c7a407964a13dc0cf9d977207/grh-compact/src/Lib/Weil/DualPair/TwoEdgesOnePair.lean\#L155-L158}{{\leanfont theta\_tripleWeil\_rat\_mul}}}
\expandafter\newcommand\csname LN@inPairHeis\endcsname{\href{https://github.com/lsernau/grh-lean/blob/72fead102f663d5c7a407964a13dc0cf9d977207/grh-compact/src/Lib/Weil/DualPair/TwoEdgesOnePair.lean\#L32-L34}{{\leanfont pairWeil\_homHeisOp}}}
\expandafter\newcommand\csname LN@inAdapted\endcsname{\href{https://github.com/lsernau/grh-lean/blob/72fead102f663d5c7a407964a13dc0cf9d977207/grh-compact/src/Lib/Weil/MixedModel/ZeroTerm.lean\#L272-L288}{{\leanfont Space.\allowbreak{}Splitting.\allowbreak{}exists\_isAdapted\_and\_frame\_eq}}}
\expandafter\newcommand\csname LN@inMixedVac\endcsname{\href{https://github.com/lsernau/grh-lean/blob/72fead102f663d5c7a407964a13dc0cf9d977207/grh-compact/src/Lib/Weil/MixedModel/ZeroTerm.lean\#L615-L619}{{\leanfont Space.\allowbreak{}Splitting.\allowbreak{}IsAdapted.\allowbreak{}mixedTerm\_omega\_vacElem}}}
\expandafter\newcommand\csname LN@inRealZero\endcsname{\href{https://github.com/lsernau/grh-lean/blob/72fead102f663d5c7a407964a13dc0cf9d977207/grh-compact/src/Lib/Weil/MixedModel/ZeroTerm.lean\#L395-L405}{{\leanfont Space.\allowbreak{}Splitting.\allowbreak{}IsAdapted.\allowbreak{}exists\_real\_zeroTerm}}}
\expandafter\newcommand\csname LN@inZeroRho\endcsname{\href{https://github.com/lsernau/grh-lean/blob/72fead102f663d5c7a407964a13dc0cf9d977207/grh-compact/src/Lib/Weil/MixedModel/ZeroTerm.lean\#L567-L568}{{\leanfont Space.\allowbreak{}Splitting.\allowbreak{}zeroTerm\_rho\_iAdelic}}}
\expandafter\newcommand\csname LN@inFrameVec\endcsname{\href{https://github.com/lsernau/grh-lean/blob/72fead102f663d5c7a407964a13dc0cf9d977207/grh-compact/src/Lib/Weil/MixedModel/ZeroTerm.lean\#L446-L449}{{\leanfont Space.\allowbreak{}Splitting.\allowbreak{}exists\_eq\_finFrameVector}}}
\expandafter\newcommand\csname LN@inFrameAdj\endcsname{\href{https://github.com/lsernau/grh-lean/blob/72fead102f663d5c7a407964a13dc0cf9d977207/grh-compact/src/Lib/Weil/MixedModel/ZeroTerm.lean\#L441-L443}{{\leanfont Space.\allowbreak{}Splitting.\allowbreak{}adjoint\_mulVec\_finFrameVector}}}
\expandafter\newcommand\csname LN@inZeroIncl\endcsname{\href{https://github.com/lsernau/grh-lean/blob/72fead102f663d5c7a407964a13dc0cf9d977207/grh-compact/src/Lib/Weil/DualPair/TwoEdgesOnePair.lean\#L58-L68}{{\leanfont zeroTerm\_omega\_of\_inclMatrix}}}
\expandafter\newcommand\csname LN@inHomTarget\endcsname{\href{https://github.com/lsernau/grh-lean/blob/72fead102f663d5c7a407964a13dc0cf9d977207/grh-compact/src/Lib/Weil/DualPair/PairSplitting.lean\#L314-L319}{{\leanfont Space.\allowbreak{}Splitting.\allowbreak{}homTarget\_zeroTerm\_pairWeil}}}
\expandafter\newcommand\csname LN@inNameBox\endcsname{\href{https://github.com/lsernau/grh-lean/blob/72fead102f663d5c7a407964a13dc0cf9d977207/grh-compact/src/Lib/Weil/MixedModel/Names.lean\#L91-L92}{{\leanfont Space.\allowbreak{}name\_boxIndicator}}}
\expandafter\newcommand\csname LN@inNameTr\endcsname{\href{https://github.com/lsernau/grh-lean/blob/72fead102f663d5c7a407964a13dc0cf9d977207/grh-compact/src/Lib/Weil/MixedModel/Names.lean\#L85-L87}{{\leanfont Space.\allowbreak{}name\_schrodingerAct\_boxIndicator}}}
\expandafter\newcommand\csname LN@inNameOm\endcsname{\href{https://github.com/lsernau/grh-lean/blob/72fead102f663d5c7a407964a13dc0cf9d977207/grh-compact/src/Lib/Weil/MixedModel/NameWeilAction.lean\#L194-L202}{{\leanfont Space.\allowbreak{}ω\_name\_on}}}
\expandafter\newcommand\csname LN@inBesselMean\endcsname{\href{https://github.com/lsernau/grh-lean/blob/72fead102f663d5c7a407964a13dc0cf9d977207/grh-compact/src/Lib/Weil/Periods/BesselIntegral.lean\#L112-L144}{{\leanfont Grading.\allowbreak{}integral\_besselMeasure\_eq\_of\_mean}}}
\expandafter\newcommand\csname LN@inGenConst\endcsname{\href{https://github.com/lsernau/grh-lean/blob/72fead102f663d5c7a407964a13dc0cf9d977207/grh-compact/src/Lib/Weil/Periods/UnfoldingStep.lean\#L254-L261}{{\leanfont Grading.\allowbreak{}IsGeneric.\allowbreak{}unipotent\_mul\_constTerm}}}
\expandafter\newcommand\csname LN@inConstRat\endcsname{\href{https://github.com/lsernau/grh-lean/blob/72fead102f663d5c7a407964a13dc0cf9d977207/grh-compact/src/Lib/Weil/Periods/UnfoldingStep.lean\#L265-L269}{{\leanfont Grading.\allowbreak{}constTermAtT\_rat\_mul}}}
\expandafter\newcommand\csname LN@inAutCoef\endcsname{\href{https://github.com/lsernau/grh-lean/blob/72fead102f663d5c7a407964a13dc0cf9d977207/grh-compact/src/Lib/Adelic/CuspForms/AlgebraicForms.lean\#L42-L47}{{\leanfont Space.\allowbreak{}isAutomorphic\_algFormCoeff}}}
\expandafter\newcommand\csname LN@inHermite\endcsname{\href{https://github.com/lsernau/grh-lean/blob/72fead102f663d5c7a407964a13dc0cf9d977207/grh-compact/src/Lib/Adelic/Witt/FrameAction.lean\#L74-L80}{{\leanfont IsValuationPair.\allowbreak{}exists\_isPrimitiveMatrix\_mul}}}
\expandafter\newcommand\csname LN@inDense\endcsname{\href{https://github.com/lsernau/grh-lean/blob/72fead102f663d5c7a407964a13dc0cf9d977207/grh-compact/src/Lib/Weil/Fock/FockHarmonicCoeff.lean\#L70-L75}{{\leanfont denseRange\_polyVector}}}
\expandafter\newcommand\csname LN@inPolyProd\endcsname{\href{https://github.com/lsernau/grh-lean/blob/72fead102f663d5c7a407964a13dc0cf9d977207/grh-compact/src/Lib/General/LinearAlgebra/SmoothMatrices.lean\#L89-L92}{{\leanfont mem\_rename\_range\_mul}}}
\expandafter\newcommand\csname LN@inMeans\endcsname{\href{https://github.com/lsernau/grh-lean/blob/72fead102f663d5c7a407964a13dc0cf9d977207/grh-compact/src/Lib/Adelic/QuotientMeasure/PeriodicMean.lean\#L60-L68}{{\leanfont quotMean\_quotMean}}}
\expandafter\newcommand\csname LN@inCoprime\endcsname{\href{https://github.com/lsernau/grh-lean/blob/72fead102f663d5c7a407964a13dc0cf9d977207/grh-compact/src/Lib/General/Algebra/AdicRingHomExt.lean\#L102-L103}{{\leanfont Nat.\allowbreak{}Coprime.\allowbreak{}norm\_natCast\_eq\_one}}}
\newcommand{\ConcordanceRows}{\crow{IsAdmissibleComb}{\href{https://github.com/lsernau/grh-lean/blob/72fead102f663d5c7a407964a13dc0cf9d977207/grh-compact/src/New/Seesaw/AdmissibleComb.lean\#L12-L13}{{\leanfont IsAdmissibleComb}}}{AdmissibleComb.lean}{12--13}
\crow{ofsum}{\href{https://github.com/lsernau/grh-lean/blob/72fead102f663d5c7a407964a13dc0cf9d977207/grh-compact/src/New/Seesaw/AdmissibleComb.lean\#L20-L22}{{\leanfont IsAdmissibleComb.\allowbreak{}of\_sum}}}{AdmissibleComb.lean}{20--22}
\crow{single}{\href{https://github.com/lsernau/grh-lean/blob/72fead102f663d5c7a407964a13dc0cf9d977207/grh-compact/src/New/Seesaw/AdmissibleComb.lean\#L24-L25}{{\leanfont IsAdmissibleComb.\allowbreak{}single}}}{AdmissibleComb.lean}{24--25}
\crow{bind}{\href{https://github.com/lsernau/grh-lean/blob/72fead102f663d5c7a407964a13dc0cf9d977207/grh-compact/src/New/Seesaw/AdmissibleComb.lean\#L32-L40}{{\leanfont IsAdmissibleComb.\allowbreak{}bind}}}{AdmissibleComb.lean}{32--40}
\crow{map}{\href{https://github.com/lsernau/grh-lean/blob/72fead102f663d5c7a407964a13dc0cf9d977207/grh-compact/src/New/Seesaw/AdmissibleComb.lean\#L42-L42}{{\leanfont IsAdmissibleComb.\allowbreak{}map}}}{AdmissibleComb.lean}{42--42}
\crow{smul}{\href{https://github.com/lsernau/grh-lean/blob/72fead102f663d5c7a407964a13dc0cf9d977207/grh-compact/src/New/Seesaw/AdmissibleComb.lean\#L44-L46}{{\leanfont IsAdmissibleComb.\allowbreak{}smul}}}{AdmissibleComb.lean}{44--46}
\crow{smulofeqzeroorisOfFinOrder}{\href{https://github.com/lsernau/grh-lean/blob/72fead102f663d5c7a407964a13dc0cf9d977207/grh-compact/src/New/Seesaw/AdmissibleComb.lean\#L48-L50}{{\leanfont IsAdmissibleComb.\allowbreak{}smul\_of\_eq\_zero\_or\_isOfFinOrder}}}{AdmissibleComb.lean}{48--50}
\crow{ofeqisOfFinOrdersmul}{\href{https://github.com/lsernau/grh-lean/blob/72fead102f663d5c7a407964a13dc0cf9d977207/grh-compact/src/New/Seesaw/AdmissibleComb.lean\#L52-L54}{{\leanfont IsAdmissibleComb.\allowbreak{}of\_eq\_isOfFinOrder\_smul}}}{AdmissibleComb.lean}{52--54}
\crow{mapbilinear}{\href{https://github.com/lsernau/grh-lean/blob/72fead102f663d5c7a407964a13dc0cf9d977207/grh-compact/src/New/Seesaw/AdmissibleComb.lean\#L56-L59}{{\leanfont IsAdmissibleComb.\allowbreak{}map\_bilinear}}}{AdmissibleComb.lean}{56--59}
\crow{smulmem}{\href{https://github.com/lsernau/grh-lean/blob/72fead102f663d5c7a407964a13dc0cf9d977207/grh-compact/src/New/Seesaw/AdmissibleComb.lean\#L64-L67}{{\leanfont IsAdmissibleComb.\allowbreak{}smul\_mem}}}{AdmissibleComb.lean}{64--67}
\crow{pure}{\href{https://github.com/lsernau/grh-lean/blob/72fead102f663d5c7a407964a13dc0cf9d977207/grh-compact/src/New/Seesaw/AdmissibleComb.lean\#L71-L77}{{\leanfont IsAdmissibleComb.\allowbreak{}pure}}}{AdmissibleComb.lean}{71--77}
\crow{normUnitsmul}{\href{https://github.com/lsernau/grh-lean/blob/72fead102f663d5c7a407964a13dc0cf9d977207/grh-compact/src/New/Seesaw/AdmissibleComb.lean\#L79-L83}{{\leanfont IsAdmissibleComb.\allowbreak{}normUnit\_smul}}}{AdmissibleComb.lean}{79--83}
\crow{selfOrConj}{\href{https://github.com/lsernau/grh-lean/blob/72fead102f663d5c7a407964a13dc0cf9d977207/grh-compact/src/New/Seesaw/AdmissibleEvaluation.lean\#L20-L20}{{\leanfont selfOrConj}}}{AdmissibleEvaluation.lean}{20--20}
\crow{mulVecmemweightedTranslations}{\href{https://github.com/lsernau/grh-lean/blob/72fead102f663d5c7a407964a13dc0cf9d977207/grh-compact/src/New/Seesaw/AdmissibleEvaluation.lean\#L22-L26}{{\leanfont mulVec\_mem\_weightedTranslations}}}{AdmissibleEvaluation.lean}{22--26}
\crow{starmemweightedTranslations}{\href{https://github.com/lsernau/grh-lean/blob/72fead102f663d5c7a407964a13dc0cf9d977207/grh-compact/src/New/Seesaw/AdmissibleEvaluation.lean\#L28-L30}{{\leanfont star\_mem\_weightedTranslations}}}{AdmissibleEvaluation.lean}{28--30}
\crow{invxFinmulVecmemofspansLattice}{\href{https://github.com/lsernau/grh-lean/blob/72fead102f663d5c7a407964a13dc0cf9d977207/grh-compact/src/New/Seesaw/AdmissibleEvaluation.lean\#L32-L42}{{\leanfont inv\_xFin\_mulVec\_mem\_of\_spansLattice}}}{AdmissibleEvaluation.lean}{32--42}
\crow{IsNeighbourVector}{\href{https://github.com/lsernau/grh-lean/blob/72fead102f663d5c7a407964a13dc0cf9d977207/grh-compact/src/New/Seesaw/AdmissibleEvaluation.lean\#L60-L62}{{\leanfont IsNeighbourVector}}}{AdmissibleEvaluation.lean}{60--62}
\crow{IsAdmissibleEval}{\href{https://github.com/lsernau/grh-lean/blob/72fead102f663d5c7a407964a13dc0cf9d977207/grh-compact/src/New/Seesaw/AdmissibleEvaluation.lean\#L67-L72}{{\leanfont IsAdmissibleEval}}}{AdmissibleEvaluation.lean}{67--72}
\crow{existsisAdmissibleEvalstd}{\href{https://github.com/lsernau/grh-lean/blob/72fead102f663d5c7a407964a13dc0cf9d977207/grh-compact/src/New/Seesaw/AdmissibleEvaluation.lean\#L81-L160}{{\leanfont exists\_isAdmissibleEval\_std}}}{AdmissibleEvaluation.lean}{81--160}
\crow{isAdmissibleEvalfinEval}{\href{https://github.com/lsernau/grh-lean/blob/72fead102f663d5c7a407964a13dc0cf9d977207/grh-compact/src/New/Seesaw/AdmissibleEvaluation.lean\#L167-L215}{{\leanfont isAdmissibleEval\_finEval}}}{AdmissibleEvaluation.lean}{167--215}
\crow{integralthetaeqlift}{\href{https://github.com/lsernau/grh-lean/blob/72fead102f663d5c7a407964a13dc0cf9d977207/grh-compact/src/New/Seesaw/DoubleLiftValue.lean\#L27-L53}{{\leanfont integral\_theta\_eq\_lift₂}}}{DoubleLiftValue.lean}{27--53}
\crow{integralconstTermAtTmul}{\href{https://github.com/lsernau/grh-lean/blob/72fead102f663d5c7a407964a13dc0cf9d977207/grh-compact/src/New/Seesaw/DoubleLiftValue.lean\#L58-L118}{{\leanfont ArrowPair.\allowbreak{}integral\_constTermAtT\_mul}}}{DoubleLiftValue.lean}{58--118}
\crow{zeroTermFin}{\href{https://github.com/lsernau/grh-lean/blob/72fead102f663d5c7a407964a13dc0cf9d977207/grh-compact/src/New/Seesaw/EdgePair.lean\#L24-L27}{{\leanfont zeroTermFin}}}{EdgePair.lean}{24--27}
\crow{zeroTermFinspec}{\href{https://github.com/lsernau/grh-lean/blob/72fead102f663d5c7a407964a13dc0cf9d977207/grh-compact/src/New/Seesaw/EdgePair.lean\#L28-L37}{{\leanfont zeroTermFin\_spec}}}{EdgePair.lean}{28--37}
\crow{zeroTermFinfinHeisOpfinFrameVector}{\href{https://github.com/lsernau/grh-lean/blob/72fead102f663d5c7a407964a13dc0cf9d977207/grh-compact/src/New/Seesaw/EdgePair.lean\#L39-L41}{{\leanfont zeroTermFin\_finHeisOp\_finFrameVector}}}{EdgePair.lean}{39--41}
\crow{zeroTermFinsupport}{\href{https://github.com/lsernau/grh-lean/blob/72fead102f663d5c7a407964a13dc0cf9d977207/grh-compact/src/New/Seesaw/EdgePair.lean\#L46-L54}{{\leanfont zeroTermFin\_support}}}{EdgePair.lean}{46--54}
\crow{zeroTermFinfinHeisOpvacElem}{\href{https://github.com/lsernau/grh-lean/blob/72fead102f663d5c7a407964a13dc0cf9d977207/grh-compact/src/New/Seesaw/EdgePair.lean\#L58-L75}{{\leanfont zeroTermFin\_finHeisOp\_vacElem}}}{EdgePair.lean}{58--75}
\crow{GradeZeroPart}{\href{https://github.com/lsernau/grh-lean/blob/72fead102f663d5c7a407964a13dc0cf9d977207/grh-compact/src/New/Seesaw/EdgePair.lean\#L90-L96}{{\leanfont GradeZeroPart}}}{EdgePair.lean}{90--96}
\crow{ratCoordMatrixsum}{\href{https://github.com/lsernau/grh-lean/blob/72fead102f663d5c7a407964a13dc0cf9d977207/grh-compact/src/New/Seesaw/EdgePair.lean\#L99-L100}{{\leanfont ratCoordMatrix\_sum}}}{EdgePair.lean}{99--100}
\crow{splitting}{\href{https://github.com/lsernau/grh-lean/blob/72fead102f663d5c7a407964a13dc0cf9d977207/grh-compact/src/New/Seesaw/EdgePair.lean\#L107-L114}{{\leanfont GradeZeroPart.\allowbreak{}splitting}}}{EdgePair.lean}{107--114}
\crow{hom}{\href{https://github.com/lsernau/grh-lean/blob/72fead102f663d5c7a407964a13dc0cf9d977207/grh-compact/src/New/Seesaw/EdgePair.lean\#L118-L122}{{\leanfont GradeZeroPart.\allowbreak{}hom}}}{EdgePair.lean}{118--122}
\crow{stackedT}{\href{https://github.com/lsernau/grh-lean/blob/72fead102f663d5c7a407964a13dc0cf9d977207/grh-compact/src/New/Seesaw/EdgePair.lean\#L126-L126}{{\leanfont stackedT}}}{EdgePair.lean}{126--126}
\crow{stackedTFin}{\href{https://github.com/lsernau/grh-lean/blob/72fead102f663d5c7a407964a13dc0cf9d977207/grh-compact/src/New/Seesaw/EdgePair.lean\#L130-L130}{{\leanfont stackedTFin}}}{EdgePair.lean}{130--130}
\crow{inclXmulstringMapmulinclXinv}{\href{https://github.com/lsernau/grh-lean/blob/72fead102f663d5c7a407964a13dc0cf9d977207/grh-compact/src/New/Seesaw/EdgePair.lean\#L133-L135}{{\leanfont inclX\_mul\_stringMap\_mul\_inclX\_inv}}}{EdgePair.lean}{133--135}
\crow{ArrowPair}{\href{https://github.com/lsernau/grh-lean/blob/72fead102f663d5c7a407964a13dc0cf9d977207/grh-compact/src/New/Seesaw/EdgePair.lean\#L141-L153}{{\leanfont ArrowPair}}}{EdgePair.lean}{141--153}
\crow{homZeroPart}{\href{https://github.com/lsernau/grh-lean/blob/72fead102f663d5c7a407964a13dc0cf9d977207/grh-compact/src/New/Seesaw/EdgePair.lean\#L159-L160}{{\leanfont ArrowPair.\allowbreak{}homZeroPart}}}{EdgePair.lean}{159--160}
\crow{splittinga}{\href{https://github.com/lsernau/grh-lean/blob/72fead102f663d5c7a407964a13dc0cf9d977207/grh-compact/src/New/Seesaw/EdgePair.lean\#L164-L164}{{\leanfont ArrowPair.\allowbreak{}splitting}}}{EdgePair.lean}{164--164}
\crow{finEval}{\href{https://github.com/lsernau/grh-lean/blob/72fead102f663d5c7a407964a13dc0cf9d977207/grh-compact/src/New/Seesaw/EdgePair.lean\#L168-L169}{{\leanfont ArrowPair.\allowbreak{}finEval}}}{EdgePair.lean}{168--169}
\crow{arrowTensorfinEval}{\href{https://github.com/lsernau/grh-lean/blob/72fead102f663d5c7a407964a13dc0cf9d977207/grh-compact/src/New/Seesaw/EdgePair.lean\#L171-L174}{{\leanfont ArrowPair.\allowbreak{}arrowTensor\_finEval}}}{EdgePair.lean}{171--174}
\crow{zeroTermsplittingtripleWeil}{\href{https://github.com/lsernau/grh-lean/blob/72fead102f663d5c7a407964a13dc0cf9d977207/grh-compact/src/New/Seesaw/EdgePair.lean\#L179-L207}{{\leanfont ArrowPair.\allowbreak{}zeroTerm\_splitting\_tripleWeil}}}{EdgePair.lean}{179--207}
\crow{finEvalinclXFin}{\href{https://github.com/lsernau/grh-lean/blob/72fead102f663d5c7a407964a13dc0cf9d977207/grh-compact/src/New/Seesaw/EdgePair.lean\#L211-L225}{{\leanfont ArrowPair.\allowbreak{}finEval\_inclXFin}}}{EdgePair.lean}{211--225}
\crow{middlePair}{\href{https://github.com/lsernau/grh-lean/blob/72fead102f663d5c7a407964a13dc0cf9d977207/grh-compact/src/New/Seesaw/EdgePair.lean\#L237-L245}{{\leanfont middlePair}}}{EdgePair.lean}{237--245}
\crow{lastPair}{\href{https://github.com/lsernau/grh-lean/blob/72fead102f663d5c7a407964a13dc0cf9d977207/grh-compact/src/New/Seesaw/EdgePair.lean\#L249-L256}{{\leanfont lastPair}}}{EdgePair.lean}{249--256}
\crow{isUnimodulargradedLatticeBasisgram}{\href{https://github.com/lsernau/grh-lean/blob/72fead102f663d5c7a407964a13dc0cf9d977207/grh-compact/src/New/Seesaw/FramesForY.lean\#L15-L24}{{\leanfont isUnimodular\_gradedLatticeBasis\_gram}}}{FramesForY.lean}{15--24}
\crow{existsframegradedGram}{\href{https://github.com/lsernau/grh-lean/blob/72fead102f663d5c7a407964a13dc0cf9d977207/grh-compact/src/New/Seesaw/FramesForY.lean\#L29-L36}{{\leanfont exists\_frame\_gradedGram}}}{FramesForY.lean}{29--36}
\crow{divmemintegralAway}{\href{https://github.com/lsernau/grh-lean/blob/72fead102f663d5c7a407964a13dc0cf9d977207/grh-compact/src/New/Seesaw/IntegralAway.lean\#L21-L23}{{\leanfont div\_mem\_integralAway}}}{IntegralAway.lean}{21--23}
\crow{localLevelleintegralAboveofcoprime}{\href{https://github.com/lsernau/grh-lean/blob/72fead102f663d5c7a407964a13dc0cf9d977207/grh-compact/src/New/Seesaw/IntegralAway.lean\#L26-L27}{{\leanfont localLevel\_le\_integralAbove\_of\_coprime}}}{IntegralAway.lean}{26--27}
\crow{normfinMultipliereqone}{\href{https://github.com/lsernau/grh-lean/blob/72fead102f663d5c7a407964a13dc0cf9d977207/grh-compact/src/New/Seesaw/IntegralAway.lean\#L29-L32}{{\leanfont Parameters.\allowbreak{}norm\_finMultiplier\_eq\_one}}}{IntegralAway.lean}{29--32}
\crow{normpadicdet}{\href{https://github.com/lsernau/grh-lean/blob/72fead102f663d5c7a407964a13dc0cf9d977207/grh-compact/src/New/Seesaw/IntegralAway.lean\#L34-L36}{{\leanfont AutChar.\allowbreak{}HasWeight.\allowbreak{}norm\_padic\_det}}}{IntegralAway.lean}{34--36}
\crow{gradedLatticeBasis}{\href{https://github.com/lsernau/grh-lean/blob/72fead102f663d5c7a407964a13dc0cf9d977207/grh-compact/src/New/Seesaw/LatticeAndTorsionDefs.lean\#L16-L19}{{\leanfont gradedLatticeBasis}}}{LatticeAndTorsionDefs.lean}{16--19}
\crow{SpansLattice}{\href{https://github.com/lsernau/grh-lean/blob/72fead102f663d5c7a407964a13dc0cf9d977207/grh-compact/src/New/Seesaw/LatticeAndTorsionDefs.lean\#L25-L26}{{\leanfont Grading.\allowbreak{}SpansLattice}}}{LatticeAndTorsionDefs.lean}{25--26}
\crow{integralAway}{\href{https://github.com/lsernau/grh-lean/blob/72fead102f663d5c7a407964a13dc0cf9d977207/grh-compact/src/New/Seesaw/LatticeAndTorsionDefs.lean\#L30-L31}{{\leanfont integralAway}}}{LatticeAndTorsionDefs.lean}{30--31}
\crow{IsAdmissibleCoeffs}{\href{https://github.com/lsernau/grh-lean/blob/72fead102f663d5c7a407964a13dc0cf9d977207/grh-compact/src/New/Seesaw/LatticeAndTorsionDefs.lean\#L35-L38}{{\leanfont IsAdmissibleCoeffs}}}{LatticeAndTorsionDefs.lean}{35--38}
\crow{weightedTranslations}{\href{https://github.com/lsernau/grh-lean/blob/72fead102f663d5c7a407964a13dc0cf9d977207/grh-compact/src/New/Seesaw/LatticeAndTorsionDefs.lean\#L46-L48}{{\leanfont weightedTranslations}}}{LatticeAndTorsionDefs.lean}{46--48}
\crow{translations}{\href{https://github.com/lsernau/grh-lean/blob/72fead102f663d5c7a407964a13dc0cf9d977207/grh-compact/src/New/Seesaw/LatticeAndTorsionDefs.lean\#L52-L52}{{\leanfont translations}}}{LatticeAndTorsionDefs.lean}{52--52}
\crow{IsSplitSquarefree}{\href{https://github.com/lsernau/grh-lean/blob/72fead102f663d5c7a407964a13dc0cf9d977207/grh-compact/src/New/Seesaw/LatticeAndTorsionDefs.lean\#L58-L59}{{\leanfont IsSplitSquarefree}}}{LatticeAndTorsionDefs.lean}{58--59}
\crow{isAdmissibleCombweilvacElem}{\href{https://github.com/lsernau/grh-lean/blob/72fead102f663d5c7a407964a13dc0cf9d977207/grh-compact/src/New/Seesaw/MovedVacuumComb.lean\#L12-L22}{{\leanfont isAdmissibleComb\_weil\_vacElem}}}{MovedVacuumComb.lean}{12--22}
\crow{inFinFin}{\href{https://github.com/lsernau/grh-lean/blob/72fead102f663d5c7a407964a13dc0cf9d977207/grh-compact/src/Lib/Weil/Construction/Projections.lean\#L32-L33}{{\leanfont Space.\allowbreak{}finHeisOp\_finHeisOp}}}{Projections.lean}{32--33}
\crow{inPsiFin}{\href{https://github.com/lsernau/grh-lean/blob/72fead102f663d5c7a407964a13dc0cf9d977207/grh-compact/src/Lib/Adelic/Adeles/AdditiveCharBasic.lean\#L15-L16}{{\leanfont isOfFinOrder\_ψf}}}{AdditiveCharBasic.lean}{15--16}
\crow{inFinVac}{\href{https://github.com/lsernau/grh-lean/blob/72fead102f663d5c7a407964a13dc0cf9d977207/grh-compact/src/Lib/Weil/Construction/Projections.lean\#L131-L134}{{\leanfont Space.\allowbreak{}finHeisOp\_vacElem}}}{Projections.lean}{131--134}
\crow{inFinMem}{\href{https://github.com/lsernau/grh-lean/blob/72fead102f663d5c7a407964a13dc0cf9d977207/grh-compact/src/Lib/Weil/Construction/FinitePart.lean\#L22-L23}{{\leanfont Space.\allowbreak{}finHeisOp\_mem\_finitePart}}}{FinitePart.lean}{22--23}
\crow{inOmFin}{\href{https://github.com/lsernau/grh-lean/blob/72fead102f663d5c7a407964a13dc0cf9d977207/grh-compact/src/Lib/Weil/Construction/Projections.lean\#L61-L62}{{\leanfont Space.\allowbreak{}omega\_finHeisOp}}}{Projections.lean}{61--62}
\crow{inOmInt}{\href{https://github.com/lsernau/grh-lean/blob/72fead102f663d5c7a407964a13dc0cf9d977207/grh-compact/src/Lib/Weil/Construction/Projections.lean\#L91-L96}{{\leanfont Space.\allowbreak{}omega\_integrated}}}{Projections.lean}{91--96}
\crow{inRhoInt}{\href{https://github.com/lsernau/grh-lean/blob/72fead102f663d5c7a407964a13dc0cf9d977207/grh-compact/src/Lib/Weil/Construction/Projections.lean\#L53-L55}{{\leanfont Space.\allowbreak{}exists\_rho\_integrated\_eq}}}{Projections.lean}{53--55}
\crow{inVacSum}{\href{https://github.com/lsernau/grh-lean/blob/72fead102f663d5c7a407964a13dc0cf9d977207/grh-compact/src/Lib/Weil/Construction/Projections.lean\#L215-L220}{{\leanfont Space.\allowbreak{}omega\_vacElem\_sum}}}{Projections.lean}{215--220}
\crow{inNormSq}{\href{https://github.com/lsernau/grh-lean/blob/72fead102f663d5c7a407964a13dc0cf9d977207/grh-compact/src/Lib/Weil/Construction/FiniteMultiplier.lean\#L181-L185}{{\leanfont norm\_finMultiplier\_sq}}}{FiniteMultiplier.lean}{181--185}
\crow{inSmulVac}{\href{https://github.com/lsernau/grh-lean/blob/72fead102f663d5c7a407964a13dc0cf9d977207/grh-compact/src/Lib/Weil/Construction/FinitePart.lean\#L57-L61}{{\leanfont Space.\allowbreak{}exists\_eq\_smul\_omega\_vacElem}}}{FinitePart.lean}{57--61}
\crow{inDual}{\href{https://github.com/lsernau/grh-lean/blob/72fead102f663d5c7a407964a13dc0cf9d977207/grh-compact/src/Lib/Weil/Construction/FiniteVacuum.lean\#L62-L64}{{\leanfont stdLattice.\allowbreak{}mem\_stdLattice\_iff\_psi₀}}}{FiniteVacuum.lean}{62--64}
\crow{inProj}{\href{https://github.com/lsernau/grh-lean/blob/72fead102f663d5c7a407964a13dc0cf9d977207/grh-compact/src/Lib/Weil/Construction/Projections.lean\#L149-L152}{{\leanfont Space.\allowbreak{}projOn\_finHeisOp}}}{Projections.lean}{149--152}
\crow{inThArrow}{\href{https://github.com/lsernau/grh-lean/blob/72fead102f663d5c7a407964a13dc0cf9d977207/grh-compact/src/Lib/Weil/DualPair/TwoEdgesOnePair.lean\#L112-L115}{{\leanfont theta\_arrowProduct}}}{TwoEdgesOnePair.lean}{112--115}
\crow{inHomArrow}{\href{https://github.com/lsernau/grh-lean/blob/72fead102f663d5c7a407964a13dc0cf9d977207/grh-compact/src/Lib/Weil/DualPair/TwoEdgesOnePair.lean\#L117-L121}{{\leanfont homHeisOp\_arrowProduct}}}{TwoEdgesOnePair.lean}{117--121}
\crow{inFinArrow}{\href{https://github.com/lsernau/grh-lean/blob/72fead102f663d5c7a407964a13dc0cf9d977207/grh-compact/src/Lib/Weil/DualPair/TwoEdgesOnePair.lean\#L123-L128}{{\leanfont finHeisOp\_arrowProduct}}}{TwoEdgesOnePair.lean}{123--128}
\crow{inArrowPair}{\href{https://github.com/lsernau/grh-lean/blob/72fead102f663d5c7a407964a13dc0cf9d977207/grh-compact/src/Lib/Weil/DualPair/TwoEdgesOnePair.lean\#L133-L135}{{\leanfont arrowProduct\_twistedPairWeil}}}{TwoEdgesOnePair.lean}{133--135}
\crow{inTensorLift}{\href{https://github.com/lsernau/grh-lean/blob/72fead102f663d5c7a407964a13dc0cf9d977207/grh-compact/src/Lib/Weil/DualPair/TwoEdgesOnePair.lean\#L151-L153}{{\leanfont arrowTensor\_map\_liftWeilFin}}}{TwoEdgesOnePair.lean}{151--153}
\crow{inLeftInv}{\href{https://github.com/lsernau/grh-lean/blob/72fead102f663d5c7a407964a13dc0cf9d977207/grh-compact/src/Lib/Weil/DualPair/TwoEdgesOnePair.lean\#L137-L143}{{\leanfont leftInverse\_arrowTensor}}}{TwoEdgesOnePair.lean}{137--143}
\crow{inRange}{\href{https://github.com/lsernau/grh-lean/blob/72fead102f663d5c7a407964a13dc0cf9d977207/grh-compact/src/Lib/Weil/DualPair/TwoEdgesOnePair.lean\#L145-L149}{{\leanfont range\_arrowTensor}}}{TwoEdgesOnePair.lean}{145--149}
\crow{inThRat}{\href{https://github.com/lsernau/grh-lean/blob/72fead102f663d5c7a407964a13dc0cf9d977207/grh-compact/src/Lib/Weil/DualPair/TwoEdgesOnePair.lean\#L155-L158}{{\leanfont theta\_tripleWeil\_rat\_mul}}}{TwoEdgesOnePair.lean}{155--158}
\crow{inPairHeis}{\href{https://github.com/lsernau/grh-lean/blob/72fead102f663d5c7a407964a13dc0cf9d977207/grh-compact/src/Lib/Weil/DualPair/TwoEdgesOnePair.lean\#L32-L34}{{\leanfont pairWeil\_homHeisOp}}}{TwoEdgesOnePair.lean}{32--34}
\crow{inAdapted}{\href{https://github.com/lsernau/grh-lean/blob/72fead102f663d5c7a407964a13dc0cf9d977207/grh-compact/src/Lib/Weil/MixedModel/ZeroTerm.lean\#L272-L288}{{\leanfont Space.\allowbreak{}Splitting.\allowbreak{}exists\_isAdapted\_and\_frame\_eq}}}{ZeroTerm.lean}{272--288}
\crow{inMixedVac}{\href{https://github.com/lsernau/grh-lean/blob/72fead102f663d5c7a407964a13dc0cf9d977207/grh-compact/src/Lib/Weil/MixedModel/ZeroTerm.lean\#L615-L619}{{\leanfont Space.\allowbreak{}Splitting.\allowbreak{}IsAdapted.\allowbreak{}mixedTerm\_omega\_vacElem}}}{ZeroTerm.lean}{615--619}
\crow{inRealZero}{\href{https://github.com/lsernau/grh-lean/blob/72fead102f663d5c7a407964a13dc0cf9d977207/grh-compact/src/Lib/Weil/MixedModel/ZeroTerm.lean\#L395-L405}{{\leanfont Space.\allowbreak{}Splitting.\allowbreak{}IsAdapted.\allowbreak{}exists\_real\_zeroTerm}}}{ZeroTerm.lean}{395--405}
\crow{inZeroRho}{\href{https://github.com/lsernau/grh-lean/blob/72fead102f663d5c7a407964a13dc0cf9d977207/grh-compact/src/Lib/Weil/MixedModel/ZeroTerm.lean\#L567-L568}{{\leanfont Space.\allowbreak{}Splitting.\allowbreak{}zeroTerm\_rho\_iAdelic}}}{ZeroTerm.lean}{567--568}
\crow{inFrameVec}{\href{https://github.com/lsernau/grh-lean/blob/72fead102f663d5c7a407964a13dc0cf9d977207/grh-compact/src/Lib/Weil/MixedModel/ZeroTerm.lean\#L446-L449}{{\leanfont Space.\allowbreak{}Splitting.\allowbreak{}exists\_eq\_finFrameVector}}}{ZeroTerm.lean}{446--449}
\crow{inFrameAdj}{\href{https://github.com/lsernau/grh-lean/blob/72fead102f663d5c7a407964a13dc0cf9d977207/grh-compact/src/Lib/Weil/MixedModel/ZeroTerm.lean\#L441-L443}{{\leanfont Space.\allowbreak{}Splitting.\allowbreak{}adjoint\_mulVec\_finFrameVector}}}{ZeroTerm.lean}{441--443}
\crow{inZeroIncl}{\href{https://github.com/lsernau/grh-lean/blob/72fead102f663d5c7a407964a13dc0cf9d977207/grh-compact/src/Lib/Weil/DualPair/TwoEdgesOnePair.lean\#L58-L68}{{\leanfont zeroTerm\_omega\_of\_inclMatrix}}}{TwoEdgesOnePair.lean}{58--68}
\crow{inHomTarget}{\href{https://github.com/lsernau/grh-lean/blob/72fead102f663d5c7a407964a13dc0cf9d977207/grh-compact/src/Lib/Weil/DualPair/PairSplitting.lean\#L314-L319}{{\leanfont Space.\allowbreak{}Splitting.\allowbreak{}homTarget\_zeroTerm\_pairWeil}}}{PairSplitting.lean}{314--319}
\crow{inNameBox}{\href{https://github.com/lsernau/grh-lean/blob/72fead102f663d5c7a407964a13dc0cf9d977207/grh-compact/src/Lib/Weil/MixedModel/Names.lean\#L91-L92}{{\leanfont Space.\allowbreak{}name\_boxIndicator}}}{Names.lean}{91--92}
\crow{inNameTr}{\href{https://github.com/lsernau/grh-lean/blob/72fead102f663d5c7a407964a13dc0cf9d977207/grh-compact/src/Lib/Weil/MixedModel/Names.lean\#L85-L87}{{\leanfont Space.\allowbreak{}name\_schrodingerAct\_boxIndicator}}}{Names.lean}{85--87}
\crow{inNameOm}{\href{https://github.com/lsernau/grh-lean/blob/72fead102f663d5c7a407964a13dc0cf9d977207/grh-compact/src/Lib/Weil/MixedModel/NameWeilAction.lean\#L194-L202}{{\leanfont Space.\allowbreak{}ω\_name\_on}}}{NameWeilAction.lean}{194--202}
\crow{inBesselMean}{\href{https://github.com/lsernau/grh-lean/blob/72fead102f663d5c7a407964a13dc0cf9d977207/grh-compact/src/Lib/Weil/Periods/BesselIntegral.lean\#L112-L144}{{\leanfont Grading.\allowbreak{}integral\_besselMeasure\_eq\_of\_mean}}}{BesselIntegral.lean}{112--144}
\crow{inGenConst}{\href{https://github.com/lsernau/grh-lean/blob/72fead102f663d5c7a407964a13dc0cf9d977207/grh-compact/src/Lib/Weil/Periods/UnfoldingStep.lean\#L254-L261}{{\leanfont Grading.\allowbreak{}IsGeneric.\allowbreak{}unipotent\_mul\_constTerm}}}{UnfoldingStep.lean}{254--261}
\crow{inConstRat}{\href{https://github.com/lsernau/grh-lean/blob/72fead102f663d5c7a407964a13dc0cf9d977207/grh-compact/src/Lib/Weil/Periods/UnfoldingStep.lean\#L265-L269}{{\leanfont Grading.\allowbreak{}constTermAtT\_rat\_mul}}}{UnfoldingStep.lean}{265--269}
\crow{inAutCoef}{\href{https://github.com/lsernau/grh-lean/blob/72fead102f663d5c7a407964a13dc0cf9d977207/grh-compact/src/Lib/Adelic/CuspForms/AlgebraicForms.lean\#L42-L47}{{\leanfont Space.\allowbreak{}isAutomorphic\_algFormCoeff}}}{AlgebraicForms.lean}{42--47}
\crow{inHermite}{\href{https://github.com/lsernau/grh-lean/blob/72fead102f663d5c7a407964a13dc0cf9d977207/grh-compact/src/Lib/Adelic/Witt/FrameAction.lean\#L74-L80}{{\leanfont IsValuationPair.\allowbreak{}exists\_isPrimitiveMatrix\_mul}}}{FrameAction.lean}{74--80}
\crow{inDense}{\href{https://github.com/lsernau/grh-lean/blob/72fead102f663d5c7a407964a13dc0cf9d977207/grh-compact/src/Lib/Weil/Fock/FockHarmonicCoeff.lean\#L70-L75}{{\leanfont denseRange\_polyVector}}}{FockHarmonicCoeff.lean}{70--75}
\crow{inPolyProd}{\href{https://github.com/lsernau/grh-lean/blob/72fead102f663d5c7a407964a13dc0cf9d977207/grh-compact/src/Lib/General/LinearAlgebra/SmoothMatrices.lean\#L89-L92}{{\leanfont mem\_rename\_range\_mul}}}{SmoothMatrices.lean}{89--92}
\crow{inMeans}{\href{https://github.com/lsernau/grh-lean/blob/72fead102f663d5c7a407964a13dc0cf9d977207/grh-compact/src/Lib/Adelic/QuotientMeasure/PeriodicMean.lean\#L60-L68}{{\leanfont quotMean\_quotMean}}}{PeriodicMean.lean}{60--68}
\crow{inCoprime}{\href{https://github.com/lsernau/grh-lean/blob/72fead102f663d5c7a407964a13dc0cf9d977207/grh-compact/src/Lib/General/Algebra/AdicRingHomExt.lean\#L102-L103}{{\leanfont Nat.\allowbreak{}Coprime.\allowbreak{}norm\_natCast\_eq\_one}}}{AdicRingHomExt.lean}{102--103}}
% the notations of the defined functions: one slot for each explicit argument of the Lean declaration
\newcommand{\AdmComb}[7]{#7\in\operatorname{Adm}(#1,#2;\allowbreak #3,#4;\allowbreak #5,#6)}
\newcommand{\soc}[2]{#1^{\langle #2\rangle}}
\newcommand{\NbV}[4]{\mathsf{Nb}_{#1,#2}(#3;#4)}
\newcommand{\AdmEv}[5]{\mathsf{AdmEv}_{#1,#2,#3}(#4;\allowbreak #5)}
\newcommand{\ztf}[1]{\operatorname{ev}^{f}_{#1}}
\newcommand{\GZP}[3]{\operatorname{GZ}(#1,#2;#3)}
\newcommand{\gzs}[1]{\Xi(#1)}
\newcommand{\stT}[3]{\mathbf{T}(#1,#2,#3)}
\newcommand{\stTf}[3]{\mathbf{T}_f(#1,#2,#3)}
\newcommand{\ArP}[5]{\operatorname{Arr}_{#4,#5}(#1,#2,#3)}
\newcommand{\hzp}[1]{\mathcal{Z}^{\mathrm{Hom}}(#1)}
\newcommand{\aps}[1]{\Xi^{\mathrm{pr}}(#1)}
\newcommand{\fev}[1]{\operatorname{fev}_{#1}}
\newcommand{\glb}[4]{\mathsf{L}_{#1}(#2,#3,#4)}
\newcommand{\Spans}[4]{\mathsf{Lat}(#1;#2,#3,#4)}
\newcommand{\intaw}[1]{O^{(#1)}}
\newcommand{\AdmCo}[6]{#6\in\operatorname{Co}(#1;#2,#3;#4,#5)}
\newcommand{\wtr}[2]{\operatorname{Tr}_{#1}[#2]}
\newcommand{\trl}[3]{\operatorname{Tr}^{#1}_{#2}[#3]}
\newcommand{\SplSq}[3]{\mathsf{SqF}_{#1}(#2;#3)}
\newcommand{\UnitMembers}{312}
\newcommand{\UnitInNote}{57}
\newcommand{\UnitProofs}{83}
\newcommand{\UnitSteps}{503}
\newcommand{\UnitAllSteps}{550}


\linespread{1.08}                       % 11 pt on 14.7 pt
\setlength{\parindent}{0pt}
\setlength{\parskip}{5.5pt plus 1.5pt minus 0.5pt}
\setlength{\abovedisplayskip}{7.5pt plus 2pt minus 3pt}
\setlength{\belowdisplayskip}{7.5pt plus 2pt minus 3pt}
\setlength{\abovedisplayshortskip}{3pt plus 2pt}
\setlength{\belowdisplayshortskip}{6pt plus 2pt minus 2pt}
\emergencystretch=1.5em
\clubpenalty=10000 \widowpenalty=10000 \displaywidowpenalty=10000
\relpenalty=7000 \binoppenalty=9000
\raggedbottom
\pagestyle{fancy}\fancyhf{}\renewcommand{\headrulewidth}{0pt}
\fancyfoot[C]{\sffamily\fontsize{7.6}{9}\selectfont\color{inkfaint}\thepage\,/\,\pageref*{LastPage}}
\setlist{topsep=2pt,partopsep=0pt,parsep=0pt,itemsep=4pt,leftmargin=2.1em,labelsep=0.55em}

\newcommand{\Q}{\mathbb{Q}}
\newcommand{\Z}{\mathbb{Z}}
\newcommand{\F}{\mathbb{F}}
\newcommand{\len}{\operatorname{len}}
\newcommand{\cok}{\operatorname{cok}}
\newcommand{\End}{\operatorname{End}}
\newcommand{\Hom}{\operatorname{Hom}}
\newcommand{\Gal}{\operatorname{Gal}}
\newcommand{\diag}{\operatorname{diag}}
\newcommand{\kw}[1]{\emph{#1}}
\newcommand{\m}[1]{\mbox{$#1$}}        % a short formula that must not be broken across lines
\newcommand{\head}[1]{\par\addvspace{17pt plus 4pt minus 2pt}\Needspace{6\baselineskip}%
  \noindent{\fontsize{13.2}{16}\selectfont\bfseries\color{ink}#1}\par\nopagebreak\vspace{-1pt}}
\newcommand{\sect}[2]{\refstepcounter{section}\label{#2}\par\addvspace{19pt plus 4pt minus 2pt}\Needspace{7\baselineskip}%
  \noindent\llap{{\fontsize{13.2}{16}\selectfont\bfseries\color{inkfaint}\thesection}\hspace{9pt}}%
  {\fontsize{13.2}{16}\selectfont\bfseries\color{ink}#1}\par\nopagebreak\vspace{-1pt}}
\newcommand{\subhead}[1]{\par\addvspace{9pt plus 2pt}\Needspace{4\baselineskip}\noindent\textbf{#1}\enspace\ignorespaces}
% where a statement lives in the Lean archive
\newcommand{\leantag}{{\sffamily\bfseries\fontsize{6.3}{8}\selectfont\color{leancol}\addfontfeatures{LetterSpace=7}LEAN}}
\newcommand{\sep}{\unskip\nobreak\hspace{5pt}{\color{inkfaint}\textperiodcentered}\hspace{5pt plus 2pt}\ignorespaces}
\newcommand{\leanline}[1]{\par\nopagebreak\vspace{-2.5pt}{\fontsize{8.9}{12.4}\selectfont\color{inksoft}\raggedright
  \noindent\llap{\leantag\hspace{9pt}}#1\par}\vspace{1pt}}
% a displayed statement in a tinted box: #1 colour, #2 text
\newcommand{\statement}[2]{\par\addvspace{7pt}\begingroup\setlength{\fboxsep}{8.5pt}\setlength{\fboxrule}{0.6pt}%
  \noindent\fcolorbox{#1!30!white}{#1!3!white}{\parbox{\dimexpr\linewidth-2\fboxsep-2\fboxrule\relax}{#2}}%
  \endgroup\par\addvspace{7pt}}
\newcommand{\nt}[2]{\par\noindent\hangindent=1.1em\hangafter=1 #1\enspace{\color{inksoft}#2}\par\vspace{1pt}}
\newcommand{\legdot}[1]{\tikz[baseline=-0.62ex]\fill[#1] (0,0) circle (2.5pt);\hspace{3pt}}
\newcommand{\crow}[4]{\csname CS@#1\endcsname & #2 & \textcolor{inksoft}{#3} & \textcolor{inksoft}{#4}\\}
\newcommand{\leanquoteline}[2]{\par\noindent\hangindent=\dimexpr#1\fontdimen2\font+2.2em\relax\hangafter=1
  \hspace*{#1\fontdimen2\font}#2}
\newcommand{\quoted}[1]{\par\vspace{3pt}{\ttfamily\fontsize{7.5}{10.4}\selectfont\color{leancol}\raggedright
  \setlength{\parskip}{0pt}\leanquote{#1}\par}\vspace{4pt}}
\newcommand{\bx}[3]{{\sffamily\bfseries\fontsize{7.3}{9.2}\selectfont\color{#1!82!black}#2}\\[1.5pt]#3}
\newcommand{\cs}[2]{\expandafter\newcommand\csname CS@#1\endcsname{#2}}

% theorems
\newtheoremstyle{plainx}{7pt plus 2pt}{7pt plus 2pt}{\itshape}{}{\bfseries}{.}{0.5em}{}
\newtheoremstyle{uprightx}{7pt plus 2pt}{7pt plus 2pt}{}{}{\bfseries}{.}{0.5em}{}
\theoremstyle{plainx}
\newtheorem{theorem}{Theorem}[section]
\newtheorem{lemma}[theorem]{Lemma}
\newtheorem{proposition}[theorem]{Proposition}
\newtheorem{corollary}[theorem]{Corollary}
\theoremstyle{uprightx}
\newtheorem{definition}[theorem]{Definition}
\newtheorem{setting}[theorem]{Setting}          % hypotheses shared by several statements; each of them cites it by number
\newtheorem{remark}[theorem]{Remark}
\newtheorem{example}[theorem]{Example}
\renewcommand{\qedsymbol}{{\setlength{\fboxsep}{0pt}\setlength{\fboxrule}{0.5pt}\fbox{\rule{0pt}{0.46em}\hspace{0.46em}}}}
% what holds throughout the note: the ambient objects only (implicit and instance arguments), never a hypothesis
\newenvironment{conventions}{\par\addvspace{5pt}\noindent\textbf{Conventions.}\enspace\ignorespaces}{\par\addvspace{3pt}}
\makeatletter
\renewenvironment{proof}[1][Proof]{\par\pushQED{\qed}\normalfont\topsep0pt\partopsep0pt
  \trivlist\item[\hskip\labelsep\itshape #1.]\ignorespaces}{\popQED\endtrivlist\@endpefalse}
\makeatother


% where each declaration of the concordance is treated: \cs{key}{place}, the place naming a \label or a lettered item
\cs{integralAway}{Definition~\ref{d:away}}
\cs{IsAdmissibleCoeffs}{Definition~\ref{d:adm}}
\cs{IsAdmissibleComb}{Definition~\ref{d:adm}}
\cs{ofsum}{Lemma~\ref{l:rules}}
\cs{single}{Lemma~\ref{l:rules}}
\cs{map}{Lemma~\ref{l:rules}}
\cs{smul}{Lemma~\ref{l:rules}}
\cs{smulofeqzeroorisOfFinOrder}{Lemma~\ref{l:rules}}
\cs{ofeqisOfFinOrdersmul}{Lemma~\ref{l:rules}}
\cs{mapbilinear}{Lemma~\ref{l:rules}}
\cs{normUnitsmul}{Lemma~\ref{l:rules}}
\cs{bind}{Theorem A}
\cs{smulmem}{Theorem A}
\cs{pure}{Theorem A}
\cs{divmemintegralAway}{Lemma~\ref{l:away}}
\cs{localLevelleintegralAboveofcoprime}{Lemma~\ref{l:away}}
\cs{normfinMultipliereqone}{Lemma~\ref{l:away}}
\cs{normpadicdet}{Lemma~\ref{l:away}}
\cs{weightedTranslations}{Definition~\ref{d:transl}}
\cs{translations}{Definition~\ref{d:transl}}
\cs{selfOrConj}{Definition~\ref{d:transl}}
\cs{mulVecmemweightedTranslations}{Lemma~\ref{l:transl}}
\cs{starmemweightedTranslations}{Lemma~\ref{l:transl}}
\cs{IsSplitSquarefree}{Definition~\ref{d:splsq}}
\cs{gradedLatticeBasis}{Definition~\ref{d:lattice}}
\cs{SpansLattice}{Definition~\ref{d:lattice}}
\cs{isUnimodulargradedLatticeBasisgram}{Proposition~\ref{p:unimod}}
\cs{existsframegradedGram}{Proposition~\ref{p:frame}}
\cs{invxFinmulVecmemofspansLattice}{Lemma~\ref{l:inlattice}}
\cs{isAdmissibleCombweilvacElem}{Proposition~\ref{p:moved}}
\cs{zeroTermFin}{Definition~\ref{d:ztf}}
\cs{zeroTermFinspec}{Proposition~\ref{p:ztfspec}}
\cs{zeroTermFinfinHeisOpfinFrameVector}{Proposition~\ref{p:ztfspec}}
\cs{zeroTermFinsupport}{Proposition~\ref{p:support}}
\cs{zeroTermFinfinHeisOpvacElem}{Proposition~\ref{p:vac}}
\cs{GradeZeroPart}{Definition~\ref{d:gzp}}
\cs{splitting}{Definition~\ref{d:gzp}}
\cs{hom}{Definition~\ref{d:gzphom}}
\cs{ratCoordMatrixsum}{Lemma~\ref{l:coord}}
\cs{inclXmulstringMapmulinclXinv}{Lemma~\ref{l:coord}}
\cs{stackedT}{Definition~\ref{d:stacked}}
\cs{stackedTFin}{Definition~\ref{d:stacked}}
\cs{ArrowPair}{Definition~\ref{d:arrow}}
\cs{homZeroPart}{Definition~\ref{d:fev}}
\cs{splittinga}{Definition~\ref{d:fev}}
\cs{finEval}{Definition~\ref{d:fev}}
\cs{arrowTensorfinEval}{Lemma~\ref{l:atfev}}
\cs{zeroTermsplittingtripleWeil}{Theorem B}
\cs{finEvalinclXFin}{Theorem B}
\cs{middlePair}{Definition~\ref{d:pairs}}
\cs{lastPair}{Definition~\ref{d:pairs}}
\cs{integralconstTermAtTmul}{Theorem C}
\cs{integralthetaeqlift}{Theorem C}
\cs{IsNeighbourVector}{Definition~\ref{d:nbv}}
\cs{IsAdmissibleEval}{Definition~\ref{d:admev}}
\cs{existsisAdmissibleEvalstd}{Proposition~\ref{p:std}}
\cs{isAdmissibleEvalfinEval}{Theorem D}
\cs{inFinFin}{Input~\ref{in:1}}
\cs{inPsiFin}{Input~\ref{in:1}}
\cs{inFinVac}{Input~\ref{in:1}}
\cs{inFinMem}{Input~\ref{in:1}}
\cs{inOmFin}{Input~\ref{in:1}}
\cs{inOmInt}{Input~\ref{in:1}}
\cs{inRhoInt}{Input~\ref{in:1}}
\cs{inVacSum}{Input~\ref{in:1}}
\cs{inNormSq}{Input~\ref{in:1}}
\cs{inSmulVac}{Input~\ref{in:1}}
\cs{inDual}{Input~\ref{in:1}}
\cs{inProj}{Input~\ref{in:1}}
\cs{inThArrow}{Input~\ref{in:2}}
\cs{inHomArrow}{Input~\ref{in:2}}
\cs{inFinArrow}{Input~\ref{in:2}}
\cs{inArrowPair}{Input~\ref{in:2}}
\cs{inTensorLift}{Input~\ref{in:2}}
\cs{inLeftInv}{Input~\ref{in:2}}
\cs{inRange}{Input~\ref{in:2}}
\cs{inThRat}{Input~\ref{in:2}}
\cs{inPairHeis}{Input~\ref{in:2}}
\cs{inAdapted}{Input~\ref{in:3}}
\cs{inMixedVac}{Input~\ref{in:3}}
\cs{inRealZero}{Input~\ref{in:3}}
\cs{inZeroRho}{Input~\ref{in:3}}
\cs{inFrameVec}{Input~\ref{in:3}}
\cs{inFrameAdj}{Input~\ref{in:3}}
\cs{inZeroIncl}{Input~\ref{in:3}}
\cs{inHomTarget}{Input~\ref{in:3}}
\cs{inNameBox}{Input~\ref{in:3}}
\cs{inNameTr}{Input~\ref{in:3}}
\cs{inNameOm}{Input~\ref{in:3}}
\cs{inBesselMean}{Input~\ref{in:4}}
\cs{inGenConst}{Input~\ref{in:4}}
\cs{inConstRat}{Input~\ref{in:4}}
\cs{inAutCoef}{Input~\ref{in:5}}
\cs{inHermite}{Input~\ref{in:6}}
\cs{inDense}{Input~\ref{in:7}}
\cs{inPolyProd}{Input~\ref{in:7}}
\cs{inMeans}{Input~\ref{in:7}}
\cs{inCoprime}{Input~\ref{in:7}}

% ------------------------------------------------------------------ macros of this note that are not notations of
% defined functions (those are made by the build, in data.tex, from the spec)
\newcommand{\CC}{\mathbb{C}}
\newcommand{\NN}{\mathbb{N}}
\newcommand{\RR}{\mathbb{R}}
\newcommand{\A}{\mathbb{A}}
\newcommand{\Af}{\mathbb{A}_f}
\newcommand{\Oh}{\hat{\mathcal{O}}}
\newcommand{\cj}[1]{\overline{#1}}
\newcommand{\cX}{\mathcal{X}}
\newcommand{\cY}{\mathcal{Y}}
\newcommand{\cZ}{\mathcal{Z}}
\newcommand{\cR}{\mathcal{R}}
\newcommand{\pt}{\mathrm{pt}}
\newcommand{\Sch}[1]{\mathcal{S}(#1)}           % the space of data of a hermitian space
\newcommand{\Wf}[1]{\mathcal{W}_f(#1)}          % its finite part
\newcommand{\om}[1]{\omega_{#1}}                % the Weil representation
\newcommand{\rh}[1]{\varrho_{#1}}               % the Heisenberg operators
\newcommand{\rhf}[1]{\varrho^{f}_{#1}}          % the finite Heisenberg operators
\newcommand{\vac}[1]{\mathrm{vac}_{#1}}         % the vacuum
\newcommand{\tht}[1]{\theta_{#1}}               % the theta functional
\newcommand{\Df}{\Delta_f}                      % the finite multiplier
\newcommand{\Hm}[2]{\Hom(#1,#2)}
\newcommand{\fls}[2]{#1^{-}\oplus #2}           % the first space with its form negated, plus the second
\newcommand{\St}[1]{\langle #1\rangle}          % the standard definite space on an index set
\newcommand{\Hyp}[2]{\mathbb{H}^{#1}\oplus #2}  % hyperbolic planes plus a space
\newcommand{\ev}[1]{\operatorname{ev}_{#1}}     % the zero term along a splitting
\newcommand{\gr}[1]{\operatorname{gr}_{#1}}     % the grade of a grading
\newcommand{\Tm}[2]{T_{#1,#2}}                  % the frequency map of two gradings
\newcommand{\inX}[1]{\imath_{#1}}               % the inclusion of the definite part
\newcommand{\inXf}[1]{\imath^{f}_{#1}}
\newcommand{\fc}[1]{\mathfrak{c}(#1)}           % the constant of a grading
\newcommand{\tW}[2]{\omega^{\mathrm{tr}}_{#1,#2}}       % the triple Weil action
\newcommand{\pW}[1]{\omega^{\mathrm{pr}}_{#1}}          % the twisted pair action
\newcommand{\pWf}[1]{\omega^{\mathrm{pr},f}_{#1}}       % its finite part
\newcommand{\at}{\operatorname{at}}
\newcommand{\atl}{\operatorname{at}^{\flat}}
\newcommand{\tsum}[5]{\tau^{#3}_{#4}[#1,#2,#5]}        % the translate sum of the note on the Gauss-sum operator
\newcommand{\Gg}[1]{G(#1)}
\newcommand{\Gfin}[1]{G_f(#1)}
\newcommand{\qG}[1]{[G(#1)]}
\newcommand{\Bes}[1]{B(#1)}
\newcommand{\qB}[1]{[B(#1)]}
\newcommand{\coef}[1]{\operatorname{coef}_{#1}}
\newcommand{\ct}[3]{\operatorname{ct}^{#3}_{#1,#2}}    % the constant term at T of two gradings, for a character
\newcommand{\AF}[3]{\mathcal{A}(#1;#2,#3)}             % algebraic forms on a standard space
\newcommand{\epsm}{\epsilon_{-}}                        % the chosen integral finite adele of norm -1
\newcommand{\wt}{\operatorname{wt}}
\newcommand{\ind}[1]{\operatorname{ind}_{#1}}
\newcommand{\inc}[1]{\operatorname{inc}_{#1}}
\newcommand{\psio}{\psi_0}
\newcommand{\gram}[2]{\operatorname{Gr}_{#1}(#2)}       % the Gram matrix of a definite part plus a string
\newcommand{\gzh}[2]{\Hom(#1,#2)^{0}}                   % the grade-zero part of a space of homomorphisms
\newcommand{\Rmid}[3]{\cR_{\mathrm{mid}}(#1,#2,#3)}
\newcommand{\Rlast}[4]{\cR_{\mathrm{last}}(#1,#2,#3;#4)}
\newcommand{\nun}{\operatorname{nu}}
\newcommand{\Fock}[1]{\mathcal{F}(#1)}                  % the real data (Fock space) of a hermitian space
\newcounter{inputno}
\newcommand{\inp}[1]{\refstepcounter{inputno}\textbf{Input~\theinputno.}}   % the argument is the number, for the reader of the source

\begin{document}
\color{ink}

{\raggedright
{\sffamily\bfseries\fontsize{7.2}{9}\selectfont\color{leancol}\addfontfeatures{LetterSpace=8}READBACK NOTE 22\par}
\vspace{5pt}
{\fontsize{21.5}{25.5}\selectfont\bfseries Two consecutive theta lifts: admissible combinations, arrow pairs and zero terms\par}
\vspace{5pt}
{\fontsize{12.4}{16}\selectfont\itshape\color{inksoft} One unit of the Lean artifact \ArtifactName, read back into
mathematics\par}}
\vspace{8pt}
{\color{rulegrey}\hrule height 0.6pt}
\vspace{7pt}
\thispagestyle{empty}

\textbf{Summary.} The seesaw identity of Chapter~V is proved by unfolding a chain of theta lifts two lifts at a
time. This note supplies that two-lift step and its arithmetic bookkeeping. First, a calculus of \kw{admissible
combinations}: finite linear combinations whose coefficients are integral away from a given integer, up to a
denominator, and bounded at every conjugate. Second, for three hermitian spaces graded by strings of lengths
$l$, \m{l-1}, \m{l-2}, an \kw{arrow pair}: a splitting of the space of homomorphisms that carries two consecutive lifts
at once; its zero term intertwines the Weil actions, so that the period of a generic function against two
consecutive constant terms is a matrix coefficient of two definite theta lifts. Third, the finite part of that
zero term sends translated vacua to admissible combinations of Gauss-sum translates of pairs of translated vacua.

\statement{leancol}{\fontsize{9.6}{13.2}\selectfont\textbf{What this note is.} The artifact \ArtifactName{} is a
body of Lean code whose final theorems state the Riemann hypothesis and its analogue for Dirichlet $L$-functions.
Its authors report that it compiles and that Lean's kernel accepts these theorems from the three standard axioms;
the record of checks of this write-up says who replayed what. That is a fact about a formal text: no
mathematician has read the argument through, and the result is a \emph{claim}. This note treats one part of one
unit of the argument (the folder of the seesaw identity): \UnitInNote{} of its \UnitMembers{} declarations,
those of eight files; the other declarations are treated in the notes on the seesaw identity, on the
Gauss-sum operator and on the characters and index sets of the data. Every statement below was compared, argument
by argument, with its Lean statement, and every proof is written out, except that two statements proved in
the formal text by one automatic step are only described (Lemma~\ref{l:away}, parts~(c) and~(d)). What is
taken for granted is stated as Inputs~1 to~7 with its owner: the construction of the Weil representation, its
mixed model and the unfolding of Bessel periods come from the background library of the artifact, which this
write-up quotes and does not read. The Lean text read is \PinText; Appendix~A is the concordance of \ConcDecls{}
declarations.}

\sect{The statements}{s:intro}
The objects are defined in Section~\ref{s:setting}; the four statements cite its Settings. Throughout, a
\kw{proposition} $p$ is a statement that may or may not hold: it switches the size bounds on.

\statement{arith}{\phantomsection\label{t:A}\textbf{Theorem A} (calculus of admissible combinations).\enspace
Let $\cX$, $\cX'$ be sets, $Z$ a complex vector space, \m{v\colon\cX\to Z} and \m{v'\colon\cX'\to Z} families of
vectors, \m{A\subseteq\cX} and \m{A'\subseteq\cX'}, $n$ and $d$ natural numbers, $p$ a proposition, $B$ a real
number and \m{z\in Z} with $\AdmComb{v}{A}{n}{d}{p}{B}{z}$ (Definition~\ref{d:adm}).

\emph{(i) Substitution.} Let $d'$, $d''$ be natural numbers and $B'$, $C$ real numbers. Suppose that
$\AdmComb{v'}{A'}{n}{d'}{p}{B'}{v(x)}$ for every \m{x\in A}, that \m{d\cdot d'} divides $d''$, and that $p$
implies \m{B\cdot|B'|\le C}. Then $\AdmComb{v'}{A'}{n}{d''}{p}{C}{z}$.

\emph{(ii) Integrality.} Let $\ell$ be a prime number that does not divide~$n$, let
\m{\jmath\colon\CC\to\CC_\ell} be a ring homomorphism, $O_\jmath$ the ring of the complex numbers $w$ with
\m{|\jmath(w)|_\ell\le1}, and $M$ an $O_\jmath$-submodule of $Z$ that contains $v(x)$ for every \m{x\in A}. Then
\m{d\cdot z\in M}.

\emph{(iii) Tensor products.} Let $n$ be a natural number and $p$ a proposition. For \m{j\in\NN} let $\cX_j$ be
an additive group and $A_j$ a subgroup of it; for \m{i\in\NN} let $F_i$ be a complex vector space,
\m{v_i\colon\cX_{2i}\times\cX_{2i+1}\to F_i} a family and \m{f_i\in F_i}. Let $d$, $m$ be natural numbers and $B$
a real number, and suppose $\AdmComb{v_i}{A_{2i}\times A_{2i+1}}{n}{d}{p}{B}{f_i}$ for every \m{i<m}. Let $v_{\otimes m}$ be
the family \m{y\mapsto v_{m-1}(y_{2m-2},y_{2m-1})\otimes\dots\otimes v_0(y_0,y_1)}, indexed by all tuples
\m{y=(y_j)_{j\in\NN}} with \m{y_j\in\cX_j}, and $A_{*}$ the set of the tuples with \m{y_j\in A_j} for all~$j$. Then the
pure tensor satisfies
\[\AdmComb{v_{\otimes m}}{A_{*}}{n}{d^{m}}{p}{|B|^{m}}{f_{m-1}\otimes\dots\otimes f_0}.\]}

\statement{arith}{\phantomsection\label{t:B}\textbf{Theorem B} (the zero term of an arrow pair is equivariant).\enspace
In Setting~\ref{set:three}, let $X''$ and $H$ be finite index sets and \m{\cR\in\ArP{\cY}{\cY'}{\cY''}{X''}{H}} an
arrow pair (Definition~\ref{d:arrow}), with its homomorphism
$\operatorname{emb}_{\cR}\colon\Gg{\St{X_0}}\to\Gg{\St{X''}}$, its splitting $\aps{\cR}$ and its finite evaluation
$\fev{\cR}$ (Definition~\ref{d:fev}). Let $\mu$, $\mu'$ be characters of $\mathrm{U}(1)(\A)$ trivial on
$\mathrm{U}(1)(E)$.

\emph{(i)} For all \m{g\in\Gg{\St{X}}}, \m{g'\in\Gg{\St{X'}}}, \m{k\in\Gg{\St{X_0}}} and
\m{f\in\Sch{\Hm{V'}{\fls{V}{V''}}}},
\[\ev{\aps{\cR}}\Bigl(\tW{\mu}{\mu'}\bigl(\inX{\cY}(g),\inX{\cY'}(g'),\inX{\cY''}(k)\bigr)f\Bigr)
=\tW{\mu}{\mu'}\bigl(g,g',\operatorname{emb}_{\cR}(k)\bigr)\,\ev{\aps{\cR}}(f).\]

\emph{(ii)} Let \m{g\in\Gfin{X}}, \m{g'\in\Gfin{X'}}, \m{k\in\Gfin{X_0}}, \m{k'\in\Gfin{X''}} and \m{x\in\Gfin{I''}}
be finite points with \m{\inXf{\cY''}(k)=x} and \m{\operatorname{emb}_{\cR}(1,k)=(1,k')}, and let
\m{\Phi\in\Wf{\Hm{V}{V'}}}, \m{\Phi'\in\Wf{\Hm{V'}{V''}}}. Then
\[\begin{aligned}
&\fev{\cR}\Bigl(\pWf{\mu}\bigl(\inXf{\cY}(g),\inXf{\cY'}(g')\bigr)\Phi\boxtimes\pWf{\mu'}\bigl(\inXf{\cY'}(g'),x\bigr)\Phi'\Bigr)\\
&\qquad=\bigl(\pWf{\mu}(g,g')\otimes\pWf{\mu'}(g',k')\bigr)\,\fev{\cR}(\Phi\boxtimes\Phi').
\end{aligned}\]}

\statement{arith}{\phantomsection\label{t:C}\textbf{Theorem C} (two unfolding steps are a double lift).\enspace
\emph{(i)} In Setting~\ref{set:three}, let $X''$ and $H$ be finite index sets and
\m{\cR\in\ArP{\cY}{\cY'}{\cY''}{X''}{H}}. There is \m{z\in\CC} such that for all characters $\mu$, $\mu'$ of
$\mathrm{U}(1)(\A)$ trivial on $\mathrm{U}(1)(E)$, all data \m{f\in\Sch{\Hm{V}{V'}}} and
\m{f'\in\Sch{\Hm{V'}{V''}}}, all \m{k\in\Gg{\St{X_0}}}, all functions \m{P\colon\Gg{V}\to\CC} and all
\m{h\in\Gg{V}}: if $P$ is generic for~$\cY$, then
\[\begin{aligned}
&\int_{\qB{\cY'}}\int_{\qB{\cY}}P(gh)\,\ct{\cY}{\cY'}{\mu}(f)(g,r)\,\ct{\cY'}{\cY''}{\mu'}(f')\bigl(r,\inX{\cY''}(k)\bigr)\,dg\,dr\\
&\quad=\int_{\qG{\St{X}}}P\bigl(\inX{\cY}(g)h\bigr)\Bigl(z\int_{\qG{\St{X'}}}\tht{}\bigl(\tW{\mu}{\mu'}\bigl(g,r,\operatorname{emb}_{\cR}(k)\bigr)f_{\cR}\bigr)\,dr\Bigr)dg,\\
&\qquad\text{where }f_{\cR}=\ev{\aps{\cR}}\bigl(\rh{}(\stT{\cY}{\cY'}{\cY''})(f\boxtimes f')\bigr).
\end{aligned}\]
Here \m{\rh{}(\stT{\cY}{\cY'}{\cY''})} is the Heisenberg operator of \m{\Hm{V'}{\fls{V}{V''}}} at the frame coordinates
of the stacked frequency map (Definition~\ref{d:stacked}), and $\tht{}$ is the theta functional of
\m{\Hm{\St{X'}}{\fls{\St{X}}{\St{X''}}}}.

\emph{(ii)} Let $E$ be a CM field with Weil data, $X$, $X'$, $X_0$ finite index sets, $\varkappa$ a weight and $e$,
$e'$, $e''$ determinant exponents, $U$ a set, and for \m{o\in U} let $\mu_o$, $\mu'_o$ be characters of
$\mathrm{U}(1)(\A)$ trivial on $\mathrm{U}(1)(E)$ and
\[\begin{aligned}
\mathcal{L}_o&\colon\Wf{\Hm{\St{X}}{\St{X'}}}\to\Hom_{\CC}\bigl(\AF{X}{\varkappa}{e},\AF{X'}{\varkappa}{e'}\bigr),\\
\mathcal{L}'_o&\colon\Wf{\Hm{\St{X'}}{\St{X_0}}}\to\Hom_{\CC}\bigl(\AF{X'}{\varkappa}{e'},\AF{X_0}{\varkappa}{e''}\bigr)
\end{aligned}\]
linear maps such that $(\mathcal{L}_o)_o$ is a family of theta lifts for $(\mu_o)_o$ and $(\mathcal{L}'_o)_o$ a
family of theta lifts for $(\mu'_o)_o$ (Input~5). Let $u$ be a real datum of
\m{\Hm{\St{X'}}{\fls{\St{X}}{\St{X_0}}}} and $t$ a linear functional on the vectors of weight~$\varkappa$ for~$X$. Let
$\mathcal{L}''_o$ be the linear map on $\Wf{\Hm{\St{X}}{\St{X'}}}\otimes\Wf{\Hm{\St{X'}}{\St{X_0}}}$ with
$\mathcal{L}''_o(\Phi\otimes\Phi')=\mathcal{L}'_o(\Phi')\circ\mathcal{L}_o(\Phi)$. Then there is a linear functional
$t''$ on the vectors of weight~$\varkappa$ for~$X_0$ such that for all \m{o\in U}, all $y$ in the tensor product, all
\m{\varphi\in\AF{X}{\varkappa}{e}} and all \m{g\in\Gg{\St{X_0}}}
\[\begin{aligned}
&\int_{\qG{\St{X}}}\coef{t}(\varphi)(x)\int_{\qG{\St{X'}}}\tht{}\Bigl(\tW{\mu_o}{\mu'_o}(x,r,g)\bigl(u\otimes\at(y)\bigr)\Bigr)\,dr\,dx\\
&\qquad=\coef{t''}\bigl(\mathcal{L}''_o(y)(\varphi)\bigr)(g).
\end{aligned}\]}

\statement{arith}{\phantomsection\label{t:D}\textbf{Theorem D} (the finite evaluation of an arrow pair is admissible).\enspace
In Setting~\ref{set:adm}, with its index sets $X$, $X''$ and its objects $N$, $\Sigma$, $\pi$, $\delta$, $\beta$,
$\kappa$, let $X'$, $X_0$, $H$ be finite index sets and $l$, $l'$, $l''$ natural numbers; let $V$, $V'$, $V''$ be
hermitian spaces over~$E$ with index sets $I$, $I'$, $I''$ and with gradings: $\cY$ of $V$ with definite part
\m{X^{+}} and a string of length~$l$, $\cY'$ of $V'$ with definite part \m{X'^{+}} and a string of length~$l'$,
$\cY''$ of $V''$ with definite part $X_0$ and a string of length~$l''$, all of norm~$N$. Let $a$ be a finite idele
and $u$ a finite adele with \m{u\bar u=1} such that \m{\Spans{\cY'}{a}{\pi}{u}}, and suppose
\m{\delta=(\bar\alpha/\alpha)\cdot(a/\bar a)}, $\alpha$ the idele of the Weil data. Let
\m{\cR\in\ArP{\cY}{\cY'}{\cY''}{X''^{+}}{H}} be an arrow pair, with the matrix $\inc{\cR}$ of its zero part
(rows \m{I\sqcup I''}, columns \m{(H\sqcup H)\sqcup(X^{+}\sqcup X''^{+})}). Suppose:

\emph{(weights)} for all \m{i\in I\sqcup I''} and \m{b\in X^{+}\sqcup X''^{+}}, either \m{\inc{\cR}(i,b)=0} or
\m{|\soc{\beta}{i}|_w=|\soc{\beta}{b}\kappa_b|_w} for every finite place~$w$;

\emph{(lattices)} for all \m{r_1,r_2\in\Af} with \m{r_1\cdot(a/\bar a)=\cj{\epsm}\,\delta} and \m{r_2r_1=1}, all $c$,
$x$ such that \m{\NbV{\pi}{\delta}{c}{x}}, and the vector
$v=\inc{\cR}\cdot\bigl((0,\,r_1uc),\ (r_1\kappa_bx_b)_{b\in X^{+}},\ (\kappa_bx_b)_{b\in X''^{+}}\bigr)$, whose
coordinates on the first copy of $H$ are $0$ and on the second copy all equal to $r_1uc$: all entries of
\m{x_{V,f}^{-1}(r_2\,v|_{I})} and of \m{x_{V'',f}^{-1}(v|_{I''})} lie in~$\Oh$.

Then $\AdmEv{N}{\Sigma}{\beta}{\fev{\cR}}{\soc{\beta}{\,\cdot\,}}$: the finite evaluation of $\cR$ is an admissible
evaluation for $(N,\Sigma,\beta)$ (Definition~\ref{d:admev}, with the index sets $X$, \m{X'^{+}}, $X''$).}

The four theorems are used in the note on the seesaw identity. Theorem~A is the bookkeeping of every
denominator and every size there; Theorem~C, applied once for each pair of consecutive lifts, turns the function left after an even
number of unfoldings into a matrix coefficient of a chain of pairs of definite lifts; Theorem~B makes the
resulting data invariant under the levels; Theorem~D, with Theorem~A(iii), expresses these data through
translated vacua with controlled denominators. The transfer of a Bessel or Fourier--Jacobi period under one
theta lift is known; the file of the seesaw theorem names as related work B.~Wang (2025), Gomez--Zhu (2014) and
Gan--Ichino (2016). The seesaw principle of Kudla, which the background library proves in a weaker form for
periods, is not used in this note.

\sect{The setting}{s:setting}
\begin{conventions}\phantomsection\label{c:conventions}
In every statement $E$ is a CM field with a CM type (one embedding into~$\CC$ chosen from each pair of complex
conjugate embeddings). The bar is complex conjugation on $E$, on its adeles~$\A$, on the finite adeles~$\Af$ and
on the ring $\Oh$ of integral finite adeles; $M^{*}$ is the conjugate transpose of a matrix. Finite places of $E$
are written $w$, with conjugate place $\bar w$; $|\cdot|_w$ is the absolute value at~$w$, with value $e^{-1}$ at a
uniformiser, and $u_w$ is the component at $w$ of a finite adele~$u$. Where a statement says so, $E$ carries
\kw{Weil data} (Definition~\ref{d:weil}). All index sets ($I$, $I'$, $I''$, $I_0$, $I_1$, $I_2$, $X$, $X'$,
$X''$, $X_0$, $H$) are finite; for an index set $X$ we put \m{X^{+}=X\sqcup\{\pt\}}, and \m{\{0,\dots,l-1\}} indexes a
\kw{string} of length~$l$. Vectors over $\Af$ indexed by a product of index sets are also read as matrices.
The letters $Z$, $Z'$, $W$, $F_i$ denote complex vector spaces, $\cX$, $\cX'$ sets that index families of
vectors, $\cX_j$ an additive group. A prime number is $\ell$ or~$N$; \m{\jmath\colon\CC\to\CC_\ell} is a ring
homomorphism and $O_\jmath$ the ring of the \m{w\in\CC} with \m{|\jmath(w)|_\ell\le1}; $\sigma$ runs over the field
automorphisms of~$\CC$. In Section~\ref{s:lattices}, $R$ and $S$ are commutative rings and $R$ has an involution,
written as a bar. A hermitian space $V$ has an index set and signs (Definition~\ref{d:spaces}); gradings are written
$\cY$, $\cY'$, $\cY''$, with strings of lengths $l$, $l'$, $l''$ and norm~$N$.
\end{conventions}

\subhead{Vocabulary of the background library.} The next five definitions recall objects that the artifact
constructs in its background library. They are not declarations of this unit; we state what the proofs use, and
quote their properties as Inputs~1 to~7.

\begin{definition}[hermitian spaces and their groups; recalled]\label{d:spaces}
A \kw{hermitian space} $V$ over $E$ consists of a finite index set $I$, \kw{signs} \m{I\to\{\pm1\}}, a hermitian
matrix \m{J_V\in M_I(E)} and a \kw{frame} \m{x_V\in\mathrm{GL}_I(\A)} such that \m{x_V^{*}J_Vx_V} is the diagonal
matrix of the signs at every infinite place and the identity at every finite place; $x_{V,f}$ is the finite part
of~$x_V$. So $x_{V,f}\Oh^{I}$ is a self-dual lattice for~$J_V$. The adelic unitary group is
\m{\Gg{V}=G_\infty(V)\times\Gfin{I}}, where $G_\infty(V)$ is the product over the infinite places of the unitary
groups of the signs and \m{\Gfin{I}=\{g\in\mathrm{GL}_I(\Af):g^{*}g=1\}}; an element \m{g=(g_\infty,g_f)} has the
\kw{matrix in rational coordinates} \m{x_Vgx_V^{-1}}, which is unitary for~$J_V$. The \kw{rational points}
$\Gamma(V)$ are the $g$ whose matrix in rational coordinates has entries in~$E$, and \m{\qG{V}=\Gamma(V)\backslash\Gg{V}}
carries the invariant measure of the library (a probability measure when $V$ is definite). The \kw{standard space}
$\St{X}$ on an index set $X$ has all signs $+1$, \m{J=1} and frame~$1$; it is definite. For hermitian spaces $V$,
$V'$ (index sets $I$, $I'$), \m{\Hm{V}{V'}} is the hermitian space with index set \m{I'\times I}, matrix
\m{J_{V'}\otimes(J_V^{-1})^{\mathsf T}} and frame \m{x_{V'}\otimes(x_V^{-1})^{\mathsf T}}; a pair
\m{(g,g')\in\Gg{V}\times\Gg{V'}} acts on it by \m{y\mapsto g'yg^{-1}}. \m{V\oplus V'} is the orthogonal sum,
$\mathbb{H}^{H}$ the hyperbolic space with index set \m{H\sqcup H} and matrix
$\bigl(\begin{smallmatrix}0&1\\1&0\end{smallmatrix}\bigr)$, and \m{\Hyp{H}{V}} their sum. An adelic matrix
\m{y\colon V\to V'} \kw{in rational coordinates} is an \m{I'\times I} matrix over~$\A$; its \kw{frame coordinates}
are \m{x_{V'}^{-1}yx_V}, and its \kw{adjoint} is \m{y^{\dagger}=J_V^{-1}y^{*}J_{V'}}.
\end{definition}

\begin{definition}[Weil data, spaces of data, the Weil representation; recalled]\label{d:weil}
\kw{Weil data} for $E$ are: the condition that $E$ is stably unramified in the sense of the library, which
provides a unit $\epsm$ of $\Oh$ with \m{\cj{\epsm}\,\epsm=-1}; an element \m{\mathbf{j}\in E} with
\m{\bar{\mathbf{j}}=-\mathbf{j}}; a finite idele $\alpha$ such that \m{\mathbf{j}\alpha\bar\alpha} generates the
inverse different at every finite place; and a splitting character $\upsilon$, a Hecke character of $E$ whose
restriction to the norm-one ideles is again written~$\upsilon$. For a hermitian space $V$ over such an $E$, with
index set~$I$, the library constructs: a Banach space $\Sch{V}$ of \kw{data}, which are functions on $\A^{I}$;
\kw{Heisenberg operators} $\rh{V}(y)$ on it, for \m{y\in\A^{I}}; the \kw{Weil representation} $\om{V}$ of $\Gg{V}$
on it; the \kw{theta functional} \m{\tht{V}\colon\Sch{V}\to\CC}, invariant under $\Gamma(V)$, with
\m{f(y)=\tht{V}(\rh{V}(y)f)} for every datum $f$ and \m{y\in\A^{I}}; the \kw{vacuum}
\m{\vac{V}\in\Sch{V}}; and a space $\Fock{V}$ of \kw{real data} (a Fock space) with its \kw{polynomial vectors}
$v_Q$, $Q$ a complex polynomial, and a bounded bilinear map \m{\Fock{V}\times\Sch{V}\to\Sch{V}},
\m{(v,f)\mapsto v\otimes f}. The \kw{finite Heisenberg operator} of \m{x\in\Af^{I}} is $\rhf{V}(x)$, a root of unity
times $\rh{V}$ at the adelic vector with real part $0$ and finite part $\alpha x$; the \kw{finite part} of
$\Sch{V}$ is the span $\Wf{V}$ of the vectors \m{\rhf{V}(x)\vac{V}}. For an adelic vector $y$ its \kw{finite
coordinates} are \m{y_f=\alpha^{-1}\cdot(\text{finite part of }y)}. The \kw{finite multiplier} of
\m{g\in\Gfin{I}} is a complex number \m{\Df(g)\ne0}; $\psio$ is the additive character
\m{u\mapsto\psi_f(\mathbf{j}\alpha\bar\alpha u)} of $\Af$, $\psi_f$ the standard one. For a set $S$ of finite
places, $P^{S}$ is the bounded operator on $\Sch{V}$ that averages $\rhf{V}(y)$ over the \m{y\in\Oh^{I}} supported at
the places of $S$ and their conjugates. For matrices \m{y\colon V\to V'} in rational coordinates, \m{\rh{}(y)}
denotes the Heisenberg operator of $\Hm{V}{V'}$ at the frame coordinates of~$y$.
\end{definition}

\statement{ana}{\inp{1}\label{in:1} (the Weil representation; background library, units of the
folders \emph{Weil/Construction} and \emph{Adelic/Adeles}). Let $V$ be a hermitian space with index set $I$
over a field with Weil data.
(a) \m{\rhf{V}(x)\rhf{V}(y)=\varsigma\cdot\rhf{V}(x+y)} with $\varsigma$ a value of $\psio$, and every value of
$\psio$ is a root of unity; \m{\rhf{V}(y)\vac{V}=\vac{V}} for \m{y\in\Oh^{I}}; $\rhf{V}(x)$ maps $\Wf{V}$ into itself.
(b) \m{\om{V}(g)\rhf{V}(x)=\rhf{V}(g_fx)\om{V}(g)} and $\om{V}(g)(v\otimes f)=(\om{V}(g_\infty)v)\otimes\om{V}(1,g_f)f$,
where $G_\infty(V)$ acts on~$\Fock{V}$; for an adelic vector $y$ and \m{v\in\Fock{V}} there is \m{v_1\in\Fock{V}}
with \m{\rh{V}(y)(v\otimes f)=v_1\otimes\rhf{V}(y_f)f} for all~$f$; \m{\rhf{V}(x)(v\otimes f)=v\otimes\rhf{V}(x)f}; and
$\om{V}(1,g)$ maps $\Wf{V}$ into itself.
(c) For \m{g\in\Gfin{I}}, $\om{V}(1,g)\vac{V}=\Df(g)|\Df(g)|^{-2}\sum_y\rhf{V}(gy)\vac{V}$, the sum over
representatives $y$ of $\Oh^{I}$ modulo \m{\Oh^{I}\cap g^{-1}\Oh^{I}}; the number of classes is $|\Df(g)|^{2}$.
(d) If \m{f\in\Wf{V}} and \m{\rhf{V}(gy)f=f} for all \m{y\in\Oh^{I}}, then \m{f=c\cdot\om{V}(1,g)\vac{V}} for some
\m{c\in\CC}.
(e) A vector \m{u\in\Af^{I}} lies in $\Oh^{I}$ if and only if \m{\psio(x^{*}u)=1} for every \m{x\in\Oh^{I}}; and
$\rhf{V}(x)\rhf{V}(y)=\psio(y^{*}x)\cdot\rhf{V}(y)\rhf{V}(x)$.
(f) If \m{P^{S}f=f}, then \m{P^{S}(\rhf{V}(y)f)} is $\rhf{V}(y)f$ if $y$ is integral at the places of $S$ and their
conjugates, and $0$ otherwise.
\emph{References.} The library gives as its sources for the Weil representation A.~Weil, Acta Math.\ 111 (1964),
and S.~Kudla, Israel J. Math.\ 87 (1994); the catalogue of this write-up has not confirmed the places: reference
to be confirmed. The formulas above are those of the lattice model.}
\leanline{\lean{inFinFin}\sep \lean{inPsiFin}\sep \lean{inFinVac}\sep \lean{inFinMem}\sep \lean{inOmFin}\sep
\lean{inOmInt}\sep \lean{inRhoInt}\sep \lean{inVacSum}\sep \lean{inNormSq}\sep \lean{inSmulVac}\sep \lean{inDual}\sep
\lean{inProj}}

\begin{definition}[arrow products and the triple action; recalled]\label{d:arrowprod}
Let $E$ carry Weil data and let $V$, $V'$, $V''$ be hermitian spaces with index sets $I$, $I'$, $I''$. $V^{-}$ is
$V$ with the matrix $-J_V$, the opposite signs and the frame $x_V$ multiplied at the finite places by the idele
\m{\cj{\epsm}\bar\alpha/\alpha}; an element $g$ of $\Gg{V}$, read in $\Gg{V^{-}}$, is written $g^{-}$. For a character
$\mu$ of $\mathrm{U}(1)(\A)$ trivial on $\mathrm{U}(1)(E)$ put \m{\mu_V=\mu\cdot\upsilon^{\lfloor|I|/2\rfloor}}. The
\kw{twisted pair action} on \m{\Sch{\Hm{V}{V'}}} is
$\pW{\mu}(g,g')=\mu_V(\det g')\,\mu_{V'}(\det g)^{-1}\,\om{\Hm{V}{V'}}(g,g')$, and $\pWf{\mu}$ is its restriction
to finite points, acting on \m{\Wf{\Hm{V}{V'}}}. The \kw{arrow product} is a bounded bilinear map
\m{(f,f')\mapsto f\boxtimes f'} from \m{\Sch{\Hm{V}{V'}}\times\Sch{\Hm{V'}{V''}}} to \m{\Sch{\Hm{V'}{\fls{V}{V''}}}}; it
induces the linear map $\at\colon\Wf{\Hm{V}{V'}}\otimes\Wf{\Hm{V'}{V''}}\to\Sch{\Hm{V'}{\fls{V}{V''}}}$, and $\atl$ is
a chosen linear left inverse of~$\at$. The \kw{triple action} of
\m{(g,g',g'')\in\Gg{V}\times\Gg{V'}\times\Gg{V''}} on \m{\Sch{\Hm{V'}{\fls{V}{V''}}}} is
\[\begin{aligned}
\tW{\mu}{\mu'}(g,g',g'')={}&\mu_V(\det g')\,\mu_{V'}(\det g)^{-1}\,\mu'_{V'}(\det g'')\,\mu'_{V''}(\det g')^{-1}\\
&\cdot\upsilon(\det g)^{|I'|}\,\upsilon(\det g')^{-|I|}\cdot\om{}(g',g^{-}\oplus g''),
\end{aligned}\]
$\om{}$ the Weil representation of \m{\Hm{V'}{\fls{V}{V''}}}.
\end{definition}

\statement{ana}{\inp{2}\label{in:2} (two lifts as one pair; background library, unit of the folder
\emph{Weil/DualPair}). In the situation of Definition~\ref{d:arrowprod}:
(a) there is \m{z\ne0} with \m{\tht{}(f\boxtimes f')=z\cdot\tht{}(f)\tht{}(f')} for all $f$, $f'$;
(b) for matrices \m{y\colon V\to V'}, \m{y'\colon V'\to V''} in rational coordinates,
$\rh{}\bigl(\begin{smallmatrix}y^{\dagger}\\y'\end{smallmatrix}\bigr)(f\boxtimes f')=\rh{}(y)f\boxtimes\rh{}(y')f'$;
(c) for a finite vector $y$ of \m{\Hm{V'}{\fls{V}{V''}}}, with blocks \m{y_1} (indices \m{I\times I'}) and
\m{y_2} (indices \m{I''\times I'}), $\rhf{}(y)(f\boxtimes f')=\varsigma\cdot\rhf{}(\epsm\,y_1^{*})f\boxtimes\rhf{}(y_2)f'$
with $\varsigma$ a root of unity, and \m{\vac{}\boxtimes\vac{}=\vac{}};
(d) $\pW{\mu}(g,g')f\boxtimes\pW{\mu'}(g',g'')f'=\tW{\mu}{\mu'}(g,g',g'')(f\boxtimes f')$, and for finite points
$\at\circ(\pWf{\mu}(p,q)\otimes\pWf{\mu'}(q,r))=\tW{\mu}{\mu'}((1,p),(1,q),(1,r))\circ\at$;
(e) $\at$ is injective, \m{\atl\circ\at=1}, and the range of $\at$ is \m{\Wf{\Hm{V'}{\fls{V}{V''}}}};
(f) \m{(g,g')\mapsto\tht{}(\tW{\mu}{\mu'}(g,g',k)f)} is invariant under left multiplication by
\m{\Gamma(V)\times\Gamma(V')};
(g) $\pW{\mu}(g,g')\rh{}(y)=\rh{}(g'yg^{-1})\pW{\mu}(g,g')$, with $g$, $g'$ read in rational coordinates.
\emph{References:} none in the catalogue; reference to be confirmed.}
\leanline{\lean{inThArrow}\sep \lean{inHomArrow}\sep \lean{inFinArrow}\sep \lean{inArrowPair}\sep \lean{inTensorLift}\sep
\lean{inLeftInv}\sep \lean{inRange}\sep \lean{inThRat}\sep \lean{inPairHeis}}

\begin{definition}[splittings and zero terms; recalled]\label{d:split}
Let $V$, $V_0$ be hermitian spaces with index sets $I$, $I_0$ and $H$ a finite index set. A \kw{splitting} $\Xi$ of
$V$ over $V_0$ along~$H$ consists of matrices $\mathsf{k}_\Xi,\mathsf{l}_\Xi\in M_{I\times H}(E)$ and
\m{\mathsf{i}_\Xi\in M_{I\times I_0}(E)} with
\m{\mathsf{k}^{*}J_V\mathsf{k}=0}, \m{\mathsf{l}^{*}J_V\mathsf{l}=0}, \m{\mathsf{l}^{*}J_V\mathsf{k}=1},
\m{\mathsf{i}^{*}J_V\mathsf{i}=J_{V_0}}, \m{\mathsf{i}^{*}J_V\mathsf{k}=0}, \m{\mathsf{i}^{*}J_V\mathsf{l}=0} and
\m{|I|=|I_0|+2|H|}. In frame coordinates these are \m{x_V^{-1}\mathsf{k}_\Xi}, \m{x_V^{-1}\mathsf{l}_\Xi} and
\m{x_V^{-1}\mathsf{i}_\Xi x_{V_0}}; their finite parts are written \m{\mathsf{k}_{\Xi,f}}, \m{\mathsf{l}_{\Xi,f}},
\m{\mathsf{i}_{\Xi,f}}. If $E$ carries Weil data, the \kw{zero term} \m{\ev{\Xi}\colon\Sch{V}\to\Sch{V_0}} is the
bounded linear map with \m{(\ev{\Xi}f)(y)} the mean of $f(x_V^{-1}\mathsf{i}_\Xi x_{V_0}y+x_V^{-1}\mathsf{k}_\Xi b)$
over \m{b\in\A^{H}} modulo~$E^{H}$. Splittings compose: one of $V$ over $V_0$ and one of $V_0$ over $V_1$ give
one of $V$ over $V_1$, and the zero term of the composite is the composite of the zero terms. For a hermitian
space $V'$, a splitting $\Xi$ of $V$ over $V_0$ along $H$ induces one of \m{\Hm{V'}{V}} over \m{\Hm{V'}{V_0}} along
\m{H\times I'}. The space \m{\Hyp{H}{V_0}} has the splitting over $V_0$ that \kw{removes the planes}. An element
\m{g\in\Gg{V}} is \kw{adapted} to $\Xi$ if the matrix with column blocks
$g^{-1}x_V^{-1}(\mathsf{l}_\Xi,\mathsf{k}_\Xi,\mathsf{i}_\Xi x_{V_0})$ has orthonormal columns at every infinite place
and integral entries at the finite places.
\end{definition}

\statement{ana}{\inp{3}\label{in:3} (zero terms and the mixed model; background library, units of
the folders \emph{Weil/MixedModel} and \emph{Weil/DualPair}). Let $\Xi$ be a splitting of $V$ over $V_0$
along~$H$, over a field with Weil data.
(a) An element of $\Gg{V}$ adapted to $\Xi$ exists. If $g$ is adapted to $\Xi$, there are \m{c\ne0} and a
function $\epsilon$ on $\A^{H}$ with
$\ev{\Xi}\bigl(\rh{V}(x_V^{-1}\mathsf{l}_\Xi b)\,\om{V}(g)\vac{V}\bigr)=c\,\epsilon(b)\,\vac{V_0}$ for all \m{b\in\A^{H}},
and there is a bounded linear map \m{\Lambda\colon\Fock{V}\to\Fock{V_0}} with
$\ev{\Xi}(\om{V}(g)(v\otimes f))=\Lambda v\otimes\ev{\Xi}(\om{V}(g)(v_{\mathrm{vac}}\otimes f))$ for all $v$,
$f$, where $v_{\mathrm{vac}}$ is the vacuum of~$\Fock{V}$; moreover \m{v_{\mathrm{vac}}\otimes f=f} for
\m{f\in\Wf{V}}.
(b) $\ev{\Xi}(\rh{V}(x_V^{-1}\mathsf{i}_\Xi x_{V_0}y)f)=\rh{V_0}(y)\ev{\Xi}(f)$ and
$\ev{\Xi}(\rh{V}(x_V^{-1}\mathsf{k}_\Xi b)f)=\ev{\Xi}(f)$.
(c) Every \m{y\in\Af^{I}} is $\mathsf{l}_{\Xi,f}a+\mathsf{i}_{\Xi,f}w+\mathsf{k}_{\Xi,f}b$ for unique $a$, $w$, $b$, and
\m{a=\mathsf{k}_{\Xi,f}^{*}y}, \m{w=\mathsf{i}_{\Xi,f}^{*}y}.
(d) Suppose that $\mathsf{k}_\Xi$ is the inclusion of the basis vectors of positive weight, for weights
\m{\wt\colon I\to\Z}. Let \m{g\in\Gg{V}}, \m{g_0\in\Gg{V_0}}, let $C$ be an \m{I\times I} matrix over $\A$ and $u$ an
idele of norm one, and suppose, in rational coordinates: $g$ and $C$ have the same columns of positive weight;
the entry \m{(b,a)} of \m{C-1} vanishes whenever \m{\wt(b)\le\min(0,\wt(a))}; \m{\det C=u}; and
\m{g\,\mathsf{i}_\Xi-\mathsf{i}_\Xi g_0} has non-zero entries only in rows of positive weight. Then
$\ev{\Xi}(\om{V}(g)f)=\upsilon(u)^{-1}\om{V_0}(g_0)\ev{\Xi}(f)$ for all~$f$.
(e) For the splitting of \m{\Hm{V'}{\Hyp{H}{V_0}}} over \m{\Hm{V'}{V_0}} induced by removing the planes, and
\m{h\in\Gg{V'}}, \m{r\in\Gg{V_0}}:
$\ev{}(\om{}(h,1\oplus r)f)=\upsilon(\det h)^{|H|}\om{}(h,r)\ev{}(f)$.
(f) For a set $S$ of finite places disjoint from its conjugate, functions $\Psi$ of the components at $S$ of a
finite vector \kw{name} operators \m{\Psi\otimes{}} on $\Sch{V}$: the indicator of the vectors integral at $S$ names
$P^{S}$; the name of its translate by $y$ is \m{\rhf{V}(y_S)\circ P^{S}} up to a root of unity, $y_S$ the part of $y$
at $S$ and its conjugate; and for \m{g\in\Gfin{I}} with component $1$ outside $S$ and its conjugate,
$\om{V}(1,g)\circ(\Psi\otimes{})=C_g\cdot(\Psi(g^{-1}\,\cdot\,)\otimes{})$ with a constant \m{C_g\in\CC}.
\emph{References.} The library gives as its sources M\oe glin--Vign\'eras--Waldspurger, LNM 1291 (1987), and
S.~Kudla, Israel J. Math.\ 87 (1994); not confirmed by the catalogue: reference to be confirmed.}
\leanline{\lean{inAdapted}\sep \lean{inMixedVac}\sep \lean{inRealZero}\sep \lean{inZeroRho}\sep \lean{inFrameVec}\sep
\lean{inFrameAdj}\sep \lean{inZeroIncl}\sep \lean{inHomTarget}\sep \lean{inNameBox}\sep \lean{inNameTr}\sep
\lean{inNameOm}}

\begin{definition}[gradings, Bessel groups, constant terms; recalled]\label{d:grading}
A \kw{grading} $\cY$ of a hermitian space $V$ (index set~$I$) with \kw{definite part} $X$ and a string of
length~$l$, of \kw{norm} \m{N\in E^{\times}}, \m{\bar N=N}, consists of a bijection
$\mathbf{e}_{\cY}\colon X\sqcup\{0,\dots,l-1\}\to I$ with \m{J_V=\gram{\mathbf{e}_{\cY}}{N}} and a constant
\m{\fc{\cY}\in E^{\times}}, \m{\cj{\fc{\cY}}=\fc{\cY}}. Here \m{\gram{\mathbf{e}}{N}} is the matrix that is the identity
on~$X$, pairs the string vectors $a$ and \m{l-1-a} to~$1$ for \m{a\ne l-1-a}, and has the entry $-N$ at the middle
string vector \m{a=l-1-a} (which exists when $l$ is odd). The \kw{grade} $\gr{\cY}$ is $0$ on $X$ and \m{l-1-2a}
at the string vector~$a$. The homomorphism \m{\inX{\cY}\colon\Gg{\St{X}}\to\Gg{V}} lets a unitary matrix act on~$X$
and fixes the string vectors (in rational coordinates), and $\inXf{\cY}$ is its restriction to finite points. The
\kw{Bessel group} $\Bes{\cY}$ is the group of the \m{g\in\Gg{V}} that, in rational coordinates, fix every string
vector of grade \m{\ge0} up to string vectors of higher grade; it contains $\inX{\cY}(\Gg{\St{X}})$ and the group
of the \kw{unipotent} elements, those $g$ for which \m{g-1} strictly raises grades. \m{\qB{\cY}} is the quotient of
$\Bes{\cY}$ by its rational points, with the measure of the library. If $E$ carries Weil data, a function
\m{P\colon\Gg{V}\to\CC} is \kw{generic} for $\cY$ if it is continuous, left invariant under the rational points of
$\Bes{\cY}$, and multiplied by an explicit character under left multiplication by the unipotent elements $n$ for
which \m{n-1} raises grades by at least~$2$. For a second grading $\cY'$ of $V'$ with definite part~$X'$, the
\kw{frequency map} \m{\Tm{\cY}{\cY'}\colon V\to V'} is the rational matrix that kills $X$ and sends the string vector
$a$ of $\cY$ to a multiple of the string vector $a$ of~$\cY'$, the multiple being one of $1$, $\fc{\cY}$,
\m{-N\fc{\cY}} according to the grade. The \kw{constant term at the frequency map} of a datum
\m{f\in\Sch{\Hm{V}{V'}}}, for a character $\mu$ as in Definition~\ref{d:arrowprod}, is the function
\m{\ct{\cY}{\cY'}{\mu}(f)} on \m{\Gg{V}\times\Gg{V'}} whose value at $(g,r)$ is the mean, over the adelic matrices
\m{V\to V'} with entries only in positions of positive degree modulo the rational ones, of the datum
\m{\rh{}(\Tm{\cY}{\cY'})\pW{\mu}(g,r)f}; the \kw{degree} of the position \m{(b,a)\in I'\times I} is
\m{\gr{\cY'}(b)-\gr{\cY}(a)}.
\end{definition}

\statement{ana}{\inp{4}\label{in:4} (Bessel periods; background library, unit of the folder
\emph{Weil/Periods}). Let $\cY$, $\cY'$ be gradings of $V$, $V'$ over a field with Weil data.
(a) Let $H$ be a finite index set and $\gamma\colon H\hookrightarrow\{b\in I:\gr{\cY}(b)=1\}$ an injection. There is
\m{C\in\CC} such that \m{\int_{\qB{\cY}}\phi=C\int_{\qG{\St{X}}}M} whenever $\phi$ is a continuous function on $\Gg{V}$,
left invariant under the rational points of $\Bes{\cY}$, such that for every \m{a\in\Gg{\St{X}}} the function
\m{n\mapsto\phi(n\,\inX{\cY}(a))} of unipotent $n$ is a continuous function, periodic for the rational
matrices, of mean $M(a)$, of the entries of \m{\fc{\cY}(n-1)} in the rows $\gamma(H)$ and the columns~$X$ (in rational
coordinates), and $M$ is continuous and left invariant under \m{\Gamma(\St{X})}.
(b) If $P$ is generic for $\cY$ and $n$ is unipotent with \m{n-1} raising grades by at least~$2$, then
$g\mapsto P(ng)\cdot(\text{constant term of }\pW{\mu}(ng,r)f\text{ at a matrix }y)$ equals an explicit value of the
Heisenberg bicharacter of $\Hm{V}{V'}$ times the same expression at~$g$, provided \m{y(n-1)} has only entries in
positions of positive degree.
(c) $\ct{\cY}{\cY'}{\mu}(f)(\gamma g,\gamma'r)=\ct{\cY}{\cY'}{\mu}(f)(g,r)$ for rational points $\gamma$, $\gamma'$ of
the two Bessel groups; and the constant term of \m{\pW{\mu}(ng,n'r)f} at a matrix $y$ is that of \m{\pW{\mu}(g,r)f}
at \m{n'^{-1}yn}, for unipotent $n$ and $n'$ with \m{n'-1} raising degrees.
\emph{References.} The library gives as its source for the unfolding of Bessel periods B.~Wang Peng Jun,
J. Number Theory 271 (2025), Prop.~3.2 and Thm.~4.5, and the authors record the place as compared; the
catalogue of this write-up has not confirmed it: reference to be confirmed.}
\leanline{\lean{inBesselMean}\sep \lean{inGenConst}\sep \lean{inConstRat}}


\statement{ana}{\inp{5}\label{in:5} (algebraic forms and theta lifts; background library, units of
the folders \emph{Adelic/CuspForms} and \emph{Weil/Periods}). For a finite index set~$X$, a weight $\varkappa$ and
a determinant exponent~$e$, \m{\AF{X}{\varkappa}{e}} is the space of algebraic modular forms of weight
$(\varkappa,e)$ on the definite unitary group of~$\St{X}$: functions on \m{\Gfin{X}} with values in the space of
vectors of weight~$\varkappa$ for~$X$ (the catalogue names B.~Gross, Algebraic modular forms, Israel J. Math.\ 113
(1999), for this notion; place not confirmed). For a linear functional $t$ on that space of vectors,
\m{\coef{t}(\varphi)} is the function \m{g\mapsto t(\bar g_\infty^{-1}\varphi(g_f))} on \m{\Gg{\St{X}}}; it is continuous
and left invariant under \m{\Gamma(\St{X})}. Let $E$ carry Weil data. A family of linear maps
$\mathcal{L}_o\colon\Wf{\Hm{\St{X}}{\St{X'}}}\to\Hom_{\CC}(\AF{X}{\varkappa}{e},\AF{X'}{\varkappa}{e'})$, \m{o\in U}, is a
\kw{family of theta lifts} for characters $(\mu_o)_{o\in U}$ if for every complex polynomial $Q$ and every functional
$t$ there is a functional $t'$ such that for all \m{o\in U}, $\Phi$, $\varphi$ and \m{h\in\Gg{\St{X'}}}
\[\int_{\qG{\St{X}}}\tht{}\bigl(\pW{\mu_o}(g,h)(v_Q\otimes\Phi)\bigr)\,\coef{t}(\varphi)(g)\,dg=\coef{t'}\bigl(\mathcal{L}_o(\Phi)\varphi\bigr)(h).\]
This is a definition of the library; the lifts to which the seesaw note applies it are theta lifts in this sense
by a result of the library that this note does not use.}
\leanline{\lean{inAutCoef}}

\statement{ana}{\inp{6}\label{in:6} (Hermite reduction; background library, unit of the folder
\emph{Adelic/Witt}). Let $I$, $H$ be finite index sets and $Y$ an \m{I\times H} matrix over $\Af$ that has a left
inverse. Then there is an invertible \m{H\times H} matrix $t$ over $\Af$ such that $Yt$ has entries in $\Oh$ and a
left inverse with entries in~$\Oh$. (The library proves this for pairs of a ring and a subring with a valuation
property, with $t$ block triangular; the catalogue classes it as elementary.)}
\leanline{\lean{inHermite}}

\statement{ana}{\inp{7}\label{in:7} (analysis; background library). (a) The polynomial vectors $v_Q$
are dense in the space $\Fock{V}$ of real data (folder \emph{Weil/Fock}). (b) If two families of variables together
exhaust the variables of a polynomial ring over~$\CC$, every polynomial is a finite sum of products of a
polynomial in the first family with a polynomial in the second (folder \emph{General/LinearAlgebra}). (c) Let
$\phi\colon\mathsf{A}\times\mathsf{A}'\to\mathsf{B}$ be a continuous surjective homomorphism of topological abelian
groups that maps the product of two cocompact lattices into a cocompact lattice of~$\mathsf{B}$. Then the mean
over $\mathsf{A}$ of the mean over $\mathsf{A}'$ of \m{F\circ\phi} is the mean of $F$ over~$\mathsf{B}$, each mean
taken modulo its lattice (folder \emph{Adelic/QuotientMeasure}). (d) If a prime $\ell$ does not divide the natural
number~$n$, then \m{|n|=1} in every normed field that is a normed algebra over the $\ell$-adic numbers (folder
\emph{General/Algebra}).}
\leanline{\lean{inDense}\sep \lean{inPolyProd}\sep \lean{inMeans}\sep \lean{inCoprime}}

\begin{setting}[three graded spaces]\label{set:three}
$E$ is a CM field with a CM type and Weil data. $V$, $V'$, $V''$ are hermitian spaces over $E$ with finite index
sets $I$, $I'$, $I''$ and their signs. $\cY$, $\cY'$, $\cY''$ are gradings of $V$, $V'$, $V''$ of the same norm
\m{N\in E}, with finite definite parts $X$, $X'$, $X_0$ and strings of lengths $l$, $l'$, $l''$.
\end{setting}

\subhead{Admissible combinations.}
\begin{definition}[integral away from $n$]\label{d:away}
For a natural number $n$, \m{\intaw{n}} is the set of the complex numbers $w$ with \m{|\jmath(w)|_\ell\le1} for every
prime number $\ell$ coprime to~$n$ and every ring homomorphism \m{\jmath\colon\CC\to\CC_\ell}. It is a subring
of~$\CC$, by the ultrametric inequality, and \m{\intaw{n}\subseteq O_\jmath} for every such $\ell$ and~$\jmath$.
\end{definition}
\leanline{\lean{integralAway}}

\begin{definition}[admissible coefficients and combinations]\label{d:adm}
Let $\cX$ be a set, \m{A\subseteq\cX}, $n$ and $d$ natural numbers, $p$ a proposition and $B$ a real number. A
finitely supported function \m{r\colon\cX\to\CC} is \kw{admissible}, written \m{\AdmCo{A}{n}{d}{p}{B}{r}}, if its support
lies in~$A$; \m{d\cdot r(x)\in\intaw{n}} for every \m{x\in\cX}; and, if $p$ holds, \m{\sum_x|\sigma(r(x))|\le B} for
every field automorphism $\sigma$ of~$\CC$. For a complex vector space $Z$, a family \m{v\colon\cX\to Z} and
\m{z\in Z} we write $\AdmComb{v}{A}{n}{d}{p}{B}{z}$ if \m{z=\sum_xr(x)v(x)} for some $r$ with
\m{\AdmCo{A}{n}{d}{p}{B}{r}}: $z$ is an \kw{admissible combination} of the $v(x)$, \m{x\in A}, \kw{away from} $n$,
with \kw{denominator}~$d$ and, under~$p$, \kw{size}~$B$. The size is an upper bound, and it is asked only when
$p$ holds.
\end{definition}
\leanline{\lean{IsAdmissibleCoeffs}\sep \lean{IsAdmissibleComb}}

\subhead{Translations and lattices.}
\begin{definition}[groups of translations]\label{d:transl}
Let $E$ be a CM field, $I$ an index set, $n$ a natural number and \m{d\colon I\to\Af} a family of finite adeles.
\m{\wtr{n}{d}} is the additive group of the \m{x\in\Af^{I}} such that \m{d_ix_i} is integral at every finite place
that does not divide~$n$, for every \m{i\in I}. For a single finite adele $d$, \m{\trl{I}{n}{d}} is the group
\m{\wtr{n}{d}} for the constant family \m{i\mapsto d} on~$I$. For index sets $I_1$, $I_2$, \m{\beta\in\Oh} and
\m{j\in I_1\sqcup I_2} we put \m{\soc{\beta}{j}=\beta} if \m{j\in I_1} and \m{\soc{\beta}{j}=\bar\beta} if \m{j\in I_2}.
\end{definition}
\leanline{\lean{weightedTranslations}\sep \lean{translations}\sep \lean{selfOrConj}}

\begin{definition}[split and squarefree]\label{d:splsq}
Let $E$ be a CM field with a CM type, $\Sigma$ a finite set of finite places of~$E$, \m{N\in E} and $\pi$ a finite
idele. We write \m{\SplSq{\Sigma}{N}{\pi}}, and say that $N$ is \kw{split and squarefree along} $\Sigma$ through~$\pi$, if
\m{N=\pi\bar\pi} in $\Af$ and, for every finite place~$w$,
\m{|\pi_w|_w=e^{-1}} if \m{w\in\Sigma} and \m{|\pi_w|_w=1} otherwise, and \m{w\in\Sigma} implies
\m{\bar w\notin\Sigma}.
\end{definition}
\leanline{\lean{IsSplitSquarefree}}

\begin{definition}[graded lattice basis]\label{d:lattice}
Let $R$ be a commutative ring with an involution, $X$ and $I$ index sets, $l$ a natural number,
\m{\mathbf{e}\colon X^{+}\sqcup\{0,\dots,l-1\}\to I} a bijection, $a$ and $\pi$ units of~$R$ and \m{u\in R} with
\m{u\bar u=1}. The matrix \m{\glb{\mathbf{e}}{a}{\pi}{u}\in M_I(R)} is defined by its columns, indexed
through~$\mathbf{e}$ (write $e_i$ for the basis vectors). If $l$ is odd, with middle string index~$z$: the column of
\m{x\in X} is~$e_x$, the column of $\pt$ is \m{\bar\pi^{-1}e_{\pt}+u(\pi\bar\pi)^{-1}e_z}, and the column of a
string index $b$ is~$ue_b$. If $l$ is even: the column of \m{x\in X^{+}} is $e_x$, and the column of a string
index $b$ is \m{ua^{-1}e_b} if \m{2b<l} and \m{u\bar ae_b} otherwise.

Let now $E$ be a CM field with a CM type, $V$ a hermitian space with index set~$I$, and $\cY$ a grading of $V$
with definite part $X^{+}$ and a string of length~$l$, of norm~$N$; let $a$, $\pi$ be finite ideles and $u$ a finite
adele with \m{u\bar u=1}. We write \m{\Spans{\cY}{a}{\pi}{u}} if
\m{x_{V,f}=\glb{\mathbf{e}_{\cY}}{a}{\pi}{u}\cdot k} for some \m{k\in\mathrm{GL}_I(\Oh)}: the columns of the finite
frame of $V$ and of the graded lattice basis span the same lattice.
\end{definition}
\leanline{\lean{gradedLatticeBasis}\sep \lean{SpansLattice}}

\subhead{Zero terms and arrow pairs.}
\begin{definition}[the finite zero term]\label{d:ztf}
Let $E$ carry Weil data and let $\Xi$ be a splitting of $V$ over $V_0$ along~$H$ (Definition~\ref{d:split}; index
sets $I$, $I_0$; signs as in Definition~\ref{d:spaces}). Choose an element of $\Gg{V}$ adapted to~$\Xi$
(Input~3(a)), let $g_\infty$ be its real component and $v_\Xi=\om{V}(g_\infty)v_{\mathrm{vac}}\in\Fock{V}$. The
\kw{finite zero term} is the bounded linear map
\[\ztf{\Xi}\colon\Sch{V}\to\Sch{V_0},\qquad \ztf{\Xi}(f)=\ev{\Xi}(v_\Xi\otimes f).\]
\end{definition}
\leanline{\lean{zeroTermFin}}

\begin{definition}[grade-zero parts]\label{d:gzp}
Let $E$ carry Weil data, let $V$, $V''$ be hermitian spaces with index sets $I$, $I''$ and their signs, and
\m{\wt\colon I\to\Z} a function (\kw{weights}). A \kw{grade-zero part} of $V$ for $\wt$ with zero part $V''$, an
element of the set \m{\GZP{V}{V''}{\wt}}, is a pair \m{\cZ=(\ind{\cZ},\inc{\cZ})} of an injection
\m{\ind{\cZ}\colon I''\hookrightarrow I} whose image is \m{\{i:\wt(i)=0\}} and a matrix
\m{\inc{\cZ}\in M_{I\times I''}(E)} such that \m{\inc{\cZ}(a,b)=0} whenever \m{\wt(a)\ne0}, such that
\m{\inc{\cZ}^{*}J_V\inc{\cZ}=J_{V''}}, and such that \m{J_V(a,b)\ne0} implies \m{\wt(a)+\wt(b)=0}. Its
\kw{splitting} \m{\gzs{\cZ}} is the splitting of $V$ over $V''$ along \m{\{i:\wt(i)>0\}} with $\mathsf{k}$ the
inclusion of the basis vectors of positive weight, \m{\mathsf{l}=J_V^{-1}\mathsf{k}} and
\m{\mathsf{i}=\inc{\cZ}}; the equations of Definition~\ref{d:split} hold because $J_V$ and $J_V^{-1}$ pair only
opposite weights.
\end{definition}
\leanline{\lean{GradeZeroPart}\sep \lean{splitting}}

\begin{definition}[the grade-zero part of a space of homomorphisms]\label{d:gzphom}
Let $E$ carry Weil data, let $V$, $V'$, $V''$ be hermitian spaces with index sets $I$, $I'$, $I''$, let $\cY'$ be a
grading of $V'$ of norm $N$ with definite part $X'$ and a string of length~$l'$, let \m{\wt\colon I\to\Z} and
\m{\cZ\in\GZP{V}{V''}{\wt}}. Suppose that $l'$ is even and that $\wt(i)$ is even for every \m{i\in I}. Then
\m{\gzh{\cY'}{\cZ}} is the grade-zero part of \m{\Hm{V'}{V}} for the weights \m{(i,j)\mapsto\wt(i)-\gr{\cY'}(j)}, with
zero part \m{\Hm{\St{X'}}{V''}}: its injection is $\ind{\cZ}$ times the inclusion of $X'$ into $I'$
along~$\mathbf{e}_{\cY'}$, and its matrix is the Kronecker product of $\inc{\cZ}$ with the inclusion matrix of~$X'$.
(The two hypotheses make \m{\wt(i)-\gr{\cY'}(j)=0} possible only for \m{j\in X'}.)
\end{definition}
\leanline{\lean{hom}}

\begin{definition}[the stacked frequency map]\label{d:stacked}
In Setting~\ref{set:three}, \m{\stT{\cY}{\cY'}{\cY''}} is the adelic matrix \m{V'\to\fls{V}{V''}} in rational
coordinates with the two row blocks \m{\Tm{\cY}{\cY'}^{\dagger}} (rows~$I$) and \m{\Tm{\cY'}{\cY''}} (rows~$I''$), and
\m{\stTf{\cY}{\cY'}{\cY''}} is the vector of the finite coordinates of its frame coordinates, a finite vector of
\m{\Hm{V'}{\fls{V}{V''}}}.
\end{definition}
\leanline{\lean{stackedT}\sep \lean{stackedTFin}}

\begin{definition}[arrow pairs]\label{d:arrow}
In Setting~\ref{set:three}, let $X''$ and $H$ be finite index sets. An \kw{arrow pair} for
\m{\cY\to\cY'\to\cY''}, an element of the set \m{\ArP{\cY}{\cY'}{\cY''}{X''}{H}}, consists of the following.
\begin{enumerate}[label=(\alph*)]
\item A grade-zero part $\cZ_{\cR}\in\GZP{\fls{V}{V''}}{\Hyp{H}{(\fls{\St{X}}{\St{X''}})}}{\gr{\cY}\sqcup\gr{\cY''}}$;
  its matrix is written $\inc{\cR}$ and its injection $\ind{\cR}$.
\item The conditions: $l$ is odd, \m{l'=l-1}, \m{l''=l'-1} (subtraction of natural numbers, truncated at~$0$: so
  \m{l''=l-2} if \m{l\ge3} and \m{l'=l''=0} if \m{l=1}), \m{\fc{\cY'}=\fc{\cY}}, and \m{|I|+|I''|} is even.
\item A homomorphism $\operatorname{emb}_{\cR}\colon\Gg{\St{X_0}}\to\Gg{\St{X''}}$ with
  \m{\det\operatorname{emb}_{\cR}(k)=\det k}, such that for all \m{g\in\Gg{\St{X}}} and \m{k\in\Gg{\St{X_0}}}, in
  rational coordinates,
  \[\bigl(\inX{\cY}(g)\oplus\inX{\cY''}(k)\bigr)\cdot\inc{\cR}=\inc{\cR}\cdot\bigl(1\oplus g^{-}\oplus\operatorname{emb}_{\cR}(k)\bigr).\]
\item An injection $\gamma_{\cR}\colon H\hookrightarrow\{b\in I':\gr{\cY'}(b)=1\}$ such that for every index $a$ of
  the zero part and every \m{b\in I'} the entry of \m{\stT{\cY}{\cY'}{\cY''}} at \m{(\ind{\cR}(a),b)} is
  $\sum_{k\in H,\ \gamma_{\cR}(k)=b}\inc{\cR}(\ind{\cR}(a),k^{*})\cdot\fc{\cY'}$, where $k^{*}$ is the index of the
  $k$-th vector of the second copy of $H$ in~$\mathbb{H}^{H}$.
\end{enumerate}
\end{definition}
\leanline{\lean{ArrowPair}}

\begin{definition}[the splitting and the finite evaluation of an arrow pair]\label{d:fev}
In Setting~\ref{set:three}, let $X''$, $H$ be finite index sets and \m{\cR\in\ArP{\cY}{\cY'}{\cY''}{X''}{H}}.
(i) \m{\hzp{\cR}=\gzh{\cY'}{\cZ_{\cR}}} is the grade-zero part of $\Hm{V'}{\fls{V}{V''}}$ with zero part
$\Hm{\St{X'}}{\Hyp{H}{(\fls{\St{X}}{\St{X''}})}}$; Definition~\ref{d:gzphom} applies because $l'$ is even and the
grades of $\cY$ and $\cY''$ are even, $l$ being odd and $l''$ odd or zero.
(ii) \m{\aps{\cR}} is the splitting of $\Hm{V'}{\fls{V}{V''}}$ over $\Hm{\St{X'}}{\fls{\St{X}}{\St{X''}}}$ obtained by
composing the splitting \m{\gzs{\hzp{\cR}}} with the splitting induced by removing the $H$ planes.
(iii) The \kw{finite evaluation} of $\cR$ is the linear map
\[\fev{\cR}=\atl\circ\ztf{\aps{\cR}}\circ\rhf{}\bigl(\stTf{\cY}{\cY'}{\cY''}\bigr)\]
from $\Sch{\Hm{V'}{\fls{V}{V''}}}$ to $\Wf{\Hm{\St{X}}{\St{X'}}}\otimes\Wf{\Hm{\St{X'}}{\St{X''}}}$.
\end{definition}
\leanline{\lean{homZeroPart}\sep \lean{splittinga}\sep \lean{finEval}}

\begin{definition}[the middle pair and the last pair]\label{d:pairs}
In Setting~\ref{set:three}, suppose that $l$ is odd, \m{l'=l-1}, \m{l''=l'-1} and \m{\fc{\cY'}=\fc{\cY}}, and that
\m{|I|+|I''|} is even. Let $z$, $z''$ be the middle string vectors of $\cY$ and~$\cY''$.

(i) If \m{l\ne1}, the \kw{middle pair} \m{\Rmid{\cY}{\cY'}{\cY''}}, an element of \m{\ArP{\cY}{\cY'}{\cY''}{X_0}{\{\pt\}}}, removes one
hyperbolic plane: its zero part sends the two basis vectors of the plane to \m{(2N)^{-1}(z-z'')} and \m{z+z''}, and
the basis vectors of $X$ and of $X_0$ to the corresponding basis vectors of $V$ and~$V''$; its homomorphism
\m{\operatorname{emb}} is the identity of \m{\Gg{\St{X_0}}}; its injection $\gamma$ picks the string vector of
$\cY'$ with index~\m{(l''-1)/2}, which has grade~$1$.

(ii) If \m{l=1}, and \m{\pi_0\in E} satisfies \m{\pi_0\cj{\pi_0}=N}, the \kw{last pair}
\m{\Rlast{\cY}{\cY'}{\cY''}{\pi_0}}, an element of \m{\ArP{\cY}{\cY'}{\cY''}{X_0^{+}}{\emptyset}}, removes no plane: its zero part sends
the basis vectors of $X$ and of $X_0$ to the corresponding basis vectors and the new point to \m{z/\pi_0}, $z$ the
only string vector of~$\cY$; its homomorphism \m{\operatorname{emb}} is the block embedding \m{k\mapsto k\oplus1}
of \m{\Gg{\St{X_0}}} into \m{\Gg{\St{X_0^{+}}}}.

In both cases the conditions of Definition~\ref{d:arrow} are verified entry by entry from the explicit matrices
\m{\gram{\mathbf{e}}{N}} and the definition of the frequency maps.
\end{definition}
\leanline{\lean{middlePair}\sep \lean{lastPair}}

\subhead{Admissible evaluations.}
\begin{setting}[data at a split prime]\label{set:adm}
$E$ is a CM field with a CM type and Weil data, and $X$, $X''$ are finite index sets. $N$ is a prime number and
$\Sigma$ a finite set of finite places of $E$ that is a \kw{split level} for~$N$ in the sense of the note on the
Gauss-sum operator: \m{\bar w\notin\Sigma} for \m{w\in\Sigma}, every \m{w\in\Sigma} has norm~$N$ and $N$ is a
uniformiser at~$w$, and a finite place divides $N$ exactly if it or its conjugate lies in~$\Sigma$. $\pi$ is a
finite idele with \m{\SplSq{\Sigma}{N}{\pi}}. $\delta$ is a finite adele with \m{\delta\bar\delta=1}, and
\m{\beta\in\Oh}, such that \m{\delta_w=1} and \m{|\beta_w|_w=1} at every finite place $w$ dividing~$N$. Finally
\m{\kappa\colon X^{+}\sqcup X''^{+}\to\Af} is a family with \m{\kappa_b\in\{1,\delta\}}, \m{\bar\beta\kappa_b\in\Oh} and
\m{\soc{\beta}{b}\kappa_b\in\Oh} for every~$b$.
\end{setting}

\begin{definition}[neighbour vectors]\label{d:nbv}
Let $E$ carry Weil data, let $X$, $H$ be finite index sets, $\pi$ a finite idele, $\delta$ and $c$ finite adeles
and $x$ a vector over $\Af$ indexed by \m{X^{+}\sqcup H^{+}}; write $\pt_1$, $\pt_2$ for the points of $X^{+}$
and~$H^{+}$. We write \m{\NbV{\pi}{\delta}{c}{x}}, and call $(c,x)$ a \kw{neighbour vector} for $(\pi,\delta)$, if: \m{\delta c=c}; \m{x_i\in\Oh} for every index $i$ other than
$\pt_1$, $\pt_2$; \m{\bar\pi x_i\in\Oh} for every~$i$; \m{c-\pi^{-1}x_{\pt_1}\in\Oh}; and
\m{\cj{\epsm}c-\pi^{-1}x_{\pt_2}\in\Oh}.
\end{definition}
\leanline{\lean{IsNeighbourVector}}

\begin{definition}[admissible evaluations]\label{d:admev}
Let $E$ carry Weil data; let $V$, $V'$, $V''$ be hermitian spaces with index sets $I$, $I'$, $I''$ and their signs,
and $X$, $X'$, $X''$ finite index sets. Let $N$ be a prime number, $\Sigma$ a finite set of finite places,
\m{\beta\in\Oh}, \m{d\colon I\sqcup I''\to\Af} a family (by default \m{d=\soc{\beta}{\,\cdot\,}}), and $L$ a map
from \m{\Sch{\Hm{V'}{\fls{V}{V''}}}} to $\Wf{\Hm{\St{X^{+}}}{\St{X'}}}\otimes\Wf{\Hm{\St{X'}}{\St{X''^{+}}}}$. We write
$\AdmEv{N}{\Sigma}{\beta}{L}{d}$, and call $L$ an \kw{admissible evaluation} for $(N,\Sigma,\beta)$ with weights~$d$,
if there is \m{C\in\CC} such that for every natural number $n$ coprime
to~$N$, every proposition~$p$ and every \m{y\in\wtr{n}{(i,i')\mapsto d_i}} (a finite vector indexed by
\m{(I\sqcup I'')\times I'})
\[\AdmComb{v_C}{A_n}{n}{D}{p}{1}{L\bigl(\rhf{}(y)\vac{}\bigr)},\]
where $v_C$ is the family
\m{t=(t_1,t_2)\mapsto C\cdot\tsum{X}{X'}{\Sigma}{N}{X''}\bigl(\rhf{}(t_1)\vac{}\otimes\rhf{}(t_2)\vac{}\bigr)}, the set $A_n$ is
\m{\trl{X'\times X^{+}}{n}{\bar\beta}\times\trl{X''^{+}\times X'}{n}{\bar\beta}}, $D$ is the number of maps from \m{(X^{+}\sqcup X''^{+})\times X'} to \m{\Oh/\beta\Oh} (zero if infinite), and
\m{\tsum{X}{X'}{\Sigma}{N}{X''}} is the translate sum of the note on the Gauss-sum operator: the operator
$\sum_y\rhf{}(\mathrm{c}(y))\otimes\rhf{}(\mathrm{r}(y))$ on
$\Wf{\Hm{\St{X^{+}}}{\St{X'}}}\otimes\Wf{\Hm{\St{X'}}{\St{X''^{+}}}}$, the sum over \m{y\in((\Z/N)^{\Sigma})^{X'}},
where the matrix $\mathrm{c}(y)$ has last column \m{-\tilde y/N} and $\mathrm{r}(y)$ has last row
\m{-\cj{\tilde y/N}}, both zero elsewhere, and \m{\tilde y/N} is the finite adele with component \m{\tilde y_w/N} at
\m{w\in\Sigma}, $\tilde y_w$ the least non-negative representative, and $0$ at the other places.
\end{definition}
\leanline{\lean{IsAdmissibleEval}}

\sect{Admissible combinations}{s:adm}
\begin{lemma}[rules]\label{l:rules}
Let $\cX$, $\cX'$ be sets, $Z$, $Z'$, $W$ complex vector spaces, \m{v\colon\cX\to Z} a family, \m{A\subseteq\cX},
\m{A'\subseteq\cX'}, $n$, $d$, $d'$ natural numbers, $p$ a proposition and $B$, $B'$ real numbers.
\begin{enumerate}[label=(\alph*)]
\item Let \m{F\subseteq\cX} be a finite subset and \m{c\colon\cX\to\CC} a function such that
  \m{d\cdot c(x)\in\intaw{n}} for every \m{x\in F}, and such that $p$ implies \m{\sum_{x\in F}|\sigma(c(x))|\le B} for
  every automorphism~$\sigma$. Then $\AdmComb{v}{F}{n}{d}{p}{B}{\sum_{x\in F}c(x)v(x)}$.
\item If \m{x\in A}, then $\AdmComb{v}{A}{n}{d}{p}{1}{v(x)}$.
\end{enumerate}
In (c) to (g) let \m{z\in Z} with $\AdmComb{v}{A}{n}{d}{p}{B}{z}$.
\begin{enumerate}[label=(\alph*),resume]
\item For a linear map \m{L\colon Z\to Z'}, $\AdmComb{L\circ v}{A}{n}{d}{p}{B}{L(z)}$.
\item If \m{c\in\intaw{n}}, and $p$ implies \m{|\sigma(c)|\le1} for every automorphism~$\sigma$, then
  $\AdmComb{v}{A}{n}{d}{p}{B}{c\cdot z}$.
\item If \m{\varsigma\in\CC} is zero or a root of unity, then $\AdmComb{v}{A}{n}{d}{p}{B}{\varsigma\cdot z}$.
\item Let \m{v'\colon\cX'\to Z}. If for every \m{x\in A} there are a root of unity $\varsigma$ and \m{y\in A'} with
  \m{v(x)=\varsigma\cdot v'(y)}, then $\AdmComb{v'}{A'}{n}{d}{p}{B}{z}$.
\item Let \m{v'\colon\cX'\to Z'} be a family, let \m{z'\in Z'} satisfy
  $\AdmComb{v'}{A'}{n}{d'}{p}{B'}{z'}$, and let \m{b\colon Z\times Z'\to W} be bilinear. Then, for the family \m{v_b\colon(x,x')\mapsto b(v(x),v'(x'))} on
  \m{\cX\times\cX'}, $\AdmComb{v_b}{A\times A'}{n}{d\cdot d'}{p}{B\cdot|B'|}{b(z,z')}$.
\item Let \m{C,a\in\CC}, suppose \m{d\ne0}, and let \m{z\in Z} with $\AdmComb{(x\mapsto C\cdot v(x))}{A}{n}{d}{p}{B}{z}$.
  Then $\AdmComb{v}{A}{n}{1}{p}{d\cdot B}{\nun(aC/d)\cdot a\cdot z}$, where \m{\nun(w)=w^{-1}} for \m{w\ne0} and
  \m{\nun(0)=1}.
\end{enumerate}
\end{lemma}
\begin{proof}
(a) Take the coefficients $c$ on $F$ and $0$ elsewhere. (b) Take the coefficient $1$ at $x$; the size is~$1$, and
\m{d\cdot1\in\intaw{n}}. (c) Apply $L$ to \m{z=\sum_xr(x)v(x)}. (d), (e), (f) follow from Theorem~A(i), proved
below without them. For (d): by (c), $\AdmComb{(x\mapsto c\cdot v(x))}{A}{n}{d}{p}{B}{c\cdot z}$, and each
$c\cdot v(x)$ is the combination of the $v$'s with the single coefficient $c$, which is admissible with
denominator $1$ and size~$1$; substitute. For (e): $0$ and the roots of unity lie in $\intaw{n}$, since a ring
homomorphism to a field sends a root of unity to a root of unity, of absolute value~$1$; the same holds for the
automorphisms~$\sigma$; apply~(d). For (f): by (b) and (e), $\AdmComb{v'}{A'}{n}{1}{p}{1}{v(x)}$ for \m{x\in A};
substitute. (g) By (c) for the linear map \m{b(\,\cdot\,,z')},
$\AdmComb{(x\mapsto b(v(x),z'))}{A}{n}{d}{p}{B}{b(z,z')}$. For \m{x\in A}, by (c) for \m{b(v(x),\,\cdot\,)}, the vector
\m{b(v(x),z')} is an admissible combination of the \m{b(v(x),v'(x'))}, \m{x'\in A'}, with denominator $d'$ and size
$B'$, hence, each of these being a single vector of the family indexed by \m{A\times A'}, an admissible combination
of that family; substitute. (h) Let \m{z=\sum_xr(x)\,C\,v(x)} with $r$ admissible. The coefficients \m{d\cdot r} are
admissible with denominator $1$ and, under~$p$, size \m{d\cdot B}; let \m{z_1=\sum_xd\,r(x)v(x)}. If \m{aC\ne0} then
$\nun(aC/d)\,a\,z=(d/(aC))\cdot aC\sum_xr(x)v(x)=z_1$; if \m{aC=0} then \m{a\,z=0}. In both cases
\m{\nun(aC/d)\,a\,z=\varsigma z_1} with \m{\varsigma\in\{0,1\}}, and (e) applies.
\end{proof}
\leanline{\lean{ofsum}\sep \lean{single}\sep \lean{map}\sep \lean{smul}\sep \lean{smulofeqzeroorisOfFinOrder}\sep
\lean{ofeqisOfFinOrdersmul}\sep \lean{mapbilinear}\sep \lean{normUnitsmul}}

\begin{proof}[Proof of Theorem A]
(i) Write \m{z=\sum_xc(x)v(x)} with $c$ admissible for $(A;n,d;p,B)$; for \m{x\in A} write
\m{v(x)=\sum_yr_x(y)v'(y)} with $r_x$ admissible for $(A';n,d';p,B')$; and write \m{d''=d\,d'\,k}. Put
\m{\rho=\sum_{x\in\operatorname{supp}c}c(x)\,r_x}, a finitely supported function on~$\cX'$. Then
\m{z=\sum_y\rho(y)v'(y)}. The support of $\rho$ lies in~$A'$, because the support of $c$ lies in $A$ and that of every
$r_x$, \m{x\in A}, in~$A'$. Next \m{d''\rho(y)=k\sum_x(d\,c(x))(d'\,r_x(y))} is a sum of products of elements of the
ring $\intaw{n}$. Finally suppose $p$ and let $\sigma$ be an automorphism of~$\CC$. Then \m{B\ge0}, and
\[\sum_y|\sigma(\rho(y))|\le\sum_x|\sigma(c(x))|\sum_y|\sigma(r_x(y))|\le\sum_x|\sigma(c(x))|\cdot|B'|\le B\cdot|B'|\le C.\]
(ii) Write \m{z=\sum_xr(x)v(x)} with $r$ admissible. Then \m{d\cdot z=\sum_x(d\,r(x))\,v(x)}; each \m{d\,r(x)} lies in
\m{\intaw{n}\subseteq O_\jmath}, as $\ell$ is coprime to~$n$, and each $v(x)$ with \m{r(x)\ne0} has \m{x\in A}, so lies
in~$M$.

(iii) Induction on~$m$. For \m{m=0} the pure tensor is \m{1\in\CC}, the single vector of the family at the tuple
\m{y=0}, whose entries lie in the subgroups~$A_j$; use Lemma~\ref{l:rules}(b). For the step from $m$ to \m{m+1},
apply Lemma~\ref{l:rules}(g) to the bilinear map \m{(f,f')\mapsto f\otimes f'}, to $f_m$ (denominator~$d$,
size~$B$) and to the pure tensor of \m{f_{m-1},\dots,f_0} (denominator $d^{m}$, size $|B|^{m}$, by induction): the
pure tensor of \m{f_m,\dots,f_0} is an admissible combination, with denominator \m{d^{m+1}} and size \m{|B|^{m+1}},
of the tensors \m{v_m(x)\otimes(\text{pure tensor at }y)}, for \m{x\in A_{2m}\times A_{2m+1}} and $y$ a tuple with
entries in the~$A_j$. Such a tensor is the vector of the family for \m{m+1} at the tuple that agrees with $y$
except at the places $2m$ and \m{2m+1}, where it has the entries of~$x$; that tuple has its entries in the~$A_j$.
Lemma~\ref{l:rules}(f), with \m{\varsigma=1}, re-indexes the combination.
\end{proof}
\leanline{\lean{bind}\sep \lean{smulmem}\sep \lean{pure}}

\begin{lemma}[integrality away from $n$]\label{l:away}
\begin{enumerate}[label=(\alph*)]
\item Let $n$, $c$, $d$, $k$ be natural numbers such that $c$ divides \m{d\cdot n^{k}}. Then \m{d/c\in\intaw{n}}.
\item Let $E$ be a CM field with a CM type, $I$ a finite index set, and $n$, $m$ natural numbers that are coprime.
  If \m{g\in\Gfin{I}} has \m{g_w\in\mathrm{GL}_I(O_w)} at every finite place $w$ that does not divide~$n$, then
  \m{g_w\in\mathrm{GL}_I(O_w)} at every finite place $w$ that divides~$m$; here $O_w$ is the ring of integers at~$w$.
\item Let $E$ carry Weil data, let $\ell$ be a prime number, \m{\jmath\colon\CC\to\CC_\ell} a ring homomorphism,
  $I$ a finite index set and \m{g\in\Gfin{I}} \kw{integral above} $\ell$: \m{g_w\in\mathrm{GL}_I(O_w)} at every
  finite place $w$ dividing~$\ell$. Then \m{|\jmath(\Df(g))|_\ell=1}.
\item Let $E$ be a CM field with a CM type, $\ell$, $\jmath$, $I$ as in (c), \m{g\in\Gfin{I}} integral above~$\ell$, and
  $\mu$ a character of $\mathrm{U}(1)(\A)$, trivial on $\mathrm{U}(1)(E)$, that has a weight and a tame part in the
  sense of the note on the characters of the data. For any signs on~$I$, with \m{\det(1,g)} the determinant of
  the element $(1,g)$ of the adelic unitary group of these signs: \m{|\jmath(\mu(\det(1,g)))|_\ell=1}.
\end{enumerate}
\end{lemma}
\begin{proof}
(a) Let $\ell$ be a prime coprime to $n$ and \m{\jmath\colon\CC\to\CC_\ell} a ring homomorphism. Write
\m{d\,n^{k}=c\,m} with $m$ a natural number. Since $\ell$ does not divide~$n$, \m{|\jmath(n)|_\ell=1}, so
$|\jmath(d)|_\ell=|\jmath(c)|_\ell\,|\jmath(m)|_\ell$ and \m{|\jmath(d/c)|_\ell=|\jmath(m)|_\ell\le1}, the image of a
natural number being integral. (For \m{c=0} the quotient is $0$ by convention, and the claim is trivial.)
(b) A finite place that divides $m$ does not divide~$n$: its prime ideal would contain \m{1\in m\Z+n\Z}.
(c), (d) In the formal text each is proved by one call of the simplifier, which we did not expand. It unfolds
$\Df(g)$ into the value of the finite part of the splitting character, and of an absolute value, at the
determinant of a matrix $\lambda$ with \m{\lambda\Oh^{I}=\Oh^{I}\cap g^{-1}\Oh^{I}}, which is a unit at the places above
$\ell$ when $g$ is integral there; and it unfolds \m{\mu(\det(1,g))} into a product of values of the tame part and of
the weight at the local determinants, which are units above~$\ell$. The lemmas of the library that say that
such values are $\ell$-adic units under $\jmath$ are quoted, not read.
\end{proof}
\leanline{\lean{divmemintegralAway}\sep \lean{localLevelleintegralAboveofcoprime}\sep \lean{normfinMultipliereqone}\sep
\lean{normpadicdet}}

\begin{lemma}[groups of translations]\label{l:transl}
Let $E$ be a CM field and $n$ a natural number.
\begin{enumerate}[label=(\alph*)]
\item Let $I$, $I_0$ be finite index sets, $M$ an \m{I\times I_0} matrix over $\Af$ with all entries in~$\Oh$,
  \m{d\colon I\to\Af} and \m{d'\colon I_0\to\Af} families, and \m{y\in\wtr{n}{d}}. Suppose that for all
  \m{i\in I}, \m{j\in I_0} and every finite place $w$ with \m{|d_{i,w}|_w\ne|d'_{j,w}|_w} the component of $M_{ij}$ at
  $\bar w$ is~$0$. Then \m{M^{*}y\in\wtr{n}{d'}}.
\item Let $I$, $X$ be index sets, \m{c\in\Oh}, \m{d\colon I\to\Af} a family, \m{x\in\wtr{n}{d}} and
  \m{f\colon X\to I} a map. Then the vector \m{(c\cdot\cj{x_{f(i)}})_{i\in X}} lies in
  \m{\wtr{n}{(\cj{d_{f(i)}})_{i\in X}}}.
\end{enumerate}
\end{lemma}
\begin{proof}
(a) \m{(M^{*}y)_j=\sum_i\cj{M_{ij}}\,y_i}, and the conditions define a group, so fix $i$, $j$ and a finite place
$w$ not dividing~$n$. If \m{|d_{i,w}|_w=|d'_{j,w}|_w}, then
$|d'_j\cj{M_{ij}}y_i|_w=|\cj{M_{ij}}|_w\,|d_iy_i|_w\le1$. Otherwise $M_{ij}$ vanishes at~$\bar w$, so $\cj{M_{ij}}$
vanishes at~$w$. (b) $|\cj{d_{f(i)}}\,c\,\cj{x_{f(i)}}|_w=|c|_w\,|d_{f(i)}x_{f(i)}|_{\bar w}\le1$ for $w$ not dividing
$n$, since then $\bar w$ does not divide $n$ either.
\end{proof}
\leanline{\lean{mulVecmemweightedTranslations}\sep \lean{starmemweightedTranslations}}

\begin{proposition}[a moved translate of the vacuum]\label{p:moved}
Let $E$ carry Weil data and let $V$ be a hermitian space with index set $I$ and its signs. Let $n$, $d$ be natural
numbers, \m{g\in\Gfin{I}}, \m{x\in\Af^{I}}, \m{A\subseteq\Af^{I}} a set, $p$ a proposition and \m{C,c\in\CC}. Suppose
that \m{c\cdot C=\Df(g)}; that \m{g(x+y)\in A} for every \m{y\in\Oh^{I}}; that \m{d\cdot c/|\Df(g)|^{2}\in\intaw{n}}; and
that $p$ implies \m{|\sigma(c)|\le1} for every automorphism~$\sigma$. Then
\[\AdmComb{\bigl(z\mapsto C\cdot\rhf{V}(z)\vac{V}\bigr)}{A}{n}{d}{p}{1}{\om{V}(1,g)\,\rhf{V}(x)\vac{V}}.\]
\end{proposition}
\begin{proof}
By Input~1(b), (c) and (a),
\[\om{V}(1,g)\rhf{V}(x)\vac{V}=\rhf{V}(gx)\,\om{V}(1,g)\vac{V}=\frac{c\,C}{|\Df(g)|^{2}}\sum_y\varsigma_y\,\rhf{V}\bigl(g(x+y)\bigr)\vac{V},\]
the sum over representatives $y$ of the \m{|\Df(g)|^{2}} classes of $\Oh^{I}$ modulo \m{\Oh^{I}\cap g^{-1}\Oh^{I}}, with
roots of unity~$\varsigma_y$. This is a sum, over a finite set, of the vectors
\m{C\varsigma_y\rhf{V}(g(x+y))\vac{V}} with the constant coefficient \m{c/|\Df(g)|^{2}}. Its $d$-fold lies in
$\intaw{n}$ by hypothesis; under $p$ the size is \m{|\Df(g)|^{2}\cdot|\sigma(c)|/|\Df(g)|^{2}\le1}, the positive integer
\m{|\Df(g)|^{2}} being fixed by~$\sigma$. So Lemma~\ref{l:rules}(a) applies, and Lemma~\ref{l:rules}(f) replaces
each vector by \m{C\rhf{V}(z)\vac{V}} with \m{z=g(x+y)\in A}.
\end{proof}
\leanline{\lean{isAdmissibleCombweilvacElem}}

\sect{Lattices of graded spaces}{s:lattices}
\begin{proposition}[the graded lattice basis is unimodular]\label{p:unimod}
Let $X$, $I$ be finite index sets, $l$ and $N$ natural numbers, $R$ a commutative ring with an involution,
\m{\mathbf{e}\colon X^{+}\sqcup\{0,\dots,l-1\}\to I} a bijection, $S$ a commutative ring, $R_0$ a subring of~$R$,
\m{f\colon S\to R} a ring homomorphism, $a$ and $\pi$ units of $R$ and \m{u\in R} with \m{u\bar u=1}. Suppose
\m{f(N)=\pi\bar\pi}. Then the matrix
$\glb{\mathbf{e}}{a}{\pi}{u}^{*}\cdot f(\gram{\mathbf{e}}{N})\cdot\glb{\mathbf{e}}{a}{\pi}{u}$ is hermitian and lies
in \m{\mathrm{GL}_I(R_0)}: its entries and those of its inverse lie in~$R_0$.
\end{proposition}
\begin{proof}
Let $g_i$ be the columns of the graded lattice basis and \m{\langle\,,\rangle} the form of \m{f(\gram{\mathbf{e}}{N})}.
For $l$ even, \m{g_b=ua^{-1}e_b} for \m{2b<l} and \m{g_b=u\bar ae_b} otherwise, so
$\langle g_b,g_{l-1-b}\rangle=\cj{ua^{-1}}\,u\bar a=1$ for \m{2b<l}, and the Gram matrix is
\m{\gram{\mathbf{e}}{N}} itself (there is no middle vector), a permutation matrix equal to its inverse. For $l$ odd
with middle index~$z$: \m{\langle g_b,g_{l-1-b}\rangle=u\bar u=1} for \m{b\ne z};
\m{\langle g_z,g_z\rangle=-u\bar uf(N)=-f(N)};
$\langle g_{\pt},g_z\rangle=\cj{u(\pi\bar\pi)^{-1}}\cdot(-f(N))\cdot u=-1$; and
$\langle g_{\pt},g_{\pt}\rangle=(\pi\bar\pi)^{-1}-u\bar uf(N)(\pi\bar\pi)^{-2}=0$. So the Gram matrix is the image of the
integer matrix that agrees with \m{\gram{\mathbf{e}}{N}} except for the entries $0$ at $(\pt,\pt)$ and $-1$ at
$(\pt,z)$ and $(z,\pt)$. On the span of $e_{\pt}$, $e_z$ it is
$\bigl(\begin{smallmatrix}0&-1\\-1&-N\end{smallmatrix}\bigr)$, with inverse
$\bigl(\begin{smallmatrix}N&-1\\-1&0\end{smallmatrix}\bigr)$; elsewhere it is a permutation matrix. Both
matrices are hermitian with entries in the image of~$\Z$, hence in~$R_0$.
\end{proof}
\leanline{\lean{isUnimodulargradedLatticeBasisgram}}

\begin{proposition}[a real frame]\label{p:frame}
Let $X$, $I$ be finite index sets, $l$ and $N$ natural numbers, $R$ a commutative ring,
\m{\mathbf{e}\colon X^{+}\sqcup\{0,\dots,l-1\}\to I} a bijection and \m{\phi\colon R\to\CC} a ring homomorphism.
Suppose \m{N>0}, and let \m{\varepsilon\colon I\to\{\pm1\}} be signs with exactly \m{\lfloor(l+1)/2\rfloor} minus
signs. Then there is a complex \m{I\times I} matrix $y$ with \m{y^{*}\phi(\gram{\mathbf{e}}{N})y} the diagonal matrix of
$\varepsilon$, and with \m{y^{*}y=1} if \m{N=1}.
\end{proposition}
\begin{proof}
Let \m{\eta_b=-1} for \m{b<\lfloor(l+1)/2\rfloor} and \m{\eta_b=1} for the other string indices. The signs equal to
$1$ on $X^{+}$ and to $\eta$ on the string have \m{\lfloor(l+1)/2\rfloor} minus signs, like~$\varepsilon$, so after
a bijection of the index sets we may assume that they are~$\varepsilon$. Let $y$ be the identity on $X^{+}$ and
have, at a string index $b$ with \m{b'=l-1-b\ne b}, the column \m{y_b=(\eta_be_b+e_{b'})/\sqrt2}, and at the middle
index (if $l$ is odd) the column \m{e_b/\sqrt N}. For \m{b\ne b'} the form pairs $e_b$ with $e_{b'}$ to~$1$, so
\m{\langle y_b,y_b\rangle=\eta_b} and $\langle y_b,y_{b'}\rangle=(\eta_b\eta_{b'}+1)/2=0$, because exactly one of
$b$, $b'$ is smaller than \m{\lfloor(l+1)/2\rfloor}. At the middle index \m{\langle y_b,y_b\rangle=-N/N=-1=\eta_b}. If
\m{N=1} the columns are orthonormal for the standard hermitian form:
\m{|\eta_b|^{2}/2+1/2=1}, \m{(\eta_b+\eta_{b'})/2=0}, and \m{1/N=1}.
\end{proof}
\leanline{\lean{existsframegradedGram}}

\begin{lemma}[vectors in the lattice]\label{l:inlattice}
Let $E$ be a CM field with a CM type, $V$ a hermitian space with index set $I$ and its signs, and $\cY$ a grading
of $V$ of norm~$N$ with definite part $X^{+}$ and a string of odd length~$l$. Let $a$, $\pi$ be finite ideles and
$u$ a finite adele with \m{u\bar u=1} such that \m{\Spans{\cY}{a}{\pi}{u}}. Let \m{c\in\Oh}, and let
\m{x\colon X^{+}\to\Af} be a vector with \m{x_i\in\Oh} for \m{i\in X} and \m{\bar\pi x_{\pt}\in\Oh}. Let
\m{v\in\Af^{I}} be the vector with the coordinates $x_i$ at the indices of $X^{+}$,
\m{u(\pi^{-1}x_{\pt}+c)} at the middle string index and $0$ at the other string indices. Then every coordinate
\m{(x_{V,f}^{-1}v)_i}, \m{i\in I}, lies in~$\Oh$.
\end{lemma}
\begin{proof}
Write \m{x_{V,f}=\glb{\mathbf{e}_{\cY}}{a}{\pi}{u}\cdot k} with \m{k\in\mathrm{GL}_I(\Oh)}. Let $w$ be the vector
with coordinates $x_i$ at \m{i\in X}, \m{\bar\pi x_{\pt}} at~$\pt$, $c$ at the middle index~$z$ and $0$ elsewhere; it
is integral. By Definition~\ref{d:lattice}, for odd~$l$,
\[\glb{\mathbf{e}_{\cY}}{a}{\pi}{u}\,w=\sum_{i\in X}x_ie_i+\bar\pi x_{\pt}\bigl(\bar\pi^{-1}e_{\pt}+u(\pi\bar\pi)^{-1}e_z\bigr)+c\,u\,e_z=v .\]
Hence \m{x_{V,f}^{-1}v=k^{-1}w} is integral.
\end{proof}
\leanline{\lean{invxFinmulVecmemofspansLattice}}

\sect{The finite zero term}{s:ztf}
\begin{proposition}\label{p:ztfspec}
Let $E$ carry Weil data and let $\Xi$ be a splitting of $V$ over $V_0$ along~$H$, with index sets $I$, $I_0$ and
their signs.
\begin{enumerate}[label=(\roman*)]
\item $\ztf{\Xi}$ maps $\Wf{V}$ into $\Wf{V_0}$.
\item For every adelic vector \m{y\in\A^{I}} and every \m{v\in\Fock{V}} there is \m{v_0\in\Fock{V_0}} such that, for all \m{f\in\Sch{V}},
  \[\ev{\Xi}\bigl(\rh{V}(y)(v\otimes f)\bigr)=v_0\otimes\ztf{\Xi}\bigl(\rhf{V}(y_f)f\bigr).\]
\item Let \m{b\in\Af^{H}}, \m{w\in\Af^{I_0}} and \m{z=\mathsf{i}_{\Xi,f}w+\mathsf{k}_{\Xi,f}b}. Then
  $\ztf{\Xi}\bigl(\rhf{V}(z)f\bigr)=\rhf{V_0}(w)\,\ztf{\Xi}(f)$ for every \m{f\in\Sch{V}}.
\end{enumerate}
\end{proposition}
\begin{proof}
Let \m{g=(g_\infty,g_f)} be the adapted element chosen in Definition~\ref{d:ztf}, so that, by Input~1(b),
$v_\Xi\otimes f=\om{V}(g)\bigl(v_{\mathrm{vac}}\otimes\om{V}(1,g_f^{-1})f\bigr)$.

(iii) \m{\rhf{V}(z)} is a root of unity times the Heisenberg operator at the adelic vector with real part $0$ and
finite part \m{\alpha z}, which in frame coordinates is the sum of an $\mathsf{i}$-part and a $\mathsf{k}$-part.
Finite translations commute with \m{v_\Xi\otimes{}}, and Input~3(b) moves them through the zero term: the
$\mathsf{k}$-part disappears and the $\mathsf{i}$-part becomes the translation by~$\alpha w$. Since $\mathsf{k}$ is
isotropic and orthogonal to~$\mathsf{i}$, \m{z^{*}z=w^{*}w}, so the root of unity in $\rhf{V}(z)$ is that of
\m{\rhf{V_0}(w)} and the two translations commute without a factor.

(i) Let \m{f\in\Wf{V}}. Then \m{f_1=\om{V}(1,g_f^{-1})f\in\Wf{V}} and \m{v_{\mathrm{vac}}\otimes f_1=f_1}
(Inputs~1(b), 3(a)), so \m{\ztf{\Xi}(f)=\ev{\Xi}(\om{V}(g)f_1)}, and it suffices to see that
\m{\ev{\Xi}\circ\om{V}(g)} maps every \m{\rhf{V}(x)\vac{V}} into~$\Wf{V_0}$. By Input~1(b),
$\om{V}(g)\rhf{V}(x)\vac{V}=\rhf{V}(g_fx)\om{V}(g)\vac{V}$. Decompose \m{g_fx} along the finite frame as in
Input~3(c), into an $\mathsf{l}$-part, an $\mathsf{i}$-part with coordinates~$w$ and a $\mathsf{k}$-part. By
Input~3(b) the zero term turns the translation by the $\mathsf{i}$-part into~$\rhf{V_0}(w)$ and forgets the
$\mathsf{k}$-part, up to roots of unity, and by Input~3(a) what remains, the zero term of an
$\mathsf{l}$-translate of \m{\om{V}(g)\vac{V}}, is a multiple of~$\vac{V_0}$. So the result is a multiple of
\m{\rhf{V_0}(w)\vac{V_0}}.

(ii) By Input~1(b), \m{\rh{V}(y)(v\otimes f)=v_1\otimes f_2} with \m{f_2=\rhf{V}(y_f)f} and some \m{v_1\in\Fock{V}}.
Again by Input~1(b), $v_1\otimes f_2=\om{V}(g)\bigl((\om{V}(g_\infty)^{-1}v_1)\otimes\om{V}(1,g_f^{-1})f_2\bigr)$.
Input~3(a) gives
\[\ev{\Xi}(v_1\otimes f_2)=\Lambda\bigl(\om{V}(g_\infty)^{-1}v_1\bigr)\otimes\ev{\Xi}\bigl(\om{V}(g)(v_{\mathrm{vac}}\otimes\om{V}(1,g_f^{-1})f_2)\bigr)
=v_0\otimes\ztf{\Xi}(f_2),\]
with \m{v_0=\Lambda(\om{V}(g_\infty)^{-1}v_1)}.
\end{proof}
\leanline{\lean{zeroTermFinspec}\sep \lean{zeroTermFinfinHeisOpfinFrameVector}}

\begin{proposition}[the frequency of a non-zero finite zero term]\label{p:support}
Let $E$ carry Weil data and let $\Xi$ be a splitting of $V$ over $V_0$ along~$H$, with index sets $I$, $I_0$ and
their signs. Let $M$, $M'$ be \m{I\times I} matrices over $\Af$ with \m{M\cdot M'^{*}=1}, let \m{f\in\Sch{V}} with
\m{\rhf{V}(Mx)f=f} for every \m{x\in\Oh^{I}}, and let \m{y\in\Af^{I}} with \m{\ztf{\Xi}(\rhf{V}(y)f)\ne0}. Then there is
\m{x_0\in\Oh^{I}} with \m{\mathsf{k}_{\Xi,f}^{*}\,(y-M'x_0)=0}.
\end{proposition}
\begin{proof}
The \m{I\times H} matrix \m{M'^{*}\mathsf{k}_{\Xi,f}} has a left inverse: if $\mathsf{c}$ is the matrix that reads off
the $\mathsf{k}$-coordinate (Input~3(c)), then $\mathsf{c}MM'^{*}\mathsf{k}_{\Xi,f}=\mathsf{c}\mathsf{k}_{\Xi,f}=1$. By
Hermite reduction over $(\Af,\Oh)$ (Input~6) there is an invertible \m{H\times H} matrix $t$ over $\Af$ such that
\m{M'^{*}\mathsf{k}_{\Xi,f}t} has entries in $\Oh$ and a left inverse $L$ with entries in~$\Oh$. Put
\m{u=t^{*}\mathsf{k}_{\Xi,f}^{*}y} and \m{x_0=L^{*}u}.

We claim \m{u\in\Oh^{H}}; by Input~1(e) it suffices that \m{\psio(x^{*}u)=1} for \m{x\in\Oh^{H}}. For such $x$ the
vector \m{M'^{*}\mathsf{k}_{\Xi,f}tx} is integral, so, by the hypothesis on $f$ and \m{MM'^{*}=1},
\m{\rhf{V}(\mathsf{k}_{\Xi,f}tx)f=f}. By Input~1(e),
$\rhf{V}(y)f=\rhf{V}(y)\rhf{V}(\mathsf{k}_{\Xi,f}tx)f=\psio(x^{*}u)\cdot\rhf{V}(\mathsf{k}_{\Xi,f}tx)\rhf{V}(y)f$. Apply
$\ztf{\Xi}$; by Proposition~\ref{p:ztfspec}(iii), with \m{w=0}, the translation along $\mathsf{k}$ disappears:
$\ztf{\Xi}(\rhf{V}(y)f)=\psio(x^{*}u)\cdot\ztf{\Xi}(\rhf{V}(y)f)$. As this vector is not zero, \m{\psio(x^{*}u)=1}.

So \m{x_0\in\Oh^{I}}, and
$t^{*}\mathsf{k}_{\Xi,f}^{*}(y-M'x_0)=u-(LM'^{*}\mathsf{k}_{\Xi,f}t)^{*}u=0$; as $t^{*}$ is invertible, the claim
follows.
\end{proof}
\leanline{\lean{zeroTermFinsupport}}

\begin{proposition}[the finite zero term of a translated vacuum]\label{p:vac}
Let $E$ carry Weil data and let $\Xi$ be a splitting of $V$ over $V_0$ along~$H$, with index sets $I$, $I_0$ and
their signs. Let \m{g\in\Gfin{I_0}}, let $c$ be an \m{H\times I_0} matrix over~$\Af$, and suppose that the matrix
\m{B=\mathsf{i}_{\Xi,f}\,g+\mathsf{k}_{\Xi,f}\,c} has all its entries in~$\Oh$. Then there is \m{C\in\CC} such that for
every \m{y\in\Af^{I}} there is \m{\varsigma\in\CC}, zero or a root of unity, with
\[\ztf{\Xi}\bigl(\rhf{V}(y)\vac{V}\bigr)=\varsigma\cdot C\cdot\om{V_0}(1,g)\,\rhf{V_0}(B^{*}y)\vac{V_0}.\]
\end{proposition}
\begin{proof}
\emph{The constant.} \m{f_0=\ztf{\Xi}(\vac{V})} lies in $\Wf{V_0}$ (Proposition~\ref{p:ztfspec}(i)). For
\m{x\in\Oh^{I_0}}, by Proposition~\ref{p:ztfspec}(iii) with \m{w=gx} and \m{b=cx},
$\rhf{V_0}(gx)f_0=\ztf{\Xi}(\rhf{V}(Bx)\vac{V})=f_0$, as \m{Bx\in\Oh^{I}} (Input~1(a)). By Input~1(d),
\m{f_0=C\cdot\om{V_0}(1,g)\vac{V_0}} for some~$C$.

\emph{The identity.} Fix~$y$. If \m{\ztf{\Xi}(\rhf{V}(y)\vac{V})=0}, take \m{\varsigma=0}. Otherwise
Proposition~\ref{p:support}, with \m{M=M'=1} and \m{f=\vac{V}}, gives \m{x_0\in\Oh^{I}} with
\m{\mathsf{k}_{\Xi,f}^{*}(y-x_0)=0}. By Input~3(c), \m{y-x_0=\mathsf{i}_{\Xi,f}w+\mathsf{k}_{\Xi,f}b} with
\m{w=\mathsf{i}_{\Xi,f}^{*}(y-x_0)}. Hence, with roots of unity \m{\varsigma_1}, \m{\varsigma_2} from Input~1(a),
\[\begin{aligned}
\ztf{\Xi}(\rhf{V}(y)\vac{V})&=\ztf{\Xi}(\rhf{V}(y)\rhf{V}(-x_0)\vac{V})=\varsigma_1\ztf{\Xi}(\rhf{V}(y-x_0)\vac{V})\\
&=\varsigma_1\rhf{V_0}(w)f_0=\varsigma_1C\,\rhf{V_0}(w)\om{V_0}(1,g)\vac{V_0}.
\end{aligned}\]
Since $g$ is unitary and \m{\mathsf{k}_{\Xi,f}^{*}(y-x_0)=0}, we have
$gB^{*}(y-x_0)=gg^{*}\mathsf{i}_{\Xi,f}^{*}(y-x_0)=w$, so by Input~1(b)
$\rhf{V_0}(w)\om{V_0}(1,g)=\om{V_0}(1,g)\rhf{V_0}(B^{*}(y-x_0))$; and \m{B^{*}x_0\in\Oh^{I_0}} fixes the vacuum, so
$\rhf{V_0}(B^{*}(y-x_0))\vac{V_0}=\varsigma_2\rhf{V_0}(B^{*}y)\vac{V_0}$. Take \m{\varsigma=\varsigma_1\varsigma_2}.
\end{proof}
\leanline{\lean{zeroTermFinfinHeisOpvacElem}}

\sect{Arrow pairs}{s:arrow}
\begin{lemma}[two computations in rational coordinates]\label{l:coord}
Let $E$ be a CM field with a CM type and $V$, $V'$ hermitian spaces with index sets $I$, $I'$ and their signs.
\begin{enumerate}[label=(\alph*)]
\item For a pair \m{g=(g_1,g_2)} of elements of $\Gg{V}$ and $\Gg{V'}$, the matrix in rational coordinates of
  \m{g_1\oplus g_2}, an element of \m{\Gg{V\oplus V'}}, is the block diagonal matrix of those of $g_1$ and~$g_2$.
\item Let $\cY$, $\cY'$ be gradings of $V$, $V'$ of norm~$N$, with definite parts $X$, $X'$ and strings of lengths
  $l$, $l'$. For \m{g\in\Gg{\St{X}}} and \m{g'\in\Gg{\St{X'}}}, in rational coordinates,
  $\inX{\cY'}(g')\cdot\Tm{\cY}{\cY'}\cdot\inX{\cY}(g)^{-1}=\Tm{\cY}{\cY'}$.
\end{enumerate}
\end{lemma}
\begin{proof}
(a) The frame of \m{V\oplus V'} is block diagonal. (b) In rational coordinates $\inX{\cY}(g)$ is block diagonal for
the decomposition into $X$ and the string, with the identity on the string, and likewise~$\inX{\cY'}(g')$; the
frequency map has its only non-zero block from string to string.
\end{proof}
\leanline{\lean{ratCoordMatrixsum}\sep \lean{inclXmulstringMapmulinclXinv}}

\begin{lemma}\label{l:atfev}
In Setting~\ref{set:three}, let $X''$, $H$ be finite index sets and \m{\cR\in\ArP{\cY}{\cY'}{\cY''}{X''}{H}}. For all
\m{\Phi\in\Wf{\Hm{V}{V'}}} and \m{\Phi'\in\Wf{\Hm{V'}{V''}}},
\[\at\bigl(\fev{\cR}(\Phi\boxtimes\Phi')\bigr)=\ztf{\aps{\cR}}\Bigl(\rhf{}\bigl(\stTf{\cY}{\cY'}{\cY''}\bigr)(\Phi\boxtimes\Phi')\Bigr),\]
where $\at$ is the map of the standard spaces $\St{X}$, $\St{X'}$, $\St{X''}$.
\end{lemma}
\begin{proof}
\m{\Phi\boxtimes\Phi'=\at(\Phi\otimes\Phi')} lies in \m{\Wf{\Hm{V'}{\fls{V}{V''}}}} (Input~2(e)); the finite Heisenberg
operators preserve finite parts (Input~1(a)); and $\ztf{\aps{\cR}}$ maps finite parts to finite parts
(Proposition~\ref{p:ztfspec}(i)). So the right side lies in \m{\Wf{\Hm{\St{X'}}{\fls{\St{X}}{\St{X''}}}}}, which is the
range of $\at$ for the standard spaces, and \m{\at\circ\atl} is the identity on that range because
\m{\atl\circ\at=1}.
\end{proof}
\leanline{\lean{arrowTensorfinEval}}

\begin{proof}[Proof of Theorem B]
(i) Put \m{\wt=\gr{\cY}\sqcup\gr{\cY''}}, write \m{l'=2j} (possible as $l$ is odd and \m{l'=l-1}), let $m_+$ be the
number of indices of \m{I\sqcup I''} of positive weight, and let
$x=\inX{\cY}(g)^{-}\oplus\inX{\cY''}(k)\in\Gg{\fls{V}{V''}}$. In rational coordinates $x$ is the identity on the basis
vectors of non-zero weight, because $\inX{\cY}$ and $\inX{\cY''}$ fix the string vectors (Lemma~\ref{l:coord}(a)).
The splitting $\aps{\cR}$ is a composite, so its zero term is the composite of two zero terms
(Definition~\ref{d:split}), and
$\tW{\mu}{\mu'}(\inX{\cY}g,\inX{\cY'}g',\inX{\cY''}k)$ is a scalar \m{c_V} times the Weil operator of
\m{(\inX{\cY'}(g'),x)} on \m{\Hm{V'}{\fls{V}{V''}}} (Definition~\ref{d:arrowprod}).

\emph{First zero term.} Apply Input~3(d) to the splitting \m{\gzs{\hzp{\cR}}}, whose frame $\mathsf{k}$ consists of
the basis vectors \m{(i,b)} of positive degree \m{\wt(i)-\gr{\cY'}(b)}, with the element \m{(\inX{\cY'}(g'),x)}, the
element $g_0=(g',1\oplus g^{-}\oplus\operatorname{emb}_{\cR}(k))$ of the group of the zero part, and the matrix $C$
that acts by $x$ on the rows above each \m{b\in I'} of negative grade and by the inverse transpose of
\m{\inX{\cY'}(g')} on the columns above each \m{i} of positive weight, and is the identity elsewhere. Its
determinant is \m{u=(\det g\cdot\det k)^{j}\cdot(\det g')^{-m_+}}, since $\cY'$ has $j$ string vectors of negative
grade, \m{\det x=\det g\cdot\det k} and \m{\det\inX{\cY'}(g')=\det g'}. The conditions of Input~3(d) are statements
about the entries of block diagonal matrices: they hold because $x$ is the identity off weight zero, because the
zero part has degree zero only above~$X'$ (Definition~\ref{d:gzphom}), and, for the last one, by condition~(c) of
Definition~\ref{d:arrow}. Hence the first zero term carries the Weil operator of \m{(\inX{\cY'}(g'),x)} to
\m{\upsilon(u)^{-1}} times that of~$g_0$.

\emph{Second zero term.} By Input~3(e), removing the $H$ planes carries the Weil operator of
$(g',1\oplus g^{-}\oplus\operatorname{emb}_{\cR}(k))$ to \m{\upsilon(\det g')^{|H|}} times that of
\m{(g',g^{-}\oplus\operatorname{emb}_{\cR}(k))}; the side conditions of the input hold for the explicit frames of the
splitting that removes the planes.

\emph{The scalars.} The right side of (i) is the scalar $c_X$ of the triple action for the standard spaces $\St{X}$,
$\St{X'}$, $\St{X''}$ at \m{(g,g',\operatorname{emb}_{\cR}(k))} times the same Weil operator. So it remains to see
$c_V=c_X\cdot\upsilon(u)\,\upsilon(\det g')^{-|H|}$. The values of $\mu$ and $\mu'$ in $c_V$ and $c_X$ agree, since
\m{\det\inX{\cY}(g)=\det g} and \m{\det\operatorname{emb}_{\cR}(k)=\det k}. It remains to compare the exponents of
$\upsilon$ at the three determinants, with \m{\mu_V=\mu\upsilon^{\lfloor|I|/2\rfloor}}. At \m{\det g}:
$(|I'|-\lfloor|I'|/2\rfloor)-(|X'|-\lfloor|X'|/2\rfloor)=2j-j=j$, as \m{|I'|=|X'|+2j}. At \m{\det k}:
\m{\lfloor|I'|/2\rfloor-\lfloor|X'|/2\rfloor=j}. At \m{\det g'}: counting the weights of \m{I\sqcup I''},
\m{|I|+|I''|=2m_++2|H|+|X|+|X''|}; this number is even by condition~(b) of Definition~\ref{d:arrow}, so
$\lfloor|I|/2\rfloor-\lfloor|I''|/2\rfloor=(|I|-|I''|)/2$ and likewise for $X$, $X''$, and the exponent is
\[\tfrac12(|I|-|I''|)-|I|-\tfrac12(|X|-|X''|)+|X|=-\tfrac12(|I|+|I''|)+\tfrac12(|X|+|X''|)=-m_+-|H|.\]
These are the exponents of \m{\upsilon(u)\upsilon(\det g')^{-|H|}}.

(ii) Since \m{\atl\circ\at=1}, Input~2(d), Lemma~\ref{l:atfev} and the hypothesis
\m{\operatorname{emb}_{\cR}(1,k)=(1,k')} give for the right side
\[\atl\Bigl(\tW{\mu}{\mu'}\bigl((1,g),(1,g'),\operatorname{emb}_{\cR}(1,k)\bigr)\,\ztf{\aps{\cR}}\bigl(\rhf{}(\mathbf{t})(\Phi\boxtimes\Phi')\bigr)\Bigr),
\qquad\mathbf{t}=\stTf{\cY}{\cY'}{\cY''}.\]
Now \m{\ztf{\aps{\cR}}(F)=\ev{\aps{\cR}}(v\otimes F)} for a fixed real datum~$v$, and the triple action of finite points
commutes with \m{v\otimes{}} (Input~1(b)). So part~(i), with \m{\inXf{\cY''}(k)=x}, turns this into the first of the
two vectors
\[\begin{aligned}
&\atl\,\ztf{\aps{\cR}}\bigl(\tW{\mu}{\mu'}(\inXf{\cY}g,\inXf{\cY'}g',x)\,\rhf{}(\mathbf{t})(\Phi\boxtimes\Phi')\bigr),\\
&\atl\,\ztf{\aps{\cR}}\bigl(\rhf{}(\mathbf{t})\,\tW{\mu}{\mu'}(\inXf{\cY}g,\inXf{\cY'}g',x)(\Phi\boxtimes\Phi')\bigr),
\end{aligned}\]
and by Input~2(d) the left side of (ii) is the second.
The two agree: on an arrow product the translation by the stacked frequency map is the arrow product of the
translations by the finite parts of the two frequency maps (Input~2(b)), the triple action is the arrow product
of the two pair actions (Input~2(d)), and each pair action commutes with the translation by its frequency map
by Input~2(g) and Lemma~\ref{l:coord}(b).
\end{proof}
\leanline{\lean{zeroTermsplittingtripleWeil}\sep \lean{finEvalinclXFin}}

\sect{Two unfolding steps}{s:double}
\begin{proof}[Proof of Theorem C(i)]
Let \m{z_0\ne0} be the constant of Input~2(a) for $V$, $V'$, $V''$. Let $C_1$ be the constant of Input~4(a) for
$\cY$ and the set of all its indices of grade~$1$; this set is empty, the grades of $\cY$ being even. Let $C_2$
be the constant of Input~4(a) for $\cY'$ and the injection $\gamma_{\cR}$. We show the claim with
\m{z=C_1C_2z_0^{-1}}. Fix $\mu$, $\mu'$, $f$, $f'$, $k$, $h$ and a generic~$P$, and put
\[F(g,r)=P(gh)\,\ct{\cY}{\cY'}{\mu}(f)(g,r)\,\ct{\cY'}{\cY''}{\mu'}(f')\bigl(r,\inX{\cY''}(k)\bigr).\]
It is continuous, and invariant under the rational points of $\Bes{\cY}$ in $g$ and of $\Bes{\cY'}$ in~$r$
(Input~4(c)).

\emph{Step 1: a product of two constant terms is a zero term.} For data $f_1$, $f_2$ and matrices
\m{y\colon V\to V'}, \m{y'\colon V'\to V''} in rational coordinates, the product of the constant term of $f_1$
at $y$ and of $f_2$ at~$y'$ equals
$z_0^{-1}\cdot\ev{\Xi_1}\bigl(\rh{}(y)f_1\boxtimes\rh{}(y')f_2\bigr)(0)$, where \m{\Xi_1=\gzs{\hzp{\cR}}}. Indeed, the
value at $0$ of this zero term is the mean of the datum over the matrices \m{V'\to\fls{V}{V''}} of positive degree
(Definition~\ref{d:split}). The map $(a,b)\mapsto\bigl(\begin{smallmatrix}a^{\dagger}\\b\end{smallmatrix}\bigr)$
carries the pairs of a matrix \m{V\to V'} and a matrix \m{V'\to V''}, both of positive degree, onto these
matrices, because the adjoint exchanges the positions \m{(b,a)} and \m{(a,b)} up to the reversal of the strings,
which changes the sign of both grades. So the mean is an iterated mean (Input~7(c)); by Input~2(b) and~(a) the
integrand is
$\tht{}\bigl(\rh{}(a)\rh{}(y)f_1\boxtimes\rh{}(b)\rh{}(y')f_2\bigr)=z_0\,\tht{}(\rh{}(a)\rh{}(y)f_1)\,\tht{}(\rh{}(b)\rh{}(y')f_2)$,
and the two means of the factors are the two constant terms (Definition~\ref{d:grading}).

\emph{Step 2: the inner integral.} Fix~$r$. For a unipotent element $n$ of $\cY$, \m{n-1} raises grades by at
least~$2$, the grades being even; so by Input~4(b), at the frequency map, where the two values of the bicharacter
cancel, \m{F(n\,\inX{\cY}(a),r)=F(\inX{\cY}(a),r)} for \m{a\in\Gg{\St{X}}}. Input~4(a), with the constant function of
mean \m{F(\inX{\cY}(a),r)}, gives
$\int_{\qB{\cY}}F(g,r)\,dg=C_1\int_{\qG{\St{X}}}F(\inX{\cY}(a),r)\,da$.

\emph{Step 3: Fubini.} The function \m{(a,r)\mapsto F(\inX{\cY}(a),r)} is continuous and invariant on a product of
two compact quotients, so the integrals over $r$ and $a$ may be exchanged.

\emph{Step 4: the integral over $\qB{\cY'}$.} Fix \m{a\in\Gg{\St{X}}} and let \m{g'\in\Gg{\St{X'}}} and $n'$ a
unipotent element of~$\cY'$. Put \m{f_1=\pW{\mu}(\inX{\cY}a,\inX{\cY'}g')f} and
\m{f_2=\pW{\mu'}(\inX{\cY'}g',\inX{\cY''}k)f'}. By Input~4(c) the two constant terms in
\m{F(\inX{\cY}(a),n'\inX{\cY'}(g'))} are the constant terms of $f_1$ at \m{n'^{-1}\Tm{\cY}{\cY'}} and of $f_2$ at
\m{\Tm{\cY'}{\cY''}n'}. Since $n'$ is unitary, $(n'^{-1}\Tm{\cY}{\cY'})^{\dagger}=\Tm{\cY}{\cY'}^{\dagger}n'$, so by Step~1 and
Input~2(b)
\[F\bigl(\inX{\cY}(a),n'\inX{\cY'}(g')\bigr)=P(\inX{\cY}(a)h)\cdot z_0^{-1}\cdot
\ev{\Xi_1}\Bigl(\rh{}\bigl(\mathbf{T}+\mathbf{T}(n'-1)\bigr)(f_1\boxtimes f_2)\Bigr)(0),\]
where we write $\mathbf{T}$ for \m{\stT{\cY}{\cY'}{\cY''}}.
The matrix \m{\stT{\cY}{\cY'}{\cY''}(n'-1)} has two kinds of entries. Outside the positions of degree
zero, which are the positions \m{(\ind{\cR}(a'),x')} with \m{x'\in X'}, it has entries only in positions of positive
degree, because the frequency maps lower degrees by at most one and \m{n'-1} raises them; translations in these
directions do not change the zero term (Input~3(b)). At the positions of degree zero, condition~(d) of
Definition~\ref{d:arrow} shows that its entries are those of \m{\inc{\cR}} applied to the vector that has, at the
second copy of \m{k\in H} and at \m{x'\in X'}, the entry of \m{\fc{\cY'}(n'-1)} at \m{(\gamma_{\cR}(k),x')}, and zeros
elsewhere. By Input~3(b) the zero term turns
this translation into the translation of the zero part by that vector; the Heisenberg cocycles that occur are
trivial. The directions of these vectors are exactly those over which the splitting that removes the planes
averages. Therefore \m{n'\mapsto F(\inX{\cY}(a),n'\inX{\cY'}(g'))} is a continuous periodic function of the entries of
\m{\fc{\cY'}(n'-1)} in the rows $\gamma_{\cR}(H)$ and the columns~$X'$, and its mean is
\[M(g')=P(\inX{\cY}(a)h)\cdot z_0^{-1}\cdot\tht{}\Bigl(\ev{\aps{\cR}}\bigl(\rh{}(\stT{\cY}{\cY'}{\cY''})(f_1\boxtimes f_2)\bigr)\Bigr).\]
Now $f_1\boxtimes f_2=\tW{\mu}{\mu'}(\inX{\cY}a,\inX{\cY'}g',\inX{\cY''}k)(f\boxtimes f')$ (Input~2(d)); this triple action
commutes with \m{\rh{}(\stT{\cY}{\cY'}{\cY''})}, as in the proof of Theorem~B(ii); and Theorem~B(i) moves it through
the zero term. So
\[M(g')=P(\inX{\cY}(a)h)\,z_0^{-1}\,\tht{}\Bigl(\tW{\mu}{\mu'}\bigl(a,g',\operatorname{emb}_{\cR}(k)\bigr)\,\ev{\aps{\cR}}\bigl(\rh{}(\mathbf{T})(f\boxtimes f')\bigr)\Bigr),\]
a continuous function of $g'$, left invariant under \m{\Gamma(\St{X'})} by Input~2(f). Input~4(a) for $\cY'$ gives
$\int_{\qB{\cY'}}F(\inX{\cY}(a),r)\,dr=C_2\int_{\qG{\St{X'}}}M(g')\,dg'$. Together with Steps 2 and~3 this is the
claim.
\end{proof}
\leanline{\lean{integralconstTermAtTmul}}

\begin{proof}[Proof of Theorem C(ii)]
Let \m{z_0\ne0} be the constant of Input~2(a) for the standard spaces $\St{X}$, $\St{X'}$, $\St{X_0}$. For
\m{o\in U}, $\varphi$, $g$ and a datum $f$ of \m{\Hm{\St{X'}}{\fls{\St{X}}{\St{X_0}}}}, the function
$(x,r)\mapsto\coef{t}(\varphi)(x)\cdot\tht{}(\tW{\mu_o}{\mu'_o}(x,r,g)f)$ on
\m{\qG{\St{X}}\times\qG{\St{X'}}} is continuous and well defined (Inputs~5 and~2(f)), on a compact space since the
standard spaces are definite; so it is integrable, and its integral $\mathcal{I}_{o,\varphi,g}(f)$ is a continuous
linear functional of~$f$.

\emph{Reduction to polynomial vectors.} The real data $u$ for which the theorem holds form a linear subspace
of the Fock space. It is closed: the map sending $u$ to the function
$(o,y,\varphi,g)\mapsto\mathcal{I}_{o,\varphi,g}(u\otimes\at(y))$ is continuous for pointwise convergence, and the
functions of the form $(o,y,\varphi,g)\mapsto\coef{t''}(\mathcal{L}''_o(y)\varphi)(g)$, $t''$ in the finite-dimensional
dual of the vectors of weight~$\varkappa$, form a finite-dimensional, hence closed, subspace. By Input~7(a) it
suffices to treat \m{u=v_Q}, and by Input~7(b) and linearity a product \m{Q=Q_1Q_2} of a polynomial $Q_1$ in the
coordinates of \m{\Hm{\St{X'}}{\St{X}^{-}}} and a polynomial $Q_2$ in those of \m{\Hm{\St{X'}}{\St{X_0}}}.

\emph{Products.} For such $Q$ the polynomial vector is an arrow product, \m{v_Q=v_{Q_1'}\boxtimes v_{Q_2}} on real
data, where $Q_1'$ is $Q_1$ read in the coordinates of \m{\Hm{\St{X}}{\St{X'}}}. Let $t'$ be the functional that
Input~5 gives for $(\mathcal{L}_o)_o$, $Q_1'$ and~$t$, and $t''$ the one for $(\mathcal{L}'_o)_o$, \m{z_0Q_2} and~$t'$.
By linearity in $y$ it suffices to take \m{y=\Phi\otimes\Phi'}. Then
$v_Q\otimes\at(y)=(v_{Q_1'}\otimes\Phi)\boxtimes(v_{Q_2}\otimes\Phi')$, and by Input~2(d) and~(a)
\[\tht{}\bigl(\tW{\mu_o}{\mu'_o}(x,r,g)(v_Q\otimes\at(y))\bigr)
=z_0\cdot\tht{}\bigl(\pW{\mu_o}(x,r)(v_{Q_1'}\otimes\Phi)\bigr)\cdot\tht{}\bigl(\pW{\mu'_o}(r,g)(v_{Q_2}\otimes\Phi')\bigr).\]
Integrating first over $x$ (Fubini) and using the property of $\mathcal{L}_o$, then over $r$ and using that of
$\mathcal{L}'_o$, we get
\[\begin{aligned}
\mathcal{I}_{o,\varphi,g}\bigl(v_Q\otimes\at(y)\bigr)&=\int_{\qG{\St{X'}}}\coef{t'}\bigl(\mathcal{L}_o(\Phi)\varphi\bigr)(r)\,
\tht{}\bigl(\pW{\mu'_o}(r,g)(v_{z_0Q_2}\otimes\Phi')\bigr)\,dr\\
&=\coef{t''}\bigl(\mathcal{L}'_o(\Phi')\mathcal{L}_o(\Phi)\varphi\bigr)(g).
\end{aligned}\qedhere\]
\end{proof}
\leanline{\lean{integralthetaeqlift}}

\sect{Admissible evaluations}{s:admev}
\begin{proposition}[the standard evaluation]\label{p:std}
In Setting~\ref{set:adm}, with its objects $N$, $\Sigma$, $\pi$, $\delta$, $\beta$, $\kappa$, let $X'$ be a finite index set,
put \m{T=X^{+}\sqcup X''^{+}}, with the points $\pt_1$ of $X^{+}$ and $\pt_2$ of~$X''^{+}$, and let
$V_{\mathrm{st}}=\Hm{\St{X'}}{\fls{\St{X^{+}}}{\St{X''^{+}}}}$, a standard space with index set \m{T\times X'}. There are an
element \m{k\in\Gfin{T\times X'}}, a \m{T\times T} matrix $\gamma$ over $\Af$ and a family \m{b\colon T\to\Af} such that:
\begin{enumerate}[label=(\arabic*)]
\item the matrix of $k$ is the Kronecker product of \m{\operatorname{diag}(\kappa)\cdot\gamma} with the identity matrix
  of~$X'$;
\item at every finite place $w$ with \m{w\notin\Sigma} and \m{\bar w\notin\Sigma}, every $b_j$ has component~$0$ and
  $\gamma$ has the identity as component;
\item for every \m{j\in T}, \m{\NbV{\pi}{\delta}{b_j}{\gamma_{\cdot j}}}, where $\gamma_{\cdot j}$ is the column $j$ of~$\gamma$;
\item $\AdmEv{N}{\Sigma}{\beta}{\bigl(f\mapsto\atl(\om{V_{\mathrm{st}}}(1,k)f)\bigr)}{\soc{\beta}{\,\cdot\,}\kappa}$
  (Definition~\ref{d:admev} for the spaces $\St{X^{+}}$, $\St{X'}$, $\St{X''^{+}}$).
\end{enumerate}
\end{proposition}
\begin{proof}
\emph{The element.} Let \m{\mathbf{1}_\Sigma\in\Oh} be the adele equal to $1$ at the places of $\Sigma$ and $0$
elsewhere. Then \m{\mathbf{1}_\Sigma\cj{\mathbf{1}_\Sigma}=0} and \m{\mathbf{1}_\Sigma\delta=\mathbf{1}_\Sigma}, as
\m{\Sigma\cap\bar\Sigma=\emptyset} and \m{\delta=1} above~$N$. Let $\lambda$ be the finite idele equal to $\pi$ at
the places of~$\Sigma$, to \m{\bar\pi^{-1}} at their conjugates and to $1$ elsewhere; so \m{\lambda\bar\lambda=1}. Let
\m{\xi\in\Oh^{T}} be the vector \m{e_{\pt_1}+\cj{\epsm\mathbf{1}_\Sigma}\,e_{\pt_2}}; then \m{\xi^{*}\xi=1}, and
\[\gamma=1+(\lambda-1)\,\xi\xi^{*}\]
is unitary: it multiplies $\xi$ by $\lambda$ and fixes the vectors orthogonal to~$\xi$. For \m{m\in T} put
$b_m=(\lambda-1)\pi^{-1}\bigl(\cj{\mathbf{1}_\Sigma}\,[m=\pt_1]+\epsm\mathbf{1}_\Sigma\,[m=\pt_2]\bigr)$, where
\m{[m=\pt_1]} is $1$ if \m{m=\pt_1} and $0$ otherwise.
The diagonal matrix \m{\operatorname{diag}(\kappa)} is unitary, because \m{\kappa_b\in\{1,\delta\}} and
\m{\delta\bar\delta=1}, and it commutes with~$\gamma$, because \m{\gamma-1} is supported above $N$, where all
$\kappa_b$ are~$1$. Let $k$ be the element of \m{\Gfin{T\times X'}} acting by \m{\operatorname{diag}(\kappa)\gamma}
on each column \m{x'\in X'}. This is~(1), and (2) holds because \m{\lambda-1} vanishes outside
\m{\Sigma\cup\bar\Sigma}.

\emph{Neighbour vectors.} \m{\delta b_m=b_m} since $b_m$ is supported above~$N$. The entries \m{\gamma_{im}} with
\m{i\ne\pt_1,\pt_2} are those of the identity. Next \m{\bar\pi(\lambda-1)\in\Oh}: it is \m{\bar\pi(\pi-1)} at $\Sigma$,
\m{1-\bar\pi} at $\bar\Sigma$ and $0$ elsewhere; so \m{\bar\pi\gamma_{im}\in\Oh}. Let
$\varpi_0=\mathbf{1}_\Sigma+(1-\mathbf{1}_\Sigma)\pi^{-1}$, an element of $\Oh$ since $\pi$ is a unit outside~$\Sigma$.
A direct computation with \m{\mathbf{1}_\Sigma\cj{\mathbf{1}_\Sigma}=0} and \m{\cj{\epsm}\epsm=-1} gives
\m{b_m-\pi^{-1}\gamma_{\pt_1m}=-\varpi_0\,[m=\pt_1]} and $\cj{\epsm}b_m-\pi^{-1}\gamma_{\pt_2m}=-\varpi_0\,[m=\pt_2]$:
for instance $\gamma_{\pt_1m}=[m=\pt_1]+(\lambda-1)([m=\pt_1]+\epsm\mathbf{1}_\Sigma[m=\pt_2])$, so the coefficient of
\m{[m=\pt_2]} cancels and that of \m{[m=\pt_1]} is
$-\pi^{-1}(1+(\lambda-1)(1-\cj{\mathbf{1}_\Sigma}))=-\pi^{-1}(1+\mathbf{1}_\Sigma(\pi-1))=-\varpi_0$. This is~(3).

\emph{The key identity.} Let \m{\bar\Sigma} be the set of the conjugates of the places of~$\Sigma$ and
$\gamma_{X'}$ the element acting by $\gamma$ on each column \m{x'\in X'}. We claim that there is \m{C_0\in\CC} with
\[\om{V_{\mathrm{st}}}(1,\gamma_{X'})\,P^{\bar\Sigma}f=C_0\sum_y\rhf{}\bigl(\cj{\mathbf{1}_\Sigma}\,\mathrm{n}(y)\bigr)P^{\bar\Sigma}f
\qquad\text{for all }f\in\Sch{V_{\mathrm{st}}},\]
the sum over \m{y\in((\Z/N)^{\Sigma})^{X'}}, where \m{\mathrm{n}(y)} is the finite vector with the entry
\m{\epsm\,\xi_i\,\cj{\tilde y_{x'}/N}} at \m{(i,x')}. By Input~3(f), $P^{\bar\Sigma}$ is named by the indicator of the
vectors integral at~$\bar\Sigma$, the element $\gamma_{X'}$, supported at $\Sigma$ and $\bar\Sigma$, transforms
this function by \m{\gamma_{X'}^{-1}} up to a constant, and translates of the indicator name the corresponding
Heisenberg translates of~$P^{\bar\Sigma}$. So the claim is an identity of functions: for every finite vector $x$
there is an element $b(x)$ such that, for every $y$, the vector \m{\gamma_{X'}x+\mathrm{n}(y)} is integral at
$\bar\Sigma$ if and only if $x$ is integral at $\bar\Sigma$ and \m{y=b(x)}. Thus the sum over $y$ of the translated
indicators has at most the one term \m{y=b(x)} at \m{\gamma_{X'}x}, and no term when $x$ is not integral
at~$\bar\Sigma$. At a place of $\bar\Sigma$ we have \m{\lambda=\bar\pi^{-1}},
\m{\xi=e_{\pt_1}+\cj{\epsm}e_{\pt_2}}, and the column $x'$ of \m{\gamma_{X'}x+\mathrm{n}(y)} is
$x_{\cdot x'}+\xi\cdot\bigl((\bar\pi^{-1}-1)x_{\pt_1x'}+\epsm\,\tilde y_{x'}/N\bigr)$. If it is integral, its coordinate
at $\pt_1$, \m{\bar\pi^{-1}x_{\pt_1x'}+\epsm\tilde y_{x'}/N}, is integral; multiplying by the uniformiser $\bar\pi$ and
using that \m{\bar\pi/N=\pi^{-1}} is a unit there, \m{x_{\pt_1x'}} is integral, and then so are all coordinates
of~$x$. Conversely, for integral $x$ the column is integral exactly if
\m{\bar\pi^{-1}x_{\pt_1x'}+\epsm\tilde y_{x'}/N} is. The lemma of the note on the Gauss-sum operator on residues
modulo~$N$ at a split level, applied after complex conjugation at~$\Sigma$, provides a residue class $b_{x'}$ such
that this holds if and only if \m{y_{x'}=b_{x'}}, the product of $N$ and \m{\bar\pi^{-1}x_{\pt_1x'}} being integral;
the family of the $b_{x'}$ is~$b(x)$.

\emph{The admissible evaluation.} We show (4) with the constant \m{C=\Df(k_0)\,C_0}, where $k_0$ acts by
\m{\operatorname{diag}(\kappa)} on each column, so that \m{k=\gamma_{X'}k_0}. Let $n$ be coprime to~$N$, $p$ a
proposition and \m{x\in\wtr{n}{(i,x')\mapsto\soc{\beta}{i}\kappa_i}}. First, Proposition~\ref{p:moved}, with \m{c=1}
and the set \m{A=\wtr{n}{(i,x')\mapsto\soc{\beta}{i}}}, shows that
\m{\om{V_{\mathrm{st}}}(1,k_0)\rhf{}(x)\vac{}} is an admissible combination of the vectors
\m{\Df(k_0)\rhf{}(z)\vac{}}, \m{z\in A}, away from~$n$, with denominator $D$ and size~$1$. Its hypotheses hold: $k_0$
multiplies the coordinate \m{(i,x')} by~$\kappa_i$, and \m{\soc{\beta}{i}\kappa_i(x+y)_{(i,x')}} is integral away from $n$
for \m{y\in\Oh^{T\times X'}}, by the choice of $x$ and because \m{\soc{\beta}{i}\kappa_i\in\Oh}; and
\m{D/|\Df(k_0)|^{2}\in\intaw{n}} by Lemma~\ref{l:away}(a) with \m{k=0}, because \m{|\Df(k_0)|^{2}}, the index of
\m{\Oh^{T\times X'}\cap k_0^{-1}\Oh^{T\times X'}} (Input~1(c)), divides~$D$: this lattice contains the kernel of the
map \m{y\mapsto(\bar y\bmod\beta)} to \m{(\Oh/\beta\Oh)^{T\times X'}}, since a vector with coordinates \m{\bar\beta\bar c},
\m{c\in\Oh}, is mapped by $k_0$ to one with coordinates \m{\bar\beta\kappa_i\bar c\in\Oh}. Second, apply the linear map
\m{\atl\circ\om{V_{\mathrm{st}}}(1,\gamma_{X'})} (Lemma~\ref{l:rules}(c)). Third, re-index by
Lemma~\ref{l:rules}(f): for \m{z\in A},
\[\Df(k_0)\,\atl\bigl(\om{V_{\mathrm{st}}}(1,\gamma_{X'})\rhf{}(z)\vac{}\bigr)
=\varsigma\cdot C\cdot\tsum{X}{X'}{\Sigma}{N}{X''}\bigl(\rhf{}(t_1)\vac{}\otimes\rhf{}(t_2)\vac{}\bigr)\]
with a root of unity $\varsigma$, where $t_1$ is \m{\epsm} times the conjugate transpose of the block of $z$ with
rows $X^{+}$ and $t_2$ is the block of $z$ with rows~$X''^{+}$; both lie in the groups of translations for
$\bar\beta$ (Lemma~\ref{l:transl}(b)). To see this identity, note that $z$ is integral at the places above~$N$:
they do not divide $n$, and $\beta$ and $\bar\beta$, the weights of the two blocks, are units there. So \m{P^{\bar\Sigma}\rhf{}(z)\vac{}=\rhf{}(z)\vac{}}
(Input~1(f)), and the key identity gives
$\om{V_{\mathrm{st}}}(1,\gamma_{X'})\rhf{}(z)\vac{}=C_0\sum_y\rhf{}(\cj{\mathbf{1}_\Sigma}\mathrm{n}(y))\rhf{}(z)\vac{}$. By
Input~2(c) each term is a root of unity times an arrow product of two translated vacua, and the terms are those of
$\at\bigl(\tsum{X}{X'}{\Sigma}{N}{X''}(\rhf{}(t_1)\vac{}\otimes\rhf{}(t_2)\vac{})\bigr)$, by the definition of the
translate sum (the first block of \m{\cj{\mathbf{1}_\Sigma}\mathrm{n}(y)} gives the last column \m{-\tilde y/N}, the second
block the last row \m{-\cj{\tilde y/N}}, with the same root of unity for all~$y$); apply $\atl$.
\end{proof}
\leanline{\lean{existsisAdmissibleEvalstd}}

\begin{proof}[Proof of Theorem D]
Apply Proposition~\ref{p:std} with the index set \m{X'^{+}} in place of~$X'$: it gives $k$, $\gamma$, $b$ with the
properties (1) to~(4), and the constant $C$ of~(4). Let $r_1$ be the finite idele \m{\cj{\epsm}\bar\alpha/\alpha}, by
which the frame of $V^{-}$ differs from that of~$V$ (Definition~\ref{d:arrowprod}); so the inverse finite frame of
\m{\fls{V}{V''}} is the block diagonal matrix of \m{r_1^{-1}x_{V,f}^{-1}} and~\m{x_{V'',f}^{-1}}, and
\m{r_1\cdot(a/\bar a)=\cj{\epsm}\delta} by the hypothesis on~$\delta$.

\emph{An integral matrix.} Let $M$ be the matrix with rows \m{(H\sqcup H)\sqcup T} and columns \m{j\in T} whose column
$j$ is: $0$ on the first copy of~$H$; \m{r_1ub_j} at every index of the second copy; \m{r_1\kappa_i\gamma_{ij}} at
\m{i\in X^{+}}; and \m{\kappa_i\gamma_{ij}} at \m{i\in X''^{+}}. Let $\mathsf{J}$ be the inclusion matrix of the basis
vectors of $X'^{+}$ into~$I'$, and
\[B=\bigl(x_{\fls{V}{V''},f}^{-1}\,\inc{\cR}\,M\bigr)\otimes\bigl(\mathsf{J}^{\mathsf T}x_{V',f}\bigr)^{\mathsf T},\]
a matrix with rows \m{(I\sqcup I'')\times I'} and columns \m{T\times X'^{+}}. It is integral. The second factor is:
\m{x_{V',f}=\glb{\mathbf{e}_{\cY'}}{a}{\pi}{u}\,k'} with \m{k'\in\mathrm{GL}_{I'}(\Oh)}, and for the even length $l'$ the
rows $X'^{+}$ of the graded lattice basis are those of the identity. The first factor is by the hypothesis
(lattices), applied to $r_1$, \m{r_2=r_1^{-1}} and the neighbour vectors \m{(b_j,\gamma_{\cdot j})} of~(3): the column
$j$ of \m{\inc{\cR}M} is the vector $v$ of that hypothesis.

\emph{The zero term.} By the definitions of the splittings, \m{B=\mathsf{i}_fk+\mathsf{k}_fc'} for the frames of
$\aps{\cR}$ and a suitable matrix~$c'$: the columns of $M$ split into the part \m{(r_1\kappa_i\gamma_{ij},\kappa_i\gamma_{ij})},
which is the frame of \m{\fls{\St{X^{+}}}{\St{X''^{+}}}} applied to \m{\operatorname{diag}(\kappa)\gamma}, and the part
along the second copy of~$H$, which is the isotropic frame of the splitting that removes the planes. So
Proposition~\ref{p:vac} gives \m{C_3\in\CC} and, for every finite vector~$y$, a number \m{\varsigma(y)}, zero or a
root of unity, with
$\ztf{\aps{\cR}}(\rhf{}(y)\vac{})=\varsigma(y)\,C_3\,\om{V_{\mathrm{st}}}(1,k)\,\rhf{}(B^{*}y)\vac{}$.

\emph{The frequency map disappears.} \m{B^{*}\,\stTf{\cY}{\cY'}{\cY''}=0}. Indeed
$(\mathsf{J}^{\mathsf T}x_{V',f})^{*}=x_{V',f}^{-1}\mathsf{J}$, by \m{x_{V',f}^{*}J_{V'}x_{V',f}=1} and because $J_{V'}$ is
the identity on~$X'^{+}$; and \m{\stT{\cY}{\cY'}{\cY''}\,\mathsf{J}=0}, the frequency maps and their adjoints having
non-zero columns only at string vectors.

\emph{Conclusion.} Let $n$ be coprime to~$N$, $p$ a proposition and \m{y\in\wtr{n}{(i,i')\mapsto\soc{\beta}{i}}}. By
Definition~\ref{d:fev}, Input~1(a) and the two preceding paragraphs,
\[\begin{aligned}
\fev{\cR}\bigl(\rhf{}(y)\vac{}\bigr)&=\varsigma'\,\atl\,\ztf{\aps{\cR}}\bigl(\rhf{}(\stTf{\cY}{\cY'}{\cY''}+y)\vac{}\bigr)\\
&=\varsigma'\,\varsigma(y)\,C_3\cdot\atl\bigl(\om{V_{\mathrm{st}}}(1,k)\,\rhf{}(B^{*}y)\vac{}\bigr)
\end{aligned}\]
with a root of unity~$\varsigma'$. By Lemma~\ref{l:transl}(a) the vector \m{B^{*}y} lies in
\m{\wtr{n}{(j,x')\mapsto\soc{\beta}{j}\kappa_j}}; its condition on the valuations is checked below. So~(4) applies to
\m{B^{*}y}, and Lemma~\ref{l:rules}(c) and~(e) carry the scalars $C_3$, \m{\varsigma(y)}, $\varsigma'$: the claim holds
with the constant \m{C_3C}.

\emph{The condition on the valuations.} Let \m{(i,i')} be a row and \m{(j,x')} a column of~$B$, and $w$ a finite
place with \m{|\soc{\beta}{i}|_w\ne|\soc{\beta}{j}\kappa_j|_w}; we show that the first factor of the entry of $B$ vanishes
at~$\bar w$. That factor is $\sum_{i_1}(x_{\fls{V}{V''},f}^{-1})_{ii_1}(\inc{\cR}M)_{i_1j}$. Fix~$i_1$ and use the
hypothesis (weights) for $i_1$ and~$j$. If \m{\inc{\cR}(i_1,j)=0}: the place $\bar w$ does not lie above~$N$, since
above $N$ all $\kappa$ are $1$ and $\beta$ is a unit, so that the two absolute values would both be~$1$; hence, by~(2),
at $\bar w$ the $b$'s vanish and $\gamma$ is the identity, the column $j$ of $M$ has at $\bar w$ only the entry at
the index~$j$, and \m{(\inc{\cR}M)_{i_1j}} is a multiple of \m{\inc{\cR}(i_1,j)=0} there. Otherwise
$|\soc{\beta}{i_1}|_w=|\soc{\beta}{j}\kappa_j|_w\ne|\soc{\beta}{i}|_w$, so $i_1$ and $i$ lie in different ones of the two
sets $I$, $I''$, and \m{(x_{\fls{V}{V''},f}^{-1})_{ii_1}=0} because this matrix is block diagonal.
\end{proof}
\leanline{\lean{isAdmissibleEvalfinEval}}

\sect{Where the results are used}{s:used}
\emph{In the unit.} The note on the seesaw identity uses: Theorem~A and Lemma~\ref{l:rules} wherever a
denominator or a size is tracked (the unfolding to an admissible combination, the combination of reduced vacua,
the passage to a good datum); Lemma~\ref{l:away}(a) for the multiplicities of the finite cuts, which divide powers
of the conductor of the twist; Theorem~C, with the middle pairs and the last pair of Definition~\ref{d:pairs} for
the gradings of the spaces of the chain of lifts, to identify the function left after $2i+2$ unfoldings with a matrix
coefficient of \m{i+1} pairs of definite lifts; Theorem~B(ii) for the invariance of the evaluated data under the
levels; Theorem~D, Lemma~\ref{l:inlattice} and Proposition~\ref{p:moved}, with Lemma~\ref{l:away}(b) to~(d), to write
the evaluated data through reduced vacua; and Proposition~\ref{p:support} in the step that bounds how far the
lines of a lattice move under one lift.

\emph{In other units.} Admissible coefficients enter the definition of a good datum in the unit of the columns,
and admissible combinations are produced in the unit of the CM curve (named data as combinations of translated
vacua) and used in the units of the exchange identity and of the rows. Propositions~\ref{p:unimod}
and~\ref{p:frame} are used in the unit that constructs the spaces of the chain of lifts and in the unit of the filling:
the first shows that the lattice spanned by the graded lattice basis is self-dual, the second provides the real
part of the frame.

\sect{Notes on the formal text}{s:notes}
\subhead{Arguments.} The lists of arguments used by the checks of this note are read from the text of the
artifact, not printed by Lean. They contain entries that Lean does not take, as the reader's dossier records: an index
type listed for \m{\intaw{n}}; an index type listed for admissible evaluations; and a first copy of the field
instances for Proposition~\ref{p:std} and for Theorem~D. The two index sets $X''$, $H$ of an arrow pair are listed
twice, as implicit and as explicit arguments, because a prefix makes them explicit for that declaration. For
admissible evaluations the lists also carry a finite index type named like the proposition of
Definition~\ref{d:admev}, with its finiteness: the definition binds that name as the proposition~$p$ and its
translate sum has the three index sets $X$, $X'$, $X''$ only, so these two entries are inferred from the text, as
the second reader of the dossier notes; a signature printed by Lean would settle them. In the other direction, the text reading of Theorem~C(ii) misses the field $E$, the index sets $X$,
$X'$, $X_0$ and their instances, which the formal statement names after a local definition; they stand in the
statement above. In Theorem~A(i) the two hypotheses on $d''$ and $C$ have default proofs by simplification, so
that the formal text can omit them at a use.

\subhead{Differences from the statements above.} (1) Coprimality of a prime $\ell$ and $n$ is the formal
hypothesis in Theorem~A(ii) and Definition~\ref{d:away}; for \m{n=0} no prime is coprime to~$n$, so
\m{\intaw{0}=\CC}. (2) The automorphisms $\sigma$ are ring automorphisms of~$\CC$. (3) The ring homomorphism $\jmath$
is an instance argument in Theorem~A(ii) and an ordinary one in Lemma~\ref{l:away}(c),~(d). (4) The finite zero term
and the left inverse $\atl$ depend on choices made once by the axiom of choice (an adapted element; a left
inverse of an injective linear map). (5) The denominator $D$ of Definition~\ref{d:admev} is a cardinality that
is~$0$ when infinite. (6) In Definition~\ref{d:arrow} the subtraction is that of natural numbers. (7) In
Theorem~C(ii) the hypothesis is the library's predicate on families of lifts, stated here as Input~5; the
characters enter through the twist \m{\mu\mapsto\mu\upsilon^{\lfloor|X|/2\rfloor}}. (8) Proposition~\ref{p:frame} takes
a ring homomorphism $\phi$ that only transports an integer matrix, and Proposition~\ref{p:unimod} a homomorphism
\m{f\colon S\to R} of the same kind. (9) The formal statements of Lemma~\ref{l:rules}(a) and of Theorem~A(iii)
are for a finite subset given as such, and for the iterated tensor product of the library, in which the factor
of index $m-1$ stands first. (10) In Lemma~\ref{l:rules}(h) the scalar is Mathlib's normalising unit, with its
value $1$ at zero; its description is recalled, Mathlib not being part of the text read. The scalar must not be
simplified to \m{d/C}: nothing is assumed non-zero, and for \m{a=0} the vector of the conclusion is zero. (11) In
Definition~\ref{d:transl} and in Theorem~D the two blocks of a vector carry the weights $\beta$ and~$\bar\beta$;
they are never replaced by one common weight.

\subhead{What the formal proofs do differently.} The unitary matrix $\gamma$ of Proposition~\ref{p:std} is given
entry by entry in the formal text, and its unitarity and the properties of neighbour vectors are found by
automation from a list of identities; the description \m{\gamma=1+(\lambda-1)\xi\xi^{*}} and the computations with
it are ours. The comparison of the scalars in Theorem~B(i) is, formally, a linear-arithmetic certificate; the
count with the weights is ours. The identity of functions behind the key identity of Proposition~\ref{p:std},
and the condition on the valuations in Theorem~D, are long simplification chains in the formal text, which we
read as the identities they establish and derived again by hand. Lemma~\ref{l:away}(c),~(d) are not expanded.
Which hypotheses the formal proofs leave unused was not determined: nothing was compiled for this note.

\subhead{The dossier and its second reading.} The reader's dossier describes Input~3(e) without the factor
\m{\upsilon(\det h)^{|H|}}. The formal statement has it, through a twist of one of the two characters, and the
scalars of Theorem~B(i) agree only with it. A second reader compared the dossier of the unit with the Lean text
after this note was drafted. Her corrections that concern this note are followed here: the zero case of
Lemma~\ref{l:rules}(h), where the dossier had simplified the scalar; the two weights of the blocks; the
index types and the proposition in Definition~\ref{d:admev}; and, in Proposition~\ref{p:std}, that the comparison
of indicators leaves at most one term.

\subhead{How this note was written.} From the structured proofs of the reader's dossier for the eight files of
this note, after its program check passed. We read these eight files of the Lean text
completely, the statements argument by argument, and line by line the proofs of Theorem~B(i),
Proposition~\ref{p:std} and Theorem~D, which the reader had marked as read from the names of the lemmas only;
the two statements with a local definition (Theorem~C) were read in full. Of the background library we read the
statements and comments of the declarations quoted as Inputs~1 to~7, and none of their proofs. No Lean was run.
The program check of the unit's dossier covered \UnitSteps{} leaf steps of \UnitProofs{} proofs.

\head{Appendix A. Concordance with the Lean artifact}
The first \UnitInNote{} rows are the declarations of the unit treated here; the rows whose place is an Input are
declarations of the background library, quoted and not read.

{\fontsize{8.6}{11.6}\selectfont\renewcommand{\leanfont}{\ttfamily\fontsize{7.7}{10}\selectfont\color{leancol}}
\setlength{\LTleft}{0pt}\setlength{\LTright}{\fill}\setlength{\tabcolsep}{0pt}
\begin{longtable}{@{}p{30mm}p{56mm}p{43mm}r@{}}
\textbf{In this note} & \textbf{Lean declaration} & \textbf{File} & \textbf{Lines}\\\midrule\endhead
\ConcordanceRows
\end{longtable}}

\subhead{Legend.} Formal names of objects and our letters: \emph{Space} with its matrix and point: a hermitian
space $V$ with $J_V$ and the frame~$x_V$; the data space, the Heisenberg and finite Heisenberg operators, the
vacuum, the finite part: $\Sch{V}$, $\rh{V}$, $\rhf{V}$, $\vac{V}$, $\Wf{V}$; \emph{integrated}: \m{v\otimes f};
\emph{Splitting} with its fields \emph{k}, \emph{l}, \emph{i}: $\Xi$ with $\mathsf{k}_\Xi$, $\mathsf{l}_\Xi$,
$\mathsf{i}_\Xi$; \emph{zeroTerm}: $\ev{\Xi}$; \emph{Grading} with \emph{indexEquiv}, \emph{c}, \emph{gradeOf},
\emph{inclX}, \emph{stringMap}: $\cY$ with $\mathbf{e}_{\cY}$, $\fc{\cY}$, $\gr{\cY}$, $\inX{\cY}$, \m{\Tm{\cY}{\cY'}};
\emph{arrowProduct}, \emph{arrowTensor}, \emph{tripleWeil}, \emph{liftWeilFin}: $\boxtimes$, $\at$, \m{\tW{\mu}{\mu'}},
$\pWf{\mu}$; the fields \emph{zeroPart}, \emph{embX}, \emph{gradeOne} of an arrow pair: $\cZ_{\cR}$,
$\operatorname{emb}_{\cR}$, $\gamma_{\cR}$; \emph{normNegOne}: $\epsm$; \emph{translateSum}: \m{\tsum{X}{X'}{\Sigma}{N}{X''}}.

\head{Notation}
\begingroup\fontsize{8.6}{10.9}\selectfont\raggedright
\setlength{\columnsep}{9mm}\setlength{\parskip}{0pt}
\begin{multicols}{2}
\nt{$E$, $\A$, $\Af$, $\Oh$}{a CM field with a CM type; its adeles, finite adeles, integral finite adeles}
\nt{$w$, $\bar w$, $|\cdot|_w$, $O_w$}{a finite place, its conjugate, its absolute value, its ring of integers}
\nt{$\mathbf{j}$, $\alpha$, $\upsilon$, $\epsm$}{Weil data: an element of $E$, an idele, the splitting character; a unit
  of $\Oh$ of norm $-1$}
\nt{$\ell$, $\jmath$, $O_\jmath$, $\sigma$}{a prime; a ring homomorphism \m{\CC\to\CC_\ell}; its ring of integral
  complex numbers; an automorphism of~$\CC$}
\nt{$\intaw{n}$}{complex numbers integral away from~$n$}
\nt{$\cX$, $\cX'$, $\cX_j$, $A$, $A'$, $A_j$}{index sets of families of vectors and subsets of them}
\nt{$Z$, $Z'$, $W$, $F_i$, $M$}{complex vector spaces; an $O_\jmath$-submodule (also: a matrix)}
\nt{$v$, $v'$, $v_i$, $z$, $z'$, $f_i$}{families of vectors and vectors (also: $v$ a real datum, $f$ a datum or a map)}
\nt{$n$, $d$, $d'$, $d''$, $m$, $k$}{natural numbers: the modulus of integrality, denominators, a number of
  factors, an exponent (also: $k$ a group element, $d$ a family of adeles)}
\nt{$p$, $B$, $B'$, $C$, $C_0$, $C_1$, $C_2$, $C_3$, $D$}{a proposition; real bounds; complex constants; a denominator}
\nt{$r$, $r_x$, $c$, $\rho$, $\varsigma$}{coefficient functions; a root of unity or zero (also: $r$ a group element,
  $c$ an adele or matrix)}
\nt{\m{\AdmCo{A}{n}{d}{p}{B}{r}}}{admissible coefficients}
\nt{$\AdmComb{v}{A}{n}{d}{p}{B}{z}$}{admissible combination}
\nt{$\nun$, $L$, $b$}{the inverse or $1$; a linear map; a bilinear map (also: $b$ an index or a vector)}
\nt{$I$, $I'$, $I''$, $I_0$, $I_1$, $I_2$, $X$, $X'$, $X''$, $X_0$, $H$, $T$, $U$}{finite index sets; a set of
  parameters $o$}
\nt{$X^{+}$, $\pt$, $\pt_1$, $\pt_2$}{an index set with one more point; the points}
\nt{$V$, $V'$, $V''$, $V_0$, $V_{\mathrm{st}}$, $J_V$, $x_V$, $x_{V,f}$}{hermitian spaces; matrix, frame, finite frame}
\nt{$\Gg{V}$, $\Gfin{I}$, $\Gamma(V)$, $\qG{V}$}{the adelic unitary group, its finite part, rational points, quotient}
\nt{$\St{X}$, \m{\Hm{V}{V'}}, \m{\fls{V}{V''}}, \m{\Hyp{H}{V}}}{standard space; homomorphisms; first form negated; with
  hyperbolic planes}
\nt{$\Sch{V}$, $\Fock{V}$, $\Wf{V}$, $\vac{V}$, $v_Q$, $Q$}{data, real data, finite part, vacuum, polynomial vectors,
  polynomials}
\nt{$\rh{V}$, $\rhf{V}$, $\om{V}$, $\tht{V}$, $\Df$, $\psio$, $\psi_f$, $P^{S}$}{Heisenberg and finite Heisenberg
  operators, Weil representation, theta functional, finite multiplier, additive characters, a projection}
\nt{$y$, $y_f$, $y^{\dagger}$, $x$, $u$, $g$, $g'$, $h$, $a$}{vectors and matrices, finite coordinates, adjoint; group
  elements (also: $a$ an idele, $u$ a unitary adele)}
\nt{$\mu$, $\mu'$, $\mu_V$, $\mu_o$}{characters of $\mathrm{U}(1)(\A)$; twisted by a power of~$\upsilon$}
\nt{$\pW{\mu}$, $\pWf{\mu}$, \m{\tW{\mu}{\mu'}}}{twisted pair action, its finite part, the triple action}
\nt{$\boxtimes$, $\at$, $\atl$}{arrow product; the induced map on finite parts; a left inverse}
\nt{$\Xi$, $\mathsf{k}_\Xi$, $\mathsf{l}_\Xi$, $\mathsf{i}_\Xi$, $\ev{\Xi}$, $\ztf{\Xi}$}{a splitting, its frames, its
  zero term, its finite zero term}
\nt{$\Lambda$, $\epsilon$, $v_{\mathrm{vac}}$, $v_\Xi$}{objects of Input~3: a map of real data, a function, the real
  vacuum and its translate}
\nt{$\cY$, $\cY'$, $\cY''$, $\mathbf{e}_{\cY}$, $l$, $l'$, $l''$, $N$, $\fc{\cY}$}{gradings: bijection, string lengths,
  norm, constant}
\nt{$\gr{\cY}$, $\inX{\cY}$, $\inXf{\cY}$, $\Bes{\cY}$, $\qB{\cY}$, \m{\Tm{\cY}{\cY'}}}{grade, inclusion of the definite
  part, Bessel group and quotient, frequency map}
\nt{\m{\ct{\cY}{\cY'}{\mu}(f)}, $P$, $\gram{\mathbf{e}}{N}$}{constant term at the frequency map; a generic function; the
  Gram matrix of a string}
\nt{\m{\wtr{n}{d}}, \m{\trl{I}{n}{d}}, \m{\soc{\beta}{j}}}{groups of translations; $\beta$ or its conjugate}
\nt{\m{\SplSq{\Sigma}{N}{\pi}}}{the split squarefree condition}
\nt{\m{\glb{\mathbf{e}}{a}{\pi}{u}}, \m{\Spans{\cY}{a}{\pi}{u}}}{graded lattice basis; the lattice condition}
\nt{$R$, $R_0$, $S$, $\phi$, $\varepsilon$, $\eta_b$, $e_i$, $g_i$}{rings and a homomorphism; signs; basis vectors and
  columns}
\nt{$\wt$, \m{\GZP{V}{V''}{\wt}}, $\cZ$, $\ind{\cZ}$, $\inc{\cZ}$, \m{\gzs{\cZ}}, \m{\gzh{\cY'}{\cZ}}}{weights; grade-zero
  parts with injection and matrix; their splitting; the part of a space of homomorphisms}
\nt{\m{\stT{\cY}{\cY'}{\cY''}}, \m{\stTf{\cY}{\cY'}{\cY''}}}{the stacked frequency map and its finite coordinates}
\nt{\m{\ArP{\cY}{\cY'}{\cY''}{X''}{H}}, $\cR$, $\cZ_{\cR}$, $\inc{\cR}$, $\ind{\cR}$, $\operatorname{emb}_{\cR}$,
  $\gamma_{\cR}$}{arrow pairs and their components}
\nt{\m{\hzp{\cR}}, \m{\aps{\cR}}, \m{\fev{\cR}}}{grade-zero part, splitting and finite evaluation of an arrow pair}
\nt{\m{\Rmid{\cY}{\cY'}{\cY''}}, \m{\Rlast{\cY}{\cY'}{\cY''}{\pi_0}}, $\pi_0$}{the middle pair and the last pair}
\nt{$\varkappa$, $e$, $e'$, $e''$, \m{\AF{X}{\varkappa}{e}}, $\varphi$, $t$, $t'$, $t''$, \m{\coef{t}}}{weight and
  determinant exponents; algebraic forms; functionals; coefficients}
\nt{$\mathcal{L}_o$, $\mathcal{L}'_o$, $\mathcal{L}''_o$, $\Phi$, $\Phi'$, $\Psi$, $\mathcal{I}$}{theta lifts and their
  composite; finite data; a function that names an operator; an integral}
\nt{$\Sigma$, $\bar\Sigma$, $\pi$, $\delta$, $\beta$, $\kappa$}{a split level and its conjugate; the objects of
  Setting~\ref{set:adm}}
\nt{\m{\NbV{\pi}{\delta}{c}{x}}}{neighbour vector}
\nt{$\AdmEv{N}{\Sigma}{\beta}{L}{d}$, \m{\tsum{X}{X'}{\Sigma}{N}{X''}}, $\tilde y$}{admissible evaluation; the translate
  sum; least non-negative representatives}
\nt{$\gamma$, $\gamma_{X'}$, $\lambda$, $\xi$, $\mathbf{1}_\Sigma$, $\varpi_0$, $\mathrm{n}(y)$, $k_0$, $B$,
  $\mathsf{J}$, $\mathsf{c}$}{objects of the proofs of Proposition~\ref{p:std} and Theorem~D}
\nt{$z_0$, $F$, $M$, $m_+$, $j$, $i$, $q$, $G$, $K$, $O$, $Y$, $s$}{further letters of the proofs: constants, functions,
  counts and indices; $G$, $O$, $Y$, $K$, $s$, $q$ occur only inside the symbols above}
\nt{$\mathsf{A}$, $\mathsf{A}'$, $\mathsf{B}$}{topological abelian groups in Input~7}
\end{multicols}
\endgroup

\end{document}
